\pdfinfoomitdate=1
\pdftrailerid{}
\pdfsuppressptexinfo=15
\documentclass[11pt,letterpaper]{article}
\usepackage[T1]{fontenc}
\usepackage[margin=1in]{geometry}
\usepackage{lmodern,microtype,amsmath,amssymb,amsthm,mathtools,booktabs,array,longtable}
\usepackage[colorlinks=true,linkcolor=black,citecolor=black,urlcolor=blue]{hyperref}
\usepackage{xurl}
\usepackage{enumitem,bookmark}
\setlist{nosep}
\setlength{\emergencystretch}{3em}
\allowdisplaybreaks[2]
\newtheorem{theorem}{Theorem}[section]
\newtheorem{lemma}[theorem]{Lemma}
\newtheorem{proposition}[theorem]{Proposition}
\newtheorem{corollary}[theorem]{Corollary}
\theoremstyle{remark}\newtheorem{remark}[theorem]{Remark}
% Macros of the three sources. Every name defined by two or three sources has
% the same definition in each (\C, \R, \e, \ii, \dd, \Oh, \Rea, \Ima).
\newcommand{\C}{\mathbb C}\newcommand{\R}{\mathbb R}
\newcommand{\N}{\mathbb N}\newcommand{\e}{\mathrm e}
\newcommand{\ii}{\mathrm i}\newcommand{\dd}{\,\mathrm d}
\newcommand{\Oh}{\mathrm O}\newcommand{\supp}{\operatorname{supp}}
\newcommand{\Rea}{\operatorname{Re}}\newcommand{\Ima}{\operatorname{Im}}
\newcommand{\asympC}{\asymp}
\newcommand{\Log}{\operatorname{Log}}
% Added in the merge (batch 106).
\newcommand{\file}[1]{\nolinkurl{#1}}
\newenvironment{writenote}[1][]{\par\smallskip\begingroup\small
 \noindent\textbf{[write]}\if\relax\detokenize{#1}\relax\else~\textbf{#1.}\fi\ }%
 {\par\endgroup\smallskip}
\newcommand{\wtag}{\textup{\textbf{[write]}}}
\newenvironment{partabstract}{\begin{quote}\small\noindent\textbf{Abstract of this Part (as delivered).}\ }{\end{quote}}
\newcommand{\partsource}[1]{\par\noindent\begin{minipage}{\textwidth}\small #1\end{minipage}\par\medskip}
% Equations are numbered within each Part as delivered: Part I keeps (k),
% Parts II and III print Report 226's and 228's equation (k) as (II.k), (III.k).
\newcommand{\parteqs}[1]{\setcounter{equation}{0}%
 \renewcommand{\theequation}{#1\arabic{equation}}%
 \renewcommand{\theHequation}{#1\arabic{equation}}}
% Added by the ProveIt write of Part IV (9 October 2026): the macros of Part IV.
\newcommand{\coef}[1]{[z^{#1}]}
\newcommand{\Repin}{\texttt{58175ca45563d9ee29875374dd77942f069268a9}}
\hypersetup{pdftitle={Iterated Euler diagonals (OEIS A290353 and A290354)},pdfauthor={Part IV: research companion prepared for Vladimir Reshetnikov, developed with OpenAI ChatGPT},pdfsubject={Merged report a290354-iterated-euler-diagonals (batch 106): the diagonal and proportional-height equivalents (Report 225), the first logarithmic correction (Report 226), the logarithmic height crossover (Report 228); Part IV (added 9 October 2026): a common density for parabolic iteration}}
\title{\textbf{Iterated Euler diagonals\\ (OEIS A290353 and A290354)}\\[7pt]
\large Kot\v{e}\v{s}ovec's asymptotic form, the first logarithmic correction,\\ and the logarithmic height crossover\\[6pt]\normalsize Part IV: a common density for parabolic iteration}
\author{Three research reports of one external research session\\(bundle Reports 225, 226 and 228), merged into one ProveIt report\\[2pt]Part IV: research companion prepared for Vladimir Reshetnikov,\\developed with OpenAI ChatGPT}
\date{Parts I--III: 5 October 2026\quad Part IV: 8 October 2026\quad merged: 6 and 9 October 2026}
\begin{document}
\hypersetup{pageanchor=false}
\maketitle
\begin{abstract}
Let $F_0(z)=z$ and $1+F_{m+1}(z)=\prod_{d\ge1}(1-z^d)^{-[z^d]F_m(z)}$, so that $A(n,m)=[z^n]F_m(z)$ counts unlabeled rooted trees with $n$ leaves, all at level $m$ (OEIS A290353), and $a_n=A(n,n)$ is OEIS A290354. This report prints three manuscripts as Parts~I--III, in the order in which each depends on the one before. Part~I proves $a_n\sim c\,(2n/\pi)^nn^{-4/3}$ with $c=T^{1/3}e^{-B}h(1)>0$, $T=\pi/2$, $B=\frac12-\frac\pi6-\frac{\pi^2}{16}-\frac{\log2}6$, where $h$ is the smooth, strictly positive density of a positive measure whose Laplace transform is the entire parametrization $P$ of $w\mapsto e^w-1$ with normalized Fatou coordinate $-2/w+\frac13\log w+O(w)$; this proves the form of the asymptotic that V\'aclav Kot\v{e}\v{s}ovec stated as a conjecture on A290354 (August 14, 2017), while the digits of $c\approx4.4923$ remain uncertified. Part~I also proves proportional-height asymptotics for $A(n,m)$ with $n/m$ in a compact subset of $(0,\infty)$, fixed and square-root height shifts with the floor phase, and a Lambert-$W$ inverse at exact sequence values. Part~II proves the first logarithmic correction $1+(-\frac{1}{18}\log^2n+A\log n+C)/n+O(n^{-7/5})$, with $A$ and $C$ exact in $h(1),h'(1),h''(1)$ and a new clock constant $\kappa$ in closed form in $\pi$, $\log2$ and Catalan's constant, the corresponding uniform proportional-height correction and a sharpened inverse. Part~III proves $A(n,m)/(2T^{-n}m^{n-1})\to e^{-1/(3\lambda)}$ when $m/(n\log n)\to\lambda\in(0,\infty)$, with every fixed inverse-logarithm order, through an exact Hankel representation of $h$ near zero, $h(0)=2$, a smooth germ $t^{t/3}h(t)$ and rational-plus-zeta coefficients; it also gives Kaneiwa's normalization and smooth height prescriptions. Validated digits of $c$, higher corrections, effective constants, the regime $n/m\to\infty$ and a conjectured identity linking $h$ with the amplitude of the labelled iterated Bell numbers (a neighbouring report) remain open. The report is unrefereed, and nothing in it is formalized.
\end{abstract}
\hypersetup{pageanchor=true}
\setcounter{tocdepth}{1}
{\small\setlength{\parskip}{0pt}\tableofcontents}
\newpage
\setlength{\parskip}{4pt}

\section*{Guide to this report}\phantomsection\label{ied:sec:guide}
\addcontentsline{toc}{section}{Guide to this report}
The three Parts are three manuscripts of the session bundle of Reports 1--243 (arrival commit \texttt{60f54ea06}), placed in this directory by commit \texttt{47fc7a069} (batch~106). They study one object, the iterated Euler transforms $F_m$ of $1+x$ and their coefficient array $A(n,m)$:
\begin{center}
\small
\begin{tabular}{@{}l l >{\raggedright\arraybackslash}p{0.5\textwidth} l@{}}
\toprule
Part & Bundle report & Delivered title & Labels\\
\midrule
I & 225 (base) & Diagonal and proportional height asymptotics for iterated Euler transforms & \texttt{ied:lead:}\\
II & 226 & First logarithmic correction for iterated Euler diagonals & \texttt{ied:corr:}\\
III & 228 & Logarithmic height crossover for iterated Euler transforms & \texttt{ied:log:}\\
IV & --- & A Common Density for Parabolic Iteration (added 9 October 2026) & \texttt{ied:par:}\\
\bottomrule
\end{tabular}
\end{center}
\emph{Arrangement.} The base of the merge is Report~225, whose \file{report225.tex} was staged as this \file{article.tex}. It proves the leading theorems and constructs the analytic foundation (the nonautonomous clock, the sectorial Fatou coordinate, the positive Laplace density $h$), so it is printed first, as Part~I, with its section, statement and equation numbers unchanged. Report~226 uses Part~I's analytic interface and refines its clock and its coefficient extraction to the first correction; it is Part~II. Report~228 uses both, Part~I's coordinate and density and Part~II's absolute derivative estimates, and pushes the transfer to degree-to-height ratios shrinking like a power of $1/\log m$; it is Part~III. The three Parts thus appear in dependency order, which is also the order of writing.

\emph{Embedded copies.} Report~226's archive carries Report~225's source and PDF under \file{supporting/}, and Report~228's carries those of Reports~225 and~226; all six files are byte-identical to the standalone deliveries (SHA-256 at placement, \texttt{cmp} again at the write). They are printed once each, as Parts~I and~II, and are not shipped. Report~226's archive also carries byte copies of eight code and data files of Report~225 (its ledger says ``copied unchanged''); they are shipped once, under Report~225's prefix.

\emph{Restatements.} Part~II's Section~\ref{ied:corr:sec:interface} and Part~III's Section~\ref{ied:log:sec:interfaces} state, in their own words, the results they import (from Part~I, and for Part~III also from Part~II). They are printed in full, with notes naming the Part~I or Part~II statement behind each item, because their displays are cited by the later proofs, because Part~II states as interface items several estimates that Part~I proves only inside its proofs (the inverse-orbit bound and the disk and neighbourhood forms of the tail and outer bounds), and because Part~III's contains material that is not a restatement (its $L^1$ consequence \eqref{ied:log:eq:L1input} and its comparison with the iterated Bell report). \textup{(Wording corrected after the independent check of 6~October 2026; the first wording read ``because they contain material that is not a restatement (Part~II's bound on inverse orbits, Part~III's $L^1$ consequence \dots)''. Part~II's bound on inverse orbits is the estimate $|G^j(u)|\ge c(|u|+j)$ of Part~I's proof of Lemma~\ref{ied:lead:lem:fatou}.)}

\emph{Numbering.} Sections are numbered continuously and statements within sections, so a statement number moves with its section only. Part~I keeps Report~225's Sections~1--11. Report~226's Section~$k$ is Section~$k+11$ here (Sections~12--23): its Theorems~1.1 and~1.2 are Theorems~\ref{ied:corr:thm:diagonal} and~\ref{ied:corr:thm:height}, its Corollary~10.1 is Corollary~\ref{ied:corr:cor:inverse}. Report~228's Section~$k$ is Section~$k+23$ here (Sections~24--36): its Theorem~1.1 is Theorem~\ref{ied:log:thm:crossover}, its Lemma~3.1, Proposition~4.1, Theorems~5.1, 6.1, 8.1 are \ref{ied:log:lem:finitefatou}, \ref{ied:log:prop:hzero}, \ref{ied:log:thm:hankel}, \ref{ied:log:thm:germ}, \ref{ied:log:thm:transfer}, and its Proposition~11.1 and Theorems~11.2--11.3 are \ref{ied:log:prop:inverseroot}, \ref{ied:log:thm:finiteinverse} and~\ref{ied:log:thm:integerprescription}. The manuscripts number equations consecutively through each document; each Part keeps its manuscript's equation numbers, prefixed in Parts~II and~III: Report~226's equation $(k)$ is $(\mathrm{II}.k)$ here and Report~228's is $(\mathrm{III}.k)$. Section~\ref{ied:sec:further} was added in the write; no delivered number moved.

Notes added when the manuscripts were merged are set in small type and tagged \textbf{[write]}, with their date (6~October 2026). They give cross-references between the Parts, mark statements that a later Part extends or answers, mark results proved twice, and record what the write checked. Superseded statements and answered questions are printed as delivered, with a dated note. No delivered statement was changed, and no delivered symbol was renamed. Each Part has a section of open questions near its end: Part~I its Section~\ref{ied:lead:sec:further}, Part~II its Section~\ref{ied:corr:sec:further}, Part~III its Section~\ref{ied:log:sec:questions}; Section~\ref{ied:sec:further} gathers them for the whole report under Vladimir's standing rule of 4~October 2026, with their sources, sketches and what is missing, and adds the question shared with the iterated Bell report.

\begin{writenote}[Part~IV, added 9 October 2026 (second write)]
Since its second write the report has four Parts. Parts~I--III, the front
matter and Section~\ref{ied:sec:further} are the report as merged on
6~October 2026 and independently checked; the abstract and the front matter
describe them alone, and nothing of them changed except this note, the
Part~IV row of the table above, dated notes at the ends of the front-matter
sections ``What is proved, and where'', ``Provenance and merge decisions''
and ``Relation to neighbouring reports and formal projects'', a dated note at
the end of Section~\ref{ied:sec:further}, and the title, author line, date and
PDF fields. Part~IV (Sections~\ref{ied:par:sec:front}--\ref{ied:par:sec:further}
and Appendix~\ref{ied:par:app:checklist}, labels \file{ied:par:}) is the
manuscript \emph{A Common Density for Parabolic Iteration} (8~October 2026,
``Research companion prepared for Vladimir Reshetnikov / Developed with
OpenAI ChatGPT''). From four analytic inputs proved in Part~I here and in the
iterated Bell report it proves the identity
$I(t)=t\,2^{t/3}h(t)$ for every $t>0$ (Theorem~\ref{ied:par:thm:identity}),
which answers item~\ref{ied:q:ibd} of Section~\ref{ied:sec:further}, and the
holomorphy of $h$ on $\Re t>0$, which answers the positive-axis part of
item~\ref{ied:q:analytic}; the germ at zero stays open. It also proves the
geometric depth aggregate (Theorem~\ref{ied:par:thm:aggregate}), the
strict-chain amplitude $C(q)$ holomorphic on $\Re q>0$, Lengyel's constant
$C=\frac12(2\log2)^{(\log2)/3}h(\log2)$ (Corollary~\ref{ied:par:cor:marked},
\eqref{ied:par:eq:Lengyel}), and the endpoint expansion of $I$
(Corollary~\ref{ied:par:cor:endpoint}). These are implications from the four
inputs, which are unrefereed results of this report and of
\file{a139383-iterated-bell-diagonals}; Section~\ref{ied:par:sec:front} prints
that trust boundary, the provenance, the checks, the non-claims and the
letters ($L$ is $\log(1+1/q)$ in Part~IV, $\log n$ in Part~III; $H(n,m)$ is the
labelled iterated Bell array, not a matching function).
\end{writenote}

\section*{What is proved, and where}\phantomsection\label{ied:sec:status}
\addcontentsline{toc}{section}{What is proved, and where}
Here $T=\pi/2$, $B=\frac12-\frac\pi6-\frac{\pi^2}{16}-\frac{\log2}6$, $f(w)=e^w-1$, $P$ is the entire solution of $P(s+1)=e^{P(s)}-1$ whose inverse is the normalized Fatou coordinate $\Phi(w)=-2/w+\frac13\log w+O(w)$, $h$ is the density with $P(s)=\int_0^\infty e^{st}h(t)\,dt$, and $\beta=n/m$.
\begingroup\small
\renewcommand{\theHtable}{ied-status.\arabic{table}}% distinct PDF destination for the uncaptioned longtable
\begin{longtable}{@{}>{\raggedright\arraybackslash}p{0.55\textwidth} >{\raggedright\arraybackslash}p{0.40\textwidth}@{}}
\toprule
Statement & Where\\
\midrule
\endhead
$P$ entire, the Laplace transform of a finite positive measure with every exponential moment, no atom at zero, support $[0,\infty)$; its density $h$ is $C^\infty$ and strictly positive on $(0,\infty)$ & Part~I, Theorem~\ref{ied:lead:thm:main}, Lemma~\ref{ied:lead:lem:fatou}, Proposition~\ref{ied:lead:prop:measure}, Corollary~\ref{ied:lead:cor:smooth}\\
$a_n\sim c(n/T)^nn^{-4/3}$, $c=T^{1/3}e^{-B}h(1)>0$, with an absolutely convergent Fourier formula for $h(1)$ (Kot\v{e}\v{s}ovec's conjectured form on A290354) & Part~I, Theorem~\ref{ied:lead:thm:main}\\
$A(n,m)\sim h(\beta)m^{-1}(m/T)^nm^{-\beta/3}T^{\beta/3}e^{-\beta B}$ uniformly for $\beta$ in a compact subset of $(0,\infty)$; ratio $e^k$ for fixed shifts; floor phase $e^{-\{\lambda n\}/\lambda}$; Gaussian-shaped profile for $|k|\le C\sqrt n$ & Part~I, Theorem~\ref{ied:lead:thm:proportional}, Corollaries~\ref{ied:lead:cor:shifts} and~\ref{ied:lead:cor:window}\\
Lambert-$W$ inverse $n=X+(\frac43\log X-\log c)/q+o(1/\log X)$, nearest-integer recovery at exact sequence values (no onset) & Part~I, Corollary~\ref{ied:lead:cor:inverse}; sharpened in Part~II, Corollary~\ref{ied:corr:cor:inverse}\\
Subsidiary forcing through fifth scaled order; two further clock terms $M,N$; the clock constant $\kappa=\frac{\pi^3}{96}+\frac{\pi^2}{12}-\frac{13\pi}{48}-\frac{11}{36}-\frac G6-\frac{\pi\log2}{24}$ ($G$ Catalan) & Part~II, Lemma~\ref{ied:corr:lem:forcing}, Section~\ref{ied:corr:sec:clocks}, \eqref{ied:corr:eq:kappa}\\
$a_n=c(n/T)^nn^{-4/3}\{1+(-\frac1{18}\log^2n+A\log n+C)/n+O(n^{-7/5})\}$, $A,C$ exact in $h(1),h'(1),h''(1)$, $\kappa$ & Part~II, Theorem~\ref{ied:corr:thm:diagonal}\\
The uniform first correction $\mathcal Q_\beta(d_m)/m$ for $\beta$ in a compact subset of $(0,\infty)$ & Part~II, Theorem~\ref{ied:corr:thm:height}\\
Finite Fatou expansions to every order with rational $a_j$ and differentiated sectorial remainders & Part~II, Lemma~\ref{ied:corr:lem:fatoufinite} (to $w^2$); Part~III, Lemma~\ref{ied:log:lem:finitefatou} (every order)\\
$h(0)=2$; exact Hankel representation of $h$ near zero; $t^{t/3}h(t)$ has a one-sided $C^\infty$ germ; power-log expansion of $h$ at zero ($h$ is not $C^1$ at $0$); $\log(h/2)$ coefficients $c_j=\rho_j-\zeta(j)/(j3^j)$, $\rho_j\in\mathbb Q$ & Part~III, Proposition~\ref{ied:log:prop:hzero}, Theorems~\ref{ied:log:thm:hankel} and~\ref{ied:log:thm:germ}, Corollary~\ref{ied:log:cor:densityseries}, Proposition~\ref{ied:log:prop:rhozeta}\\
Transfer $mA(n,m)r_m^n=h(\beta)+O(\log m/(m\beta^{30}))$ for all $n$; relative equivalent on $(\log m)^{-D}\le\beta\le B_0$ & Part~III, Theorem~\ref{ied:log:thm:transfer}\\
$A(n,m)/(2T^{-n}m^{n-1})\to e^{-1/(3\lambda)}$ for $m/(n\log n)\to\lambda$, every fixed order in $1/(\lambda\log n)$, $c_2=(9-\pi^2)/108$; Kaneiwa's normalization shifts $K$ to $K+1$ & Part~III, Theorem~\ref{ied:log:thm:crossover}, Section~\ref{ied:log:sec:kaneiwa}\\
Smooth suppression-ratio inverses, rounded height prescriptions, observed-ratio inverses & Part~III, Proposition~\ref{ied:log:prop:inverseroot}, Theorems~\ref{ied:log:thm:finiteinverse}, \ref{ied:log:thm:integerprescription}\\
\midrule
\multicolumn{2}{@{}l}{\emph{Open in all three Parts}}\\
\multicolumn{2}{@{}p{0.96\textwidth}@{}}{Validated digits of $c$, $h(1)$, $h'(1)$, $h''(1)$; corrections beyond the first and an all-orders theorem; effective constants, onsets and certified thresholds; the regime $n/m\to\infty$ and the large-argument behaviour of $h$; real analyticity of $h$ on $(0,\infty)$ and analyticity of the germ at zero; polynomially small ratios and the transition to Kaneiwa's fixed-degree regime; other arrays and initial functions; the conjectured identity $I(\beta)=\beta2^{\beta/3}h(\beta)$ with the iterated Bell amplitude. See Section~\ref{ied:sec:further}.}\\
\bottomrule
\end{longtable}
\addtocounter{table}{-1}% the uncaptioned longtable must not consume a table number
\endgroup
The finite computations of the three companions (exact integer recurrences, exact rational formal coordinates, symbolic clock identities) check identities, normalizations and coefficient algebra; no asymptotic statement rests on them, and every decimal value of $c$, $h(1)$, $h'(1)$, $h''(1)$ in the report is an uncertified floating diagnostic.
\begin{writenote}[Part~IV and this table, 9 October 2026]
Two entries of the last row are now proved in Part~IV, conditionally on its
four inputs (Section~\ref{ied:par:sec:inputs}): the identity
$I(\beta)=\beta2^{\beta/3}h(\beta)$ with the iterated Bell amplitude, for
every $\beta>0$ (Theorem~\ref{ied:par:thm:identity}), and real analyticity of
$h$ on $(0,\infty)$, indeed holomorphy on $\Re t>0$
(\eqref{ied:par:eq:hhol}). Analyticity of the germ at zero, and every other
entry of the row, stay open. Part~IV adds: the geometric depth aggregate
$\sum_me^{-Lm}H(n,m)$ of the labelled array (Theorem~\ref{ied:par:thm:aggregate}),
the strict-chain amplitude $C(q)$ on $\Re q>0$ and Lengyel's constant in terms
of $h(\log2)$ (Corollary~\ref{ied:par:cor:marked}), the endpoint expansion of
$I$ (Corollary~\ref{ied:par:cor:endpoint}) and density formulas for the
constants $K_1$, $K_2$ of the chain-length cumulants
(Corollary~\ref{ied:par:cor:K}).
\end{writenote}

\section*{The OEIS statements, and the inverses}\phantomsection\label{ied:sec:oeis}
\addcontentsline{toc}{section}{The OEIS statements, and the inverses}
The entries as read by the write on 5~October 2026 (A290354 revision \#26 of June 21, 2018; A290353 revision \#28 of September 3, 2026):
\begin{itemize}
\item \textbf{A290354} (Alois P. Heinz, July 28, 2017): ``a(n) is the n-th term of the n-th Euler transform of the sequence with g.f. 1+x'', also ``the number of unlabeled rooted trees with exactly n leaves, all in level n''; formula ``a(n) = A290353(n,n)'' and the formula line, verbatim,
\begin{center}\small\texttt{Conjecture: a(n) \textasciitilde{} c * 2\textasciicircum{}n * n\textasciicircum{}(n-4/3) / Pi\textasciicircum{}n, where c = 4.4923... - \_Vaclav Kotesovec\_, Aug 14 2017}\end{center}
Since $2^nn^{n-4/3}/\pi^n=(n/T)^nn^{-4/3}$, Theorem~\ref{ied:lead:thm:main} proves this form, with $c=T^{1/3}e^{-B}h(1)>0$ identified analytically. It does not certify the digits: the value $c\approx4.4923299$ (Part~I, Section~\ref{ied:lead:sec:computation}; Part~II's diagnostics give the same) is a floating-point evaluation without error control, and agrees with the OEIS's $4.4923\ldots$ only numerically, as Part~I says. The entry's b-file covers $n=0..414$; the manuscripts compare only the 22 displayed terms ($n=0..21$) and an independently generated fixture through $414$, as they state. The entry cross-references A139383, the diagonal of the labelled iterated Bell numbers (see ``Relation to neighbouring reports'').
\item \textbf{A290353} (Alois P. Heinz, July 28, 2017): ``Square array A(n,k), n>=0, k>=0, read by antidiagonals, where column k is the k-th Euler transform of the sequence with g.f. 1+x''; ``A(n,k) is the number of unlabeled rooted trees with exactly n leaves, all in level k''; generating function of column $k>0$: $\prod_{j>0}1/(1-x^j)^{A(j,k-1)}$. Its $k$ is the manuscripts' height $m$ (Part~II writes $b_{n,m}$). The entry's cross-references name rows $2,3,4$ as A001477, A000217, A000330, that is $A(2,m)=m$, $A(3,m)=m(m+1)/2$, $A(4,m)=m(m+1)(2m+1)/6$; these are the coefficients $[q^2]F_h$, $[q^3]F_h$, $[q^4]F_h$ computed in Part~II's Lemma~\ref{ied:corr:lem:forcing} (checked by the write against the exact recurrence for $m\le12$). The entry carries no asymptotic statement.
\end{itemize}
\emph{The inverses.} Part~I's Corollary~\ref{ied:lead:cor:inverse} and Part~II's Corollary~\ref{ied:corr:cor:inverse} start from $X\log(X/T)=Y$. This is the factorial core \texttt{p0:prop:factorial-core} of the transseries volume \cite{TSvol} with $\kappa=1$, its $d$ equal to $-\log T$ and $L=Y$: its solution $\exp(W_0(Y/T)+\log T)=Te^{W_0(Y/T)}$ equals $Y/W_0(Y/T)=X$; after taking logarithms it is the Lambert core \texttt{p0:thm:lambert-core} with $a=b=1$ and $L=\log(Y/T)$ in the variable $\log X-\log T$ \textup{($L$ named after the independent check of 6~October 2026)}. The corrections $X+\delta_0$ (Part~I) and $X+\delta_0+\delta_1/X$ (Part~II) are proved directly, by Taylor expansion of an explicit smooth model and the mean-value theorem at the true integer; they parallel the first terms of \texttt{p0:thm:perturbed-inversion} but are not instances of it, because the perturbation $-\frac43\log x+\log c+Q(\log x)/x$ is not a small analytic perturbation in that theorem's sense. The nearest-integer recovery at exact values $y=a_n$ resembles part~(3) of \texttt{p0:thm:staircase}, but is not an instance either: its smooth root $x(y)$ solves a model equation and is not an admissible interpolation of $(a_n)$, and the manuscripts state no threshold for a general $y$ between two terms. Part~III's inverses (Section~\ref{ied:log:sec:inverse}) involve no Lambert~$W$: the model $\tau=\chi_n\beta-C(\beta)$ has the linear core $\chi_n\beta=\tau$ (the case $b=0$ of \texttt{p0:thm:lambert-core}), perturbed by a germ $C$ known to be $C^\infty$ but not known to be analytic, so its finite reversion (Theorem~\ref{ied:log:thm:finiteinverse}) is proved by Taylor's theorem and the mean-value theorem, again an analogue rather than an instance of \texttt{p0:thm:perturbed-inversion}; its rounded height prescriptions are not instances of \texttt{p0:thm:staircase} (no monotonicity of $R(n,m)$ in $m$ is available). None of the manuscripts claims novelty for the inversion mechanics.

\section*{Notation across the three Parts}\phantomsection\label{ied:sec:notation}
\addcontentsline{toc}{section}{Notation across the three Parts}
Common to all Parts: $F_h$ (the height index is $h$ in Parts~I--II and $m$ in Part~III's recurrence), $f(w)=e^w-1$, $T=\pi/2$, $B$, $\Phi$, $P$, the density $h$ and its measure $\nu$, $\beta=n/m$, the radius $r$, and the contour parameter $s$. Part~II's $b_{n,m}$ is Parts~I and~III's $A(n,m)$. Part~II's $d_\star=\frac13\log T-B$, so its $D_n=n-\frac13\log n+d_\star$ is Part~I's $D_n$, and Part~III's $d_m$ is Part~II's $d_m$. The symbols below change meaning between the Parts or inside one; each Part's text fixes its meaning there, and each Part opens with a short table of its own readings. No symbol was renamed.
\begingroup\small
\renewcommand{\theHtable}{ied-notation.\arabic{table}}% distinct PDF destination for the uncaptioned longtable
\begin{longtable}{@{}>{\raggedright\arraybackslash}p{0.13\textwidth} >{\raggedright\arraybackslash}p{0.83\textwidth}@{}}
\toprule
Symbol & Meanings\\
\midrule
\endhead
$h$ & The density of $\nu$ (all Parts) and the iteration height in $F_h$, $R_h$, $E_h$ (Parts~I, II; Part~II's $t=hz$). \emph{Tempting false reading:} in $F_h(z)$ and $[q^2]F_h=h$ the letter is the height, not the density.\\
$K$ & Parts~I, II: an integer entrance index ($M_0\le K\le n$ in Lemma~\ref{ied:lead:lem:entrance}; $K=\lfloor\sqrt n\rfloor$ in Proposition~\ref{ied:lead:prop:local} and \eqref{ied:corr:eq:scales}; $K=\lceil At\rceil$ in Proposition~\ref{ied:lead:prop:tail}); Part~II also $K_0$. Part~III: the constant $K=\frac{1-\gamma+\log(T/2)}3-B$ \eqref{ied:log:eq:constants}, shifted to $K+1$ by Kaneiwa's normalization; $R_{\mathrm K}$ is Kaneiwa's ratio. \emph{False reading:} Part~III's $K$ is not an index.\\
$\kappa$ & Parts~II, III: the clock constant \eqref{ied:corr:eq:kappa}, $\kappa\approx-0.2543$. Part~I: a small fixed $\kappa>0$ bounding $K\le\kappa n$ (before \eqref{ied:lead:eq:entryphase}).\\
$\rho_j$ & Part~II: $\rho_1=h'(1)/h(1)$, $\rho_2=h''(1)/h(1)$ and $\rho_j(\beta)=h^{(j)}(\beta)/h(\beta)$ \eqref{ied:corr:eq:rhos}. Part~III: the rational numbers $\rho_j$, $j\ge2$, of \eqref{ied:log:eq:cs}. \emph{False reading:} Part~III's $\rho_2=1/12$ is not Part~II's $\rho_2\approx-0.153$.\\
$c$, $C$, $\mathcal C$ & $c=T^{1/3}e^{-B}h(1)$ the amplitude (Parts~I, II). Part~I: $C_z(w)$ the clock \eqref{ied:lead:eq:clock}, $C_\lambda$ \eqref{ied:lead:eq:floor}, $C_0$ a fixed bound (Lemma~\ref{ied:lead:lem:subsidiary}). Part~II: $C$ the constant term of the correction \eqref{ied:corr:eq:AC}, $\mathcal C_z(w)$ the extended clock \eqref{ied:corr:eq:C3}, $C_0$ a real-width bound, $C_{30}$. Part~III: $C_0=1-\gamma-\log2$, $\mathcal C$ Catalan's constant (Part~II's $G$), $C(\beta)$ the smooth germ \eqref{ied:log:eq:Cgerm}; $c_j$ the log-density coefficients \eqref{ied:log:eq:cs}. All Parts: $c_k$ the coefficients in the exact recurrence; $C$, $c$ generic constants. The iterated Bell report's $C\approx2.8654$ is a different amplitude.\\
$G$, $\mathcal G$, $g$ & Part~I: $G(u)=2/\log(1+2/u)$ (Lemma~\ref{ied:lead:lem:fatou}), Part~III's $\mathcal G(u)$; $G_n(s)=F_n(T/(D_n-s))$ \eqref{ied:lead:eq:local}; $g(w)=\log(1+w)$ (also Part~II). Part~II: $G$ Catalan's constant; $G_n(s)=F_n(re^{s/n})$ (exponential, not Part~I's rational parametrization); $G_3(y)$. Part~III: $G_m(s)=F_m(r_me^{s/m})$; $g(t)=t^{t/3}h(t)$ the germ.\\
$L$ & Parts~I, II: the clock function $L(y)$ \eqref{ied:lead:eq:JL}, and a fixed integer $L$ in Propositions~\ref{ied:lead:prop:local} and~\ref{ied:corr:prop:correctedlocal}; Part~I also $L(x)=x\log(x/T)$ in the proof of Corollary~\ref{ied:lead:cor:inverse}. Part~III: $L=\log n$.\\
$W$ & The positive Lambert function (Parts~I, II). Part~III: $W(x)=P(-x)$ (proof of Theorem~\ref{ied:log:thm:hankel}); Part~III uses no Lambert function.\\
$\lambda$ & Part~I: $m=\lfloor\lambda n\rfloor$ (Corollary~\ref{ied:lead:cor:shifts}), so $\lambda\approx m/n$. Part~III: $m/(n\log n)\to\lambda$. \emph{False reading:} Part~III's $\lambda$ is not a slope $m/n$. The iterated Bell report's $\lambda$ is $m/n$.\\
$\alpha$, $\beta$ & $\beta=n/m$ in all Parts; Part~I also uses $0<\alpha<\beta<\pi$ for angles (proof of Lemma~\ref{ied:lead:lem:fatou}), $\beta_0,\beta_1$ interval ends, $\alpha_n=\{\lambda n\}$. Part~II: $\alpha=4/3$ (Section~\ref{ied:corr:sec:inverse}), $\alpha=n-H(z)-\kappa z$ (proof of Proposition~\ref{ied:corr:prop:correctedlocal}). Part~III: $B_0$ the upper end.\\
$H$ & Parts~I, II: the matching function $H(z)=T/z+\frac13\log(1/z)+B$. Part~III: $H=H_{\theta,a}$ a Hankel contour, $H_j(X)$ the log-polynomials of \eqref{ied:log:eq:densityseries}. The iterated Bell report's $H(n,m)$ is its array.\\
$a$, $b$ & $a_n$ the diagonal (all). Part~II: $a(y)=(1+y^2)/2$ (Part~I's $A(y)$), $b(y,t)$, $u$, $v$ step coefficients, $b_{n,m}$ the array. Parts~I--III: $b_j=\sum_{d\mid j}d[z^d]F_h$. Part~III: $a_j$ the rational Fatou coefficients (Part~I's Section~\ref{ied:lead:sec:computation} lists the first three), $b_j$ the germ's Taylor coefficients. Part~I: $b$ a small constant (Proposition~\ref{ied:lead:prop:tail}).\\
$A$ & $A(n,m)$ the array (Parts~I, III). Part~II: $A$ the $\log n$ coefficient \eqref{ied:corr:eq:AC}. Part~I: $A(y)$, and $A$ a large constant (Proposition~\ref{ied:lead:prop:tail}). Part~III: $\mathcal A(w)$ \eqref{ied:log:eq:Acoord}.\\
$d$, $D$ & Part~I: $d=n-h$ (Lemma~\ref{ied:lead:lem:real}), $D_n$. Part~II: $d=-\frac13\log n+d_\star$, $d_m$, $d_y=1+y^2$, $D_n=n+d$. Part~III: $d_m$, $D_N(u)$, $D_t(v)$, $D$ the band exponent \eqref{ied:log:eq:logband}.\\
$R$, $r$ & Parts~I, II: $R_h(z)=\sum_{j\ge2}F_h(z^j)/j$; Part~II also $R_n(s)$. Part~III: the ratio $R(n,m)$ \eqref{ied:log:eq:R}, $R_2$, $R_{\mathrm K}$, the interval $\mathcal R$, and a radius $R$. $r$, $r_m$, $r_n$ radii (all); Part~III also a target ratio $r$.\\
$Q$ & Part~I: $Q(u)=u+\frac13\log u$ (proof of Lemma~\ref{ied:lead:lem:fatou}), $Q_m(t)$ (proof of Theorem~\ref{ied:lead:thm:proportional}). Part~II: $Q_2$, $Q_3$, $\mathcal Q_\beta(d)$, $Q_n(s)$, $Q(u)=-u^2/18+Au+C$. Part~III: $Q_m(s)$, Part~II's $Q_n$ at height $m$.\\
$M$, $N$ & Part~II: the clock functions $M(y)$, $N(y)$; also $M=n^{1/20}$, $M_n$, $M_0$. Part~I: constants $M$, $M_0$. Part~III: $M=m^{1/20}$; $N$ a truncation order.\\
$q$ & Part~I: the entrance constant $q$, $q=\log(X/T)+1$, $q_m$. Part~II: the derivative order $q$ ($=30$), $q_X$. Part~III: a truncation degree $q$.\\
$E$ & Parts~I, II: $E_h$ the clock error. Part~III: $E(m,\beta)$ the transfer error \eqref{ied:log:eq:relative}.\\
$\delta$, $\tau$ & Part~I: $\delta$ the tail range, $\tau$ a first crossing time (Lemma~\ref{ied:lead:lem:real}). Part~II: $\delta_0,\delta_1$ inverse coefficients. Part~III: $\tau=-\log r$.\\
$X$, $y$ & Parts~I, II: $X=Y/W(Y/T)$; $y=w/z$ the clock ratio; $y=a_n$ in the inverses. Part~III: $X=\log t$; $y=m/n$, $y_*$, $y_J$ height variables.\\
$\ell$ & Part~II: $\ell=\log c$. Part~III: $\ell_n=E_{n-1}/(n-1)!$ Kaneiwa's coefficient.\\
$I$, $\Psi$ & Part~II: $I$ a compact interval of ratios. Part~I: $\Psi$ the reciprocal Abel coordinate; Part~III also $\Psi_{\mathrm{IBD}}=-\Phi+\frac13\log2$, the iterated Bell report's coordinate, and $I(\beta)$ that report's amplitude (Section~\ref{ied:log:sec:interfaces}).\\
\bottomrule
\end{longtable}
\addtocounter{table}{-1}% the uncaptioned longtable must not consume a table number
\endgroup

\section*{Provenance and merge decisions}\phantomsection\label{ied:sec:provenance}
\addcontentsline{toc}{section}{Provenance and merge decisions}
Three manuscripts of the session bundle of Reports 1--243. Their author lines read ``Report 225'', ``Report 226'' and ``Report 228'' (PDF author field ``Research report''); none names a person, a tool or an addressee, says it is AI-assisted, or carries ``prepared for private review'' wording. All three are dated 5 October 2026. Report~225 cites the iterated Bell report at commit \texttt{65969409c} and Report~228 at commit \texttt{5fdc0d68a}; both commits exist in the repository, and that report's \file{article.tex} is unchanged at HEAD (blob \texttt{ab5014ec}), so their section numbers refer to its current text. Report~226 pins no commit. They arrived in commit \texttt{60f54ea06} and were placed by \texttt{47fc7a069} (batch~106).
\begin{itemize}
\item \textbf{Part~I}, bundle Report~225, archive \file{Report225.zip} (\file{src/report225.tex}, 639 lines, 16 pages); the base. Contributes the subsidiary-forcing estimate, the nonautonomous clock and its matching constant $B$, the noncircular real entrance, the sectorial Fatou coordinate and entire parametrization, the local limit and angular tails, the positive measure and smooth strictly positive density, the diagonal equivalent, proportional height, fixed and square-root shifts, the floor phase, the Lambert inverse, and Kaneiwa's and Bechtloff Weising's antecedents.
\item \textbf{Part~II}, bundle Report~226, archive \file{Report226.zip} (\file{src/report226.tex}, 715 lines, 17 pages). Its \file{supporting/} holds Report~225's source and PDF byte for byte, and eight of its code and data files are byte copies of Report~225's. Contributes fifth-order forcing, the clock functions $M$ and $N$, the closed form of $\kappa$, the weighted fourth-power defect, the finite Fatou expansion and forced matching, the growing-strip local limit, thirtieth-derivative extraction, the first correction on the diagonal and uniformly in compact ratios, and the sharpened inverse.
\item \textbf{Part~III}, bundle Report~228, archive \file{Report228.zip} (\file{src/report228.tex}, 729 lines, 17 pages). Its \file{supporting/} holds the sources and PDFs of Reports~225 and~226 byte for byte. Contributes the $L^1$ form of Part~II's estimates, finite Fatou expansions to every order, $h(0)=2$, the Hankel representation, the endpoint germ and its rational-plus-zeta coefficients, the transfer to logarithmic bands, the crossover theorem, Kaneiwa's normalization, the smooth inverses, and the local identification with the iterated Bell coordinate together with the conjectured amplitude identity.
\end{itemize}
Where the merge had to choose:
\begin{enumerate}
\item \emph{Order and base.} Dependency order I--II--III, with the base first (see the Guide). Statements that a later Part extends stay as delivered, with dated notes.
\item \emph{Restatements.} Part~II's Section~\ref{ied:corr:sec:interface} restates Part~I's Lemma~\ref{ied:lead:lem:subsidiary}, the clock and entrance (Sections~\ref{ied:lead:sec:subsidiary}--\ref{ied:lead:sec:matching}, \eqref{ied:lead:eq:entryphase}), the coordinate (Lemma~\ref{ied:lead:lem:fatou}), the tail and outer bounds (Proposition~\ref{ied:lead:prop:tail}, Lemma~\ref{ied:lead:lem:outer}) and the density (Proposition~\ref{ied:lead:prop:measure}), and re-derives Corollary~\ref{ied:lead:cor:smooth}. Part~III's Section~\ref{ied:log:sec:interfaces} restates Part~I's coordinate and density and Part~II's absolute estimates, then derives a new $L^1$ bound. The intake dossier suggested replacing these sections by pointers, as the iterated Bell report did for its Part~II; the write printed them in full, as the batch-105 model report did, because nearly every display in them is cited by the later proofs, because Part~II's states as items estimates that Part~I proves only inside proofs, and because Part~III's contains material found nowhere else in the merge (\eqref{ied:log:eq:L1input} and the iterated Bell comparison). \textup{(Wording corrected after the independent check of 6~October 2026; the first wording read ``and because each contains material found nowhere else in the merge''.)} No passage of one manuscript repeats another verbatim.
\item \emph{Front matter and bibliography.} The three title blocks and abstracts are printed at the heads of their Parts; one table of contents replaces three. The three bibliographies are merged into one, keeping every detail of every entry; where two manuscripts describe an entry differently both descriptions are kept. The citation keys were renamed (Report~228's \texttt{oeis} is A290353, Reports~225 and~226's \texttt{oeis} is A290354); Reports~226 and~228's entries for Reports~225 and~226 now point to Parts~I and~II. The write added the entry for the transseries volume \cite{TSvol}.
\item \emph{Part~III's table.} Report~228's \texttt{\textbackslash input\{crossover\_table.tex\}} is replaced by the table's text, byte for byte apart from its two leading comment lines, which are kept as \LaTeX\ comments; the generated file is shipped as \file{data/228-logheight-results-crossover_table.tex}.
\item \emph{Build-specific lines.} The manuscripts' \texttt{\textbackslash pdfmapfile} lines (font-map additions for their TeX~Live build) and Report~228's table-of-contents format tweak are not carried into the merged preamble; they affect no text.
\item \emph{Write additions.} This front matter, the reading-conventions tables of the Parts, Section~\ref{ied:sec:further} with Remarks~\ref{ied:rem:ibdnumbers} and~\ref{ied:rem:labelled}, the dated notes, labels for unlabelled sections, and the label prefixes. Everything else is delivered text.
\end{enumerate}

\begin{writenote}[A stale blob, 9 October 2026]
The statement above that the iterated Bell report's \file{article.tex} is
``unchanged at HEAD (blob \texttt{ab5014ec})'' was true at this report's first
write (\file{84fc1aaed}) and is stale since the same day: the reciprocal note
proposed there was added to that report by commit \file{3f9fc9d4f}, whose blob
\file{2b1d8fbe81} is that report's \file{article.tex} at HEAD and at Part~IV's
pin. The change is an added note; the section numbers the manuscripts cite are
unchanged. Part~IV (Section~\ref{ied:par:sec:front}) is the fourth manuscript;
its pin is ProveIt \file{58175ca455}, at which this report's text is Parts~I--III
as merged and checked.
\end{writenote}

\section*{What was checked, and what was not}\phantomsection\label{ied:sec:trust}
\addcontentsline{toc}{section}{What was checked, and what was not}
The report is unrefereed and not formalized. At intake (dossier of batch~106, 5~October 2026) the three manuscripts were read in full and no step was found false. Checked independently there: the first twelve terms of A290354 and the normalized ratios $a_n(T/n)^nn^{4/3}=4.6013$, $4.5529$, $4.5178$ at $n=40,80,200$ that Part~I quotes, from an own implementation of the recurrence; the closed form of $\kappa$ against a direct quadrature of the defining equation \eqref{ied:corr:eq:M}, $-0.2543481698511$ against $-0.2543481698560$ (difference from the truncated tail); that with Part~II's diagnostic constants the corrected prediction tracks the exact $a_n$ far better than the leading term (residual $1.4\times10^{-4}$ against $2.5\times10^{-2}$ at $n=200$); and $2^{1/3}h(1)$ against the iterated Bell amplitude at $\beta=1$ (Remark~\ref{ied:rem:ibdnumbers}). The delivered suites pass on copies (README). The write checked the rows $A(2,m)$, $A(3,m)$, $A(4,m)$ against Part~II's Lemma~\ref{ied:corr:lem:forcing}, the reversion coefficients \eqref{ied:log:eq:betainverse}--\eqref{ied:log:eq:yinverse}, $c_2$ and $c_3$ from $S(t)$ and the Gamma expansion, the linear coefficient $C_0/3$ and $K$ in \eqref{ied:log:eq:logmaster}, the reduction of Theorem~\ref{ied:corr:thm:height} to Theorem~\ref{ied:corr:thm:diagonal} at $\beta=1$, the floor amplitude $C_\lambda$ of \eqref{ied:lead:eq:floor}, the two orders of the inverse \eqref{ied:corr:eq:inversecoeffs}, and recomputed $c=4.4923298968$ from the shipped diagnostic value of $h(1)$ \textup{(added after the independent check of 6~October 2026: an independent evaluation of $h(1)$ gives $c=4.4923298971$; Remark~\ref{ied:rem:ibdnumbers})}.

What was \emph{not} done: the analytic estimates (the clocks and their defects, the sector bootstraps, the tails, the Hankel deformations) were read and their displayed algebra spot-checked, but they are the manuscripts' proofs, corroborated numerically, not an independent verification. The write did not read Kaneiwa (1980), Bechtloff Weising, Dudko--Sauzin, Prellberg--Mishna or Nagaev--Vakhtel; those citations and their page and theorem numbers are as the manuscripts and their source ledgers give them.

\section*{Relation to neighbouring reports and formal projects}\phantomsection\label{ied:sec:neighbours}
\addcontentsline{toc}{section}{Relation to neighbouring reports and formal projects}
\textbf{The labelled analogue.} \file{SetTheory/Cardinals/docs/reports/generating-functions-and-asymptotics/oeis-sequence-asymptotics/a139383-iterated-bell-diagonals} (labels \texttt{ibd:}) treats $H(n,m)=n![z^n](e^z-1)^{\circ m}$, OEIS A139383 and A261280: a fixed-shift diagonal theorem to every order (its Part~I) and proportional depth $m=\lambda n+\delta$ for $\lambda$ in compact subsets of $(0,\infty)$ (its Part~II), by the same parabolic map $w\mapsto e^w-1$ without subsidiary forcing. No theorem is shared: the arrays differ, and Part~I here says the labelled theorem ``cannot be directly transferred'' (Section~\ref{ied:lead:sub:prior}), because the forcing changes the leading scale ($(n/T)^n n^{-4/3}$ here, $n^{2n-5/6}/(2^{n-1}e^n)$ there) and the clock constant. The two reports share the normalized orbit: Part~III records $\Psi_{\mathrm{IBD}}(w)=-\Phi(w)+\frac13\log2$ and the unproved relation $I(\beta)=\beta2^{\beta/3}h(\beta)$ between that report's amplitude $I$ and the density $h$ here (Section~\ref{ied:log:sec:interfaces}). The values agree to eight significant digits at $\beta=1$ (Remark~\ref{ied:rem:ibdnumbers}), and the relation, if proved, would bear on that report's open endpoint $\lambda\to\infty$ (Section~\ref{ied:sec:further}, item~\ref{ied:q:ibd}, and Remark~\ref{ied:rem:labelled}). That report answers no question asked here, and this one answers none asked there; a reciprocal note is proposed for it.

\textbf{Other neighbours.} No other report of the collection treats A290353 or A290354 (repository search, 6~October 2026; A290354 occurs elsewhere only in a cross-reference line of the iterated Bell report's OEIS data). The batch-106 reports \file{a252782-diagonal-euler-transforms} (one Euler transform with weights $j^n$) and \file{a005121-strict-partition-chains} (the depth sum $\sum_m2^{-m-1}H(n,m)$ of the labelled array) share only words with this one. The transseries volume supplies the factorial core of which the leading inverses are instances (``The OEIS statements, and the inverses''). Placement in the collection confers no formal status: no Lean or Rocq file in the repository treats these sequences, and no statement of this report is formalized.
\begin{writenote}[Part~IV and the neighbours, 9 October 2026]
The two sentences above that no theorem is shared with the iterated Bell
report and that the relation $I(\beta)=\beta2^{\beta/3}h(\beta)$ is unproved
describe Parts~I--III. Part~IV proves that relation for every $\beta>0$ from
inputs of both reports (Theorem~\ref{ied:par:thm:identity}), and gives the
leading amplitude of \file{a005121-strict-partition-chains} (Lengyel's constant)
and its marked amplitude on $\Re q>0$ (Corollary~\ref{ied:par:cor:marked}).
Notes for those two reports, recording what Part~IV answers there (the endpoint
amplitude of $I$ and the exploratory values of $I(1)$ in the iterated Bell
report; the heuristic of \file{spc:q:depth} and, in part, \file{spc:q:disk} in
the chain report), are proposed separately; they are not applied by this write.
\end{writenote}
\clearpage

\part{Diagonal and proportional height asymptotics for iterated Euler transforms}\label{ied:lead:part}
\parteqs{}
\partsource{Bundle Report~225, archive \file{Report225.zip}, delivered as \file{src/report225.tex} (16 pages), the base of the merge. Delivered title block ``Diagonal and proportional height asymptotics for iterated Euler transforms / Report 225 / 5 October 2026''. Printed as delivered, with \textbf{[write]} notes; section, statement and equation numbers are those of the delivery.}
\begin{center}
\small
\begin{tabular}{@{}>{\raggedright\arraybackslash}p{0.46\textwidth} >{\raggedright\arraybackslash}p{0.46\textwidth}@{}}
\toprule
Reading conventions of Part~I & Elsewhere in the report\\
\midrule
$A(n,m)=[z^n]F_m$; $a_n=A(n,n)$ & Part~II: $b_{n,m}$ for $A(n,m)$\\
$K$ an integer entrance index; $\kappa$ a small constant in $K\le\kappa n$ & Part~III: $K$ a constant; Parts~II, III: $\kappa$ the clock constant\\
$L(y)$ the clock function; $C_z(w)$ the clock; $E_h$ & Part~II: $\mathcal C_z(w)$; Part~III: $L=\log n$, $E(m,\beta)$\\
$G(u)$, $g(w)=\log(1+w)$, $G_n(s)=F_n(T/(D_n-s))$ & Part~II: $G$ Catalan, $G_n(s)=F_n(re^{s/n})$; Part~III: $\mathcal G(u)$, $g(t)=t^{t/3}h(t)$\\
$m=\lfloor\lambda n\rfloor$ (Corollary~\ref{ied:lead:cor:shifts}) & Part~III: $m/(n\log n)\to\lambda$\\
\bottomrule
\end{tabular}
\end{center}
\begin{partabstract}
We study the number $a_n$ of unlabeled rooted trees with $n$ leaves, all at level $n$, equivalently the diagonal of the iterated Euler transform in OEIS A290354. We prove
\[
 a_n\sim c\left(\frac{2n}{\pi}\right)^n n^{-4/3},\qquad c>0,
\]
and identify $c$ by a normalized repelling parametrization of $w\mapsto\e^w-1$. The main estimate is a uniform intermediate-angle bound, obtained from an explicit nonautonomous clock, a positive-coefficient Cauchy entrance, and a forced sector bootstrap. A positive Laplace measure identifies the coefficient amplitude and proves its strict positivity without assuming a density expansion at zero. Uniform proportional-height asymptotics for the array A290353, fixed and square-root height shifts, and a Lambert $W$ inverse corollary are included. Numerical diagnostics give $c\approx4.4923$, agreeing with Kotesovec's conjecture; the decimal computation is not a validated enclosure. This is an unrefereed mathematical research report, not a proof-assistant formalization.
\end{partabstract}
\section{Definition and statement}\label{ied:lead:sec:statement}
Let
\begin{equation}\label{ied:lead:eq:rec}
 F_0(z)=z,\qquad
 F_{h+1}(z)=\exp\left(\sum_{j\ge1}\frac{F_h(z^j)}j\right)-1.
\end{equation}
Each $F_h$ is analytic in $|z|<1$, has nonnegative coefficients, vanishes at zero, and has linear coefficient one. These facts follow by induction; the sum over $j$ converges normally on compact subdisks because $F_h(z)=\Oh(z)$ at zero. For $n\ge1$ put $a_n=[z^n]F_n(z)$ and set $a_0=1$ separately. The constant-term exception matches the OEIS convention and has no asymptotic effect.

For clarity, write $A(n,m)=[z^n]F_m(z)$ for $n\ge1$, and put $A(0,m)=1$ by convention. This is the square array A290353~\cite{A290353}, with degree $n$ and height $m$; the sequence A290354 is its main diagonal. A tree whose leaves all have level $m+1$ is a root carrying a nonempty unordered multiset of trees whose leaves all have level $m$. If there are $[z^d]F_m(z)$ types of a subtree with $d$ leaves, the ordinary multiset generating function is
\begin{equation}\label{ied:lead:eq:eulerproduct}
 1+F_{m+1}(z)=\prod_{d\ge1}(1-z^d)^{-[z^d]F_m(z)}.
\end{equation}
Taking its formal logarithm gives \eqref{ied:lead:eq:rec}. This also proves integrality and explains why there is no factorial normalization. The first diagonal values are
\[
1,\ 1,\ 2,\ 6,\ 30,\ 170,\ 1337,\ 12166,\ 133476,\ 1676364.
\]
The separate constant row is a generating-function convention, rather than a claim about a nonempty tree with no leaves.


Set
\begin{equation}\label{ied:lead:eq:constants}
 T=\frac\pi2,\qquad
 B=\frac12-\frac\pi6-\frac{\pi^2}{16}-\frac{\log2}{6},\qquad
 D_n=n-\frac13\log n+\frac13\log T-B.
\end{equation}
The logarithms used near the positive axis have their real determination there.

\begin{theorem}\label{ied:lead:thm:main}
There is an entire function $P$ satisfying
\begin{equation}\label{ied:lead:eq:P}
 P(s+1)=\e^{P(s)}-1,
\end{equation}
whose inverse coordinate near zero is normalized by
\[
 \Phi(w)=-\frac2w+\frac13\log w+\Oh(w),\qquad
 \Phi(\e^w-1)=\Phi(w)+1.
\]
It is the Laplace transform of a finite positive measure $\nu$ on $[0,\infty)$ with every exponential moment. This measure has a smooth, strictly positive density $h$ on $(0,\infty)$. Moreover,
\begin{equation}\label{ied:lead:eq:asymptotic}
 a_n\sim c(n/T)^n n^{-4/3},\qquad
 c=T^{1/3}\e^{-B}h(1)>0.
\end{equation}
For every real $\sigma$,
\begin{equation}\label{ied:lead:eq:constantintegral}
 h(1)=\frac{\e^{-\sigma}}{2\pi}
 \int_{-\infty}^{\infty}P'(\sigma+\ii t)\e^{-\ii t}\dd t,
\end{equation}
and this integral is absolutely convergent.
\end{theorem}

The normalization of $\Phi$ is essential: translating a Fatou coordinate translates $P$ and changes the density amplitude. Formula \eqref{ied:lead:eq:constants} is tied to the zero constant term in the displayed normalization.

\subsection{Prior work and the scope of the result}\label{ied:lead:sub:prior}
The OEIS entry attributes \eqref{ied:lead:eq:asymptotic}, with $c=4.4923\ldots$, to V\'aclav Kot\v{e}\v{s}ovec on 14 August 2017 and labels it a conjecture~\cite{A290354}. Kaneiwa's 1980 paper~\cite{Kaneiwa}, Propositions 1 and 2 on printed page 19, proves that the unlabeled $r$-fold partition number $p(r;n)$ is a polynomial of degree $n-1$ in $r$ and identifies its leading coefficient through $\tan x+\sec x$. His indexing gives
\[
 p(r;n)=[z^n]F_{r+1}(z),\qquad a_n=p(n-1;n)\quad(n\ge1).
\]
Thus polynomiality, the tangent leading profile, and its $\pi/2$ pole are prior results. Fixed-$n$ polynomial information does not alone justify the growing-$n$ diagonal limit. Kaneiwa credits Bell's 1938 polynomiality result for the separately defined \emph{labeled} function $\widetilde p(r;n)$ on page 18; that sentence should not be read as attributing his unlabeled Proposition 1 to Bell.

Bechtloff Weising's higher-order Bell symmetric functions~\cite{BW}, Definition 2.1 and Proposition 2.12, specialize in one variable to $1+F_{m+1}$; in particular $a_n=[m_{(n)}]B_n^{(n-1)}$. This is an algebraic and species connection, not a diagonal asymptotic theorem in that source. The related iterated-Bell report~\cite{IBD} concerns $n![z^n](\e^z-1)^{\circ m}$ and cannot be directly transferred: the subsidiary forcing in \eqref{ied:lead:eq:rec} changes both the leading scale and the clock constant. No general priority claim for parabolic-coordinate methods is made here.

\begin{writenote}[The OEIS statement and the neighbouring report, 6~October 2026]
The front matter quotes A290354's formula line verbatim as read on 5~October 2026: ``Conjecture: a(n) $\sim$ c * 2\textasciicircum n * n\textasciicircum(n-4/3) / Pi\textasciicircum n, where c = 4.4923... - \_Vaclav Kotesovec\_, Aug 14 2017''. Theorem~\ref{ied:lead:thm:main} proves this form with $c$ identified; the digits $4.4923$ are matched only by an uncertified evaluation (Subsection~\ref{ied:lead:sub:notcertified}). Kaneiwa's indexing $p(r;n)=[z^n]F_{r+1}$ is used again in Part~III, Section~\ref{ied:log:sec:kaneiwa}, which compares his leading coefficient with the crossover. The iterated Bell report \cite{IBD} is the labelled analogue; Part~III, Section~\ref{ied:log:sec:interfaces}, records the common normalized orbit and a conjectured identity between its amplitude and the density $h$ here (Section~\ref{ied:sec:further}, item~\ref{ied:q:ibd}).
\end{writenote}

\section{Subsidiary terms and the nonautonomous clock}\label{ied:lead:sec:subsidiary}
Write
\[
 R_h(z)=\sum_{j\ge2}\frac{F_h(z^j)}j,\qquad f(w)=\e^w-1.
\]
The subsidiary estimate below is independent of the size of the main value $F_h(z)$.

\begin{lemma}\label{ied:lead:lem:subsidiary}
Uniformly for $h|q|$ sufficiently small and $|q|$ small,
\begin{equation}\label{ied:lead:eq:smallq}
 F_h(q)=q+hq^2+\Oh\bigl((h+h^2)|q|^3\bigr).
\end{equation}
Consequently, for each fixed $C_0$, uniformly for $h\le C_0/|z|$ as $z\to0$,
\begin{equation}\label{ied:lead:eq:R}
 R_h(z)=\frac{z^2}2+\frac{z^3}3+\frac{hz^4}2+\Oh(|z|^4).
\end{equation}
\end{lemma}
\begin{proof}
For real $0<q<1$, coefficient positivity and the absence of a constant term give
$F_h(q^j)\le q^{j-1}F_h(q)$. Hence the exponent defining $F_{h+1}(q)$ is at most $F_h(q)[-\log(1-q)]/q$. While $F_h(q)\le2q$, the one-step increment is at most $C F_h(q)^2$. The elementary reciprocal comparison, or induction against $q/(1-C'hq)$ with $C'$ fixed large, proves $F_h(q)\le2q$ for $hq\le c$ and $c$ sufficiently small. Positivity gives the corresponding complex bound.

Summing one-step $\Oh(|q|^2)$ increments gives $F_h(q)=q+\Oh(h|q|^2)$. Substituting this back into the recurrence gives the sharper increment
$F_{h+1}(q)-F_h(q)=q^2+\Oh((h+1)|q|^3)$, because the quadratic main term and the leading subsidiary term each contribute $q^2/2$. Summation proves \eqref{ied:lead:eq:smallq}. Apply it at $q=z^j$ for $j\ge2$. The $j=2$ error is $\Oh(h^2|z|^6)=\Oh(|z|^4)$; the baseline tail $j\ge4$ and the remaining corrections are bounded by geometric series. This proves \eqref{ied:lead:eq:R}.
\end{proof}

\begin{writenote}[Extended in Part~II, 6~October 2026]
Part~II's Lemma~\ref{ied:corr:lem:forcing} extends \eqref{ied:lead:eq:R} to fifth scaled order, from the exact coefficients $[q^2]F_h=h$, $[q^3]F_h=h(h+1)/2$, $[q^4]F_h=h(h+1)(2h+1)/6$, which are the rows $n=2,3,4$ of A290353 (checked by the write against the exact recurrence for $h\le12$).
\end{writenote}

Define, in a narrow sector about the positive real axis,
\begin{equation}\label{ied:lead:eq:JL}
 J(y)=2\arctan y-\frac\pi2,\qquad L(1)=0,\qquad
 L'(y)=\frac{y^3/3-y-4/3-2J(y)}{(1+y^2)^2}.
\end{equation}
Put
\begin{equation}\label{ied:lead:eq:clock}
 C_z(w)=\frac{J(w/z)}z+L(w/z).
\end{equation}
The clock domain will be $|\arg(w/z)|\le\eta$, $|w/z|\ge1/2$, with a slightly wider fixed sector available for Taylor segments. In particular $C_z(z)=0$.

\begin{lemma}\label{ied:lead:lem:clock}
Suppose $h\le C_0/|z|$, $w_h=F_h(z)$ remains sufficiently small, and $w_h/z$ remains in the stated clock domain through time $k$. Then, with $E_h=C_z(w_h)-h$,
\begin{equation}\label{ied:lead:eq:defect}
 E_{h+1}-E_h
 =\frac{z^2\bigl(L(w_h/z)-E_h\bigr)}{1+(w_h/z)^2}
   +\Oh(|w_h|^2+|z|^2),
\end{equation}
with uniform constants. Consequently
\begin{equation}\label{ied:lead:eq:clockbound}
 |E_k|\le C\left(\sum_{h<k}|F_h(z)|^2+k|z|^2\right).
\end{equation}
\end{lemma}
\begin{proof}
By Lemma~\ref{ied:lead:lem:subsidiary},
\begin{align}\label{ied:lead:eq:wstep}
 w_{h+1}={}&w_h+\frac{w_h^2+z^2}2+\frac{w_h^3}6
 +\frac{w_hz^2}2+\frac{z^3}3+\frac{hz^4}2
 +\Oh((|w_h|+|z|)^4).
\end{align}
Write $y=w_h/z$, $A(y)=(1+y^2)/2$ and
$B_1(y,hz)=y^3/6+y/2+1/3+hz/2$. The increment of $y$ is
$zA+z^2B_1+\Oh((|w_h|+|z|)^4/|z|)$. Since $J'=1/A$, the coefficient of $z$ in the clock increment is
\[
 \frac{B_1}{A}-\frac y2+A L'(y)=\frac{hz-J(y)}{1+y^2}.
\]
But $hz-J(y)=z(L(y)-E_h)$, proving the cancellation in \eqref{ied:lead:eq:defect}.

For completeness the remainder is analytic and uniform. Derivatives of $J(w/z)/z$ of order $r\ge1$ are $\Oh((|w|+|z|)^{-r-1})$, and derivatives of $L(w/z)$ are $\Oh((|w|+|z|)^{-r})$. The displacement in \eqref{ied:lead:eq:wstep}, divided by $|w|$, is $\Oh(|w|+|z|)$ because $|w/z|$ is bounded below. Nested clock cones therefore keep the full Taylor segment away from $w=\pm\ii z$. The omitted recurrence terms, the mixed quadratic Taylor terms, and the cubic Taylor remainder all contribute $\Oh((|w|+|z|)^2)$ to the clock increment.

Finally $L(y)/(1+y^2)$ is bounded on the clock cone and $|1+y^2|$ is bounded below. The multiplicative discrete Gronwall factor is at most $\exp(Ck|z|^2)$, uniformly bounded in the stated range. Since $E_0=0$, this gives \eqref{ied:lead:eq:clockbound}.
\end{proof}

\section{Matching and the positive real orbit}\label{ied:lead:sec:matching}
\begin{lemma}\label{ied:lead:lem:matching}
Uniformly in a fixed narrow sector about the positive real axis,
\begin{equation}\label{ied:lead:eq:Lasym}
 L(y)=\frac13\log y+B+\Oh(y^{-2})\qquad(|y|\to\infty).
\end{equation}
Consequently, when $|w/z|$ is sufficiently large in that sector,
\begin{equation}\label{ied:lead:eq:match}
 C_z(w)=H(z)-\frac2w+\frac13\log w
  +\Oh\left(\frac{|z|^2}{|w|^3}+\frac{|z|^2}{|w|^2}\right),
 \quad H(z)=\frac Tz+\frac13\log\frac1z+B.
\end{equation}
For real $0<x\le w\le\epsilon$, with $\epsilon$ small fixed,
\begin{equation}\label{ied:lead:eq:realcompare}
 H(x)-C_x(w)\asymp\frac1w,\qquad C_x'(w)>0.
\end{equation}
\end{lemma}
\begin{proof}
The derivative in \eqref{ied:lead:eq:JL} is $1/(3y)+\Oh(y^{-3})$; in particular there is no $y^{-2}$ term. To compute its constant, integrate the part $(y^3/3-y)/(1+y^2)^2$ first. After subtraction of $(\log y)/3$ its contribution from $1$ to infinity is $-(\log2)/6-1/3$. In the rest set $u=\arctan y$. Its contribution is
\[
 \int_{\pi/4}^{\pi/2}(\pi-4/3-4u)\cos^2u\dd u
 =\frac56-\frac\pi6-\frac{\pi^2}{16}.
\]
Their sum is $B$ in \eqref{ied:lead:eq:constants}. The arctangent expansion gives \eqref{ied:lead:eq:match}.

For real $w\ge x$, the arctangent identity gives
\[
 H(x)-C_x(w)=\frac2x\arctan\frac xw-\frac13\log w
              +\Oh((x/w)^2).
\]
The first term is between fixed positive multiples of $1/w$; the remaining terms are smaller for small $w$, uniformly for $x/w\le1$. The expansion of $L$ may be bounded on all $y\ge1$ by enlarging its constant on a compact interval. Finally the main derivative $2/(x^2+w^2)$ dominates the $\Oh(1/(x+w))$ derivative of the correction for small $\epsilon$. This proves \eqref{ied:lead:eq:realcompare}.
\end{proof}

\begin{lemma}\label{ied:lead:lem:real}
Fix $\sigma\in\R$ and let $r=T/(D_n-\sigma)$. There are constants $M_0,c_0,c,C>0$ such that, for all sufficiently large $n$, if $d=n-h\ge M_0$ and $x>0$ satisfies $|x-r|\le c_0d/n^2$, then every real iterate through $h$ stays in the small-clock range and
\begin{equation}\label{ied:lead:eq:realbound}
 \frac cd\le F_h(x)\le\frac Cd.
\end{equation}
\end{lemma}
\begin{proof}
We have $H(r)=n-\sigma+\Oh(\log n/n)$ and $H'(x)=\Oh(n^2)$ when $x\asymp n^{-1}$. Make $c_0$ small, then $M_0$ large compared with $|\sigma|$. Uniformly in the displayed range,
\begin{equation}\label{ied:lead:eq:gap}
 d/2\le H(x)-h\le2d.
\end{equation}
It remains to justify that the small-clock hypotheses are not circular. Let $\tau$ be the first time $F_\tau(x)\ge\epsilon$, if it occurs by time $h$. Start with a slightly larger admissible clock radius, so that one-step overshoot gives $F_\tau(x)\le2\epsilon$. Positivity gives
\[
 F_{j+1}(x)-F_j(x)\ge\frac12F_j(x)^2,
 \qquad \sum_{j<\tau}F_j(x)^2\le2(F_\tau(x)-x).
\]
Lemma~\ref{ied:lead:lem:clock} therefore gives $E_\tau=\Oh(\epsilon+\tau x^2)$; here $\tau\le h\le n$ and $x\asymp n^{-1}$. Lemma~\ref{ied:lead:lem:matching} and monotonicity of $C_x$ imply
$\tau\ge H(x)-C/\epsilon$. If $M_0$ is large enough, this contradicts \eqref{ied:lead:eq:gap}. Thus no such first crossing occurs. Apply \eqref{ied:lead:eq:clockbound} again through $h$ and use \eqref{ied:lead:eq:realcompare}; the clock error is $\Oh(F_h(x)+nx^2)$ and can be absorbed into the $1/F_h(x)$ comparison. Equation \eqref{ied:lead:eq:gap} now gives \eqref{ied:lead:eq:realbound}.
\end{proof}

\begin{lemma}\label{ied:lead:lem:entrance}
With $r$ as in Lemma~\ref{ied:lead:lem:real}, there is $q>0$ such that for every integer $M_0\le K\le n$ and every $|z-r|\le qK/n^2$, the orbit through $n-K$ lies in the clock cone, satisfies $|F_h(z)|\asymp(n-h)^{-1}$, and obeys
\begin{equation}\label{ied:lead:eq:entranceclock}
 C_z(F_{n-K}(z))=n-K+\Oh(K^{-1}+n^{-1}).
\end{equation}
\end{lemma}
\begin{proof}
At height $h\le n-K$, use the disk about $r$ of radius
$R_h=c_0(n-h)/n^2$. Positivity and Lemma~\ref{ied:lead:lem:real} bound its supremum by $F_h(r+R_h)\le C/(n-h)$. On its inner half, Cauchy's estimate gives
$|F_h'|\le Cn^2/(n-h)^2$. Hence
\[
 |F_h(z)-F_h(r)|\le C|z-r|\frac{n^2}{(n-h)^2}.
\]
Divide by the positive lower bound in \eqref{ied:lead:eq:realbound}. Taking $q\le c_0/2$ and then decreasing it makes the relative error at most $CqK/(n-h)\le Cq$. Also $|z-r|/r\le CqK/n$. Thus $F_h(z)/z$ lies in the required fixed narrow cone, with modulus at least $1/2$, for every $h\le n-K$. The two-sided comparison follows. Finally
$\sum_{h<n-K}|F_h(z)|^2\le C/K$, so Lemma~\ref{ied:lead:lem:clock} gives \eqref{ied:lead:eq:entranceclock}.
\end{proof}

\section{The repelling coordinate with uniform sectors}\label{ied:lead:sec:fatou}
The needed coordinate is constructed here to make the sector and branch choices explicit.
\begin{lemma}\label{ied:lead:lem:fatou}
For each closed angular sector strictly contained in $|\arg w|<\pi$, there is a sufficiently small radius on which an analytic function $\Phi$ satisfies
\begin{equation}\label{ied:lead:eq:fatou}
 \Phi(w)=-2/w+\tfrac13\log w+\Oh(w),\quad
 \Phi'(w)=2/w^2+\Oh(1/w),\quad
 \Phi(f(w))=\Phi(w)+1,
\end{equation}
whenever both arguments lie in the chosen domain. These functions agree on common smaller sectors. There is a local inverse on every sufficiently distant closed subsector of $|\arg(-s)|<\pi$, with
\begin{equation}\label{ied:lead:eq:Pinverse}
 P_{\rm loc}(s)=-2/s+\Oh(\log|s|/|s|^2),\qquad
 P_{\rm loc}'(s)=2/s^2+\Oh(\log|s|/|s|^3).
\end{equation}
For small $w$ in a smaller sector, $P_{\rm loc}(\Phi(w))=w$. The local inverse extends to the entire $P$ in Theorem~\ref{ied:lead:thm:main} and agrees with it on its sector of definition.
\end{lemma}
\begin{proof}
Let $g(w)=\log(1+w)$ and conjugate by $u=2/w$. The resulting map satisfies
\[
 G(u)=\frac2{\log(1+2/u)}=u+1-\frac1{3u}+\Oh(u^{-2}).
\]
For $Q(u)=u+(\log u)/3$,
\begin{equation}\label{ied:lead:eq:Qdefect}
 Q(G(u))-Q(u)=1+\Oh(u^{-2}).
\end{equation}
We give the uniform geometric estimate behind summation of this defect. Fix $0<\alpha<\beta<\pi$. If $|\arg u|\le\alpha$, then
\[
 |u+j|\ge\cos(\alpha/2)(|u|+j),\qquad j\ge0,
\]
and the straight ray has distance at least $c_{\alpha,\beta}(|u|+j)$ from the complement of $|\arg v|<\beta$. Distinguish an inner analyticity radius $R_0$ from a larger starting threshold $R_1$. Bootstrap
\[
 |G^j(u)-(u+j)|\le a(|u|+j),
\]
where $a$ is less than half the preceding radial and angular margins. This gives $|G^j(u)|\ge c(|u|+j)>R_0$ if $R_1$ is large enough. Summing the one-step $\Oh(1/|G^j(u)|)$ error improves the bound to
\[
 |G^j(u)-(u+j)|\le C\left(\frac1{|u|}+\log(1+j/|u|)\right).
\]
Divided by $|u|+j$, this is uniformly $\Oh(R_1^{-1})$ because $\log(1+x)/(1+x)$ is bounded. Increase $R_1$ to obtain a strict improvement and $cR_1>2R_0$. This closes the induction, including times when the straight ray decreases in modulus. Only the angular sector is invariant; the starting-radius boundary need not be.

The defects in \eqref{ied:lead:eq:Qdefect} are now absolutely summable with total $\Oh(1/|u|)$. Therefore
\[
 \Psi(u)=\lim_{m\to\infty}\bigl(Q(G^m(u))-m\bigr)
 =Q(u)+\Oh(1/u)
\]
exists locally uniformly, is analytic, and satisfies $\Psi(G(u))=\Psi(u)+1$ wherever both sides are defined. Nested sectors and thresholds justify shifting the limit. Cauchy estimates on smaller sectors give derivative bounds. Define
\[
 \Phi(w)=-\Psi(2/w)+\frac13\log2.
\]
This proves \eqref{ied:lead:eq:fatou}; $g\circ f$ is the identity near zero.

For the inverse, solve $\Psi(u)=-s+(\log2)/3$. A disk of radius $C\log|s|$ centered at the right side stays inside a slightly larger allowed sector for large $|s|$. On its boundary the logarithmic correction to the identity map is smaller than the radius for $C$ large. Rouch\'e's theorem gives exactly one root, and $\Psi'(u)=1+\Oh(1/u)$ is nonzero. The roots depend analytically on $s$, agree on overlaps by uniqueness, and give \eqref{ied:lead:eq:Pinverse}. If a small $w$ already satisfies $\Phi(w)=s$ in a smaller sector, then $u=2/w=-s+\Oh(\log|s|)$ lies in the same uniqueness disk. Thus $P_{\rm loc}(\Phi(w))=w$; global injectivity of $\Phi$ is neither assumed nor needed.

First restrict $P_{\rm loc}$ to a strict left half-plane $\Rea s<-R_2$ so that $f(P_{\rm loc}(s))=P_{\rm loc}(s+1)$ holds when both parameters lie there. On $\Rea s<m-R_2$, define
$P(s)=f^{\circ m}(P_{\rm loc}(s-m))$. Successive definitions agree on overlaps, so they produce an entire function obeying \eqref{ied:lead:eq:P}. Agreement with $P_{\rm loc}$ on every connected larger inverse sector follows from the identity theorem, since there is agreement on a nonempty open part of the left half-plane. Positivity first far to the left and then forward iteration show $P(s)>0$ for real $s$.
\end{proof}

\begin{writenote}[Finite expansions of $\Phi$ elsewhere in the report, 6~October 2026]
Part~II's Lemma~\ref{ied:corr:lem:fatoufinite} proves $\Phi(w)=-2/w+\frac13\log w-w/36+w^2/540+\Oh(w^3)$ with differentiated sectorial remainders, and Part~III's Lemma~\ref{ied:log:lem:finitefatou} the expansion to every order with rational coefficients $a_j$, both by summing the Abel defect along the inverse orbits constructed in this proof. The coefficients displayed in Subsection~\ref{ied:lead:sub:notcertified} are those $a_1,a_2,a_3$.
\end{writenote}

Combining Lemmas~\ref{ied:lead:lem:matching}, \ref{ied:lead:lem:entrance}, and \ref{ied:lead:lem:fatou} gives the crucial entrance phase. Restrict $M_0\le K\le\kappa n$ with $\kappa>0$ small fixed, so that $|F_{n-K}(z)/z|\asymp n/K$ lies in the large-ratio matching range. Uniformly on the entrance disk,
\begin{equation}\label{ied:lead:eq:entryphase}
 \Phi(F_{n-K}(z))=n-K-H(z)
   +\Oh(K^{-1}+n^{-1}+K^3/n^2).
\end{equation}
Both matching scales below satisfy this restriction.

\section{The local limit and angular tails}\label{ied:lead:sec:local}
\begin{proposition}\label{ied:lead:prop:local}
Locally uniformly for $s\in\C$,
\begin{equation}\label{ied:lead:eq:local}
 G_n(s):=F_n\left(\frac T{D_n-s}\right)\longrightarrow P(s).
\end{equation}
Here $G_n$ is defined on each fixed compact set for all sufficiently large $n$.
\end{proposition}
\begin{proof}
Fix a compact set $S$ and take $K=\lfloor\sqrt n\rfloor$. The point $z=T/(D_n-s)$ differs from the real reference $T/D_n$ by $\Oh_S(n^{-2})$, so Lemma~\ref{ied:lead:lem:entrance} applies. Exactly,
\[
 H\left(\frac T{D_n-s}\right)=n-s+\frac13\log\frac{D_n-s}{n}.
\]
Thus \eqref{ied:lead:eq:entryphase} gives entrance phase $s-K+o(1)$ uniformly on $S$.

Choose a fixed integer $L$ larger than the inverse-sector threshold and $2\max_{s\in S}(|s|+1)$. For $0\le j\le K-L$, the prototype $s-K+j$ lies in a fixed interior left sector, has modulus comparable to $K-j$, and has real part bounded away from zero. Bootstrap phase error at most one. By \eqref{ied:lead:eq:R}, the actual $w$-step differs from $f(w)$ by $\Oh(n^{-2})$. The disk of this radius about the unforced image lies in the coordinate domain: the image has distance at least $c/(K-j)$ from its angular boundary and zero. On that disk $|\Phi'|\le C(K-j)^2$. Lemma~\ref{ied:lead:lem:fatou} identifies the perturbed point with the same inverse branch. Thus the phase error at step $j$ is at most $C(K-j)^2/n^2$, and the total is $\Oh(K^3/n^2)=o(1)$. Together with the initial uniform error this strictly improves the bootstrap for large $n$.

At time $n-L$ the value tends uniformly to $P(s-L)$. In the remaining fixed $L$ steps the subsidiary forcing is still $\Oh(n^{-2})$. Finite compositions of the entire function $f$ preserve convergence on these compact images, proving \eqref{ied:lead:eq:local}.
\end{proof}

\begin{proposition}\label{ied:lead:prop:tail}
Fix $\sigma\in\R$ and put $r=T/(D_n-\sigma)$. There are $M,\delta,C>0$ such that, for all large $n$ and $M\le t\le\delta n$,
\begin{equation}\label{ied:lead:eq:tail}
 |F_n(r\e^{\ii t/n})|\le C/t,\qquad
 \left|\frac{r\e^{\ii t/n}}n F_n'(r\e^{\ii t/n})\right|\le C/t^2.
\end{equation}
The same holds for negative $t$ with $|t|$ in place of $t$.
\end{proposition}
\begin{proof}
Choose $A$ large and $K=\lceil At\rceil$. Let $z_0=r\e^{\ii t/n}$. Since $|z_0-r|\le Ct/n^2$, choosing $A$ after the entrance constant $q$ and then choosing $b>0$ small places
\[
 |z-z_0|\le bt/n^2
 \quad\hbox{inside}\quad |z-r|\le qK/n^2.
\]
Choose $\delta$ small enough that $K/n\le\kappa$ and $K\le n/2$. Expansion of $H$, using $H'(z)=\Oh(n^2)$ on this disk and $H(r)=n-\sigma+\Oh(\log n/n)$, gives
\begin{equation}\label{ied:lead:eq:initialphase}
 s_0:=\Phi(F_{n-K}(z))
 =-K+\ii t+\Oh(1+bt+t^2/n+t\log n/n+K^3/n^2).
\end{equation}
The bounded term includes fixed $\sigma$. First choose $b$ and $\delta$ to control proportional errors, then $M$ to control fixed constants and the local inverse radius, and finally $n$ to control $\log n/n$. The error in \eqref{ied:lead:eq:initialphase} is at most $t/10$.

Use the fixed prototype $p_j=-K+j+\ii t$, $0\le j\le K$, and bootstrap $|s_j-p_j|\le t/4$. Then
\[
 \Ima s_j\ge3t/4,\quad\Rea s_j\le t/4,\quad
 |\arg(-s_j)|\le\pi/2+\arctan(1/3)<2\pi/3,
 \quad ct\le|s_j|\le CK.
\]
After increasing $M$, the inverse estimate gives
\[
 |\arg P_{\rm loc}(s_j)|\le2\pi/3,\qquad
 c/K\le |P_{\rm loc}(s_j)|\le C/t\le\epsilon/4.
\]
Construct $\Phi$ beforehand on the wider sector $|\arg w|<3\pi/4$, $|w|<\epsilon$, with further harmless room for nested-sector estimates. The same bounds hold for $s_j+1$, with the next prototype index. The unforced image $f(P_{\rm loc}(s_j))=P_{\rm loc}(s_j+1)$ has distance at least $c'/K$ from zero and the angular boundary, and at least $\epsilon/2$ from the outer circle.

The actual forcing is $\Oh(n^{-2})$ by Lemma~\ref{ied:lead:lem:subsidiary}. Its full perturbation disk lies inside the wider coordinate domain for large $n$, since $K\le n/2$. On this disk $|\Phi'|\le CK^2$, so one step changes the ideal unit phase increment by at most $CK^2/n^2$. The inverse-identification clause of Lemma~\ref{ied:lead:lem:fatou}, with a fixed slight angular enlargement, ensures that the actual perturbed value is $P_{\rm loc}(s_{j+1})$. It does not switch branches.

The total additional phase error is at most $CK^3/n^2\le C'A^3\delta^2t$. Make this at most $t/20$. Combined with the initial $t/10$ error it is at most $3t/20$, strictly improving the $t/4$ bootstrap. Therefore $|F_n(z)|\le C/t$ throughout the disk about $z_0$. Cauchy's derivative estimate there yields $|F_n'(z_0)|\le Cn^2/t^2$, proving \eqref{ied:lead:eq:tail}. Conjugation proves the negative-angle case.

The noncircular order of constants is: fixed nested sectors and clock cones; small $\epsilon$ and inner analyticity radii; $c_0,q$; large $A$; small $b$; small $\delta$; large $M$; large $n$. All may depend on fixed $\sigma$, but none depends on $n,t$ once chosen.
\end{proof}

\begin{lemma}\label{ied:lead:lem:outer}
For each fixed $\delta>0$, uniformly for $\delta\le|\theta|\le\pi$ on the same moving circle,
\begin{equation}\label{ied:lead:eq:outer}
 F_n(r\e^{\ii\theta})=\Oh(n^{-1}),\qquad
 F_n'(r\e^{\ii\theta})=\Oh(1).
\end{equation}
\end{lemma}
\begin{proof}
Put $z=\zeta/n$ and $v_h=nF_h(z)$. On a bounded $v$-domain, the subsidiary estimate gives
\[
 v_{h+1}=v_h+\frac{v_h^2+\zeta^2}{2n}+\Oh(n^{-2}),
 \qquad v_0=\zeta.
\]
The limiting ordinary differential equation has solution
\[
 v(u)=\zeta\tan(\zeta u/2+\pi/4),\qquad0\le u\le1.
\]
A pole at positive real $u\le1$ requires $\zeta=(\pi/2+2k\pi)/u$. On $|\zeta|=T$ only $\zeta=T$ causes a pole; the negative pole family has modulus at least $3T$. Thus the closed arc away from $T$ has compact complex neighborhoods on which the full ODE trajectory is bounded.

Choose a bounded disk strictly containing those trajectories. Its vector field is uniformly Lipschitz, and the local discretization error is $\Oh(n^{-2})$. Discrete Gronwall gives $\Oh(n^{-1})$ error while the numerical orbit stays in the disk; this estimate prevents a first exit for large $n$. Hence $F_n=\Oh(n^{-1})$ uniformly on the complex neighborhoods. Cauchy in the $\zeta$ variable gives $F_n'=\Oh(1)$. The moving arc is contained in those neighborhoods because $nr\to T$.
\end{proof}

\section{The positive measure and its density}\label{ied:lead:sec:measure}
\begin{proposition}\label{ied:lead:prop:measure}
The entire function $P$ is the Laplace transform of a finite positive measure $\nu$ on $[0,\infty)$ with every exponential moment, no atom at zero, and support $[0,\infty)$. Its restriction to $(0,\infty)$ has a smooth strictly positive density $h$.
\end{proposition}
\begin{proof}
Put $r_n=T/D_n$ and define the finite positive measures
\[
 \nu_n=\sum_{m\ge1}[z^m]F_n(z)\,r_n^m\,\delta_{m/n}.
\]
Their actual Laplace transforms are $F_n(r_n\e^{s/n})$. The rational parameter in Proposition~\ref{ied:lead:prop:local} is
\[
 D_n-\frac{T}{r_n\e^{s/n}}
 =D_n(1-\e^{-s/n})=s+\Oh_S(\log n/n)
\]
on every compact $S$. Hence these transforms converge locally uniformly to $P(s)$.

The masses tend to $P(0)>0$. L\'evy's continuity theorem applied to the normalized characteristic functions gives a finite weak limit $\nu$, supported on $[0,\infty)$. For every real $u$, the transforms at $u+1$ are uniformly bounded. On $t>R$, the integral of $\e^{ut}$ is bounded by $\e^{-R}$ times the integral of $\e^{(u+1)t}$. This proves exponential uniform integrability, so $P(u)=\int\e^{ut}\nu(\dd t)$. Analyticity and the exponential moments extend this identity to every complex $u$.

By Lemma~\ref{ied:lead:lem:fatou}, $P(-x)\sim2/x$. Thus $\nu(\{0\})=0$. Its support has positive points arbitrarily close to zero: a positive gap would imply exponential decay of $P(-x)$, contradicting $2/x$.

For fixed real $\sigma$, Lemma~\ref{ied:lead:lem:fatou} also gives
$P'(\sigma+\ii t)=\Oh(t^{-2})$ as $|t|\to\infty$, since a vertical line eventually lies in an interior inverse sector. This is the Fourier transform of the finite measure $u\e^{\sigma u}\nu(\dd u)$ and is integrable. Fourier inversion gives this weighted measure a continuous density. Dividing by $u\e^{\sigma u}$ shows that $\nu(\dd u)=h(u)\dd u$ on $(0,\infty)$, with $h$ continuous and nonnegative.

The functional equation \eqref{ied:lead:eq:P} and uniqueness of finite measures give
\begin{equation}\label{ied:lead:eq:measureidentity}
 (\e^u-1)\nu=\sum_{k\ge2}\frac{\nu^{*k}}{k!}.
\end{equation}
Both sides are finite: the left side has mass $P(1)-P(0)$, and the positive series on the right has finite total mass. In particular $\nu*\nu\le2(\e^u-1)\nu$, so its closed support $S$ is closed under addition. Positive support points tending to zero have integer multiples approaching any prescribed positive number. Therefore $S=[0,\infty)$.

Fix $t>0$. Full support and continuity yield an interval $I\subset(0,t)$ on which $h>0$. Full support also yields a point $b\in t-I$ with $h(b)>0$. Then $h(t-b)>0$, and continuity implies $(h*h)(t)>0$. The convolution is finite and continuous on $(0,\infty)$: near any fixed $t$, split the defining integral into two small endpoint pieces and a compact interior. On an endpoint piece one factor is bounded away from zero's neighborhood while the other has arbitrarily small $L^1$ mass near zero; the interior follows by dominated convergence. Thus the density inequality from \eqref{ied:lead:eq:measureidentity}, initially almost everywhere, holds everywhere by continuity:
\[
 (\e^t-1)h(t)\ge\tfrac12(h*h)(t)>0.
\]
This proves strict positivity, without a pointwise density expansion at zero.
\end{proof}

\begin{writenote}[The density at zero, 6~October 2026]
Part~III supplies the expansion this proof does without: $h(0)=2$ (Proposition~\ref{ied:log:prop:hzero}), an exact Hankel representation near zero (Theorem~\ref{ied:log:thm:hankel}), a $C^\infty$ germ $t^{t/3}h(t)$ (Theorem~\ref{ied:log:thm:germ}) and the power-log expansion \eqref{ied:log:eq:hfirst}; in particular $h$ is continuous at $0$ but not $C^1$ there.
\end{writenote}

\begin{corollary}\label{ied:lead:cor:smooth}
The density $h$ is $C^\infty$ on $(0,\infty)$.
\end{corollary}
\begin{proof}
Fix real $\sigma$. Cauchy's estimates in the inverse sectors of Lemma~\ref{ied:lead:lem:fatou}, applied on disks of radius proportional to $|s|$, give
\[
 P^{(q)}(\sigma+\ii t)=\Oh_q(|t|^{-q-1})\qquad(|t|\to\infty)
\]
for every integer $q\ge1$. This is the Fourier transform of the finite measure $u^q\e^{\sigma u}\nu(\dd u)$. Given $k\ge0$, choose $q>k$. Then $|t|^kP^{(q)}(\sigma+\ii t)$ is integrable. Differentiation under Fourier inversion gives a $C^k$ density for this weighted measure. On $(0,\infty)$ it agrees with $u^q\e^{\sigma u}h(u)$ by uniqueness and continuity. Divide by the positive smooth weight on every compact subinterval of $(0,\infty)$. Since $k$ was arbitrary, $h$ is smooth there.
\end{proof}

\begin{writenote}[Proved twice, 6~October 2026]
Part~II restates this corollary with the same argument after its interface list (Section~\ref{ied:corr:sec:interface}, around \eqref{ied:corr:eq:Pderivative}).
\end{writenote}

\section{Coefficient extraction and the amplitude}\label{ied:lead:sec:extraction}
\begin{proof}[Proof of Theorem~\ref{ied:lead:thm:main}]
All structural claims about $P$ and $\nu$ were proved above. Fix real $\sigma$ and put $r=T/(D_n-\sigma)$. Differentiating the Cauchy coefficient formula, or integrating it by parts over the full circle, gives exactly
\begin{equation}\label{ied:lead:eq:cauchy}
 n a_n r^n=\frac1{2\pi}\int_{-\pi n}^{\pi n}
 \left[\frac znF_n'(z)\right]\e^{-\ii t}\dd t,
 \qquad z=r\e^{\ii t/n}.
\end{equation}
For the derivative normalization, at $z=T/(D_n-s)$ one has
\[
 \frac znF_n'(z)=\frac{D_n-s}{n}G_n'(s).
\]
Locally uniform holomorphic convergence in Proposition~\ref{ied:lead:prop:local} gives $G_n'\to P'$. On the circle, the actual parameter is
$s_n=D_n-(D_n-\sigma)\e^{-\ii t/n}\to\sigma+\ii t$, uniformly for bounded $t$. The central integrand therefore tends to $P'(\sigma+\ii t)\e^{-\ii t}$.

Proposition~\ref{ied:lead:prop:tail} bounds the intermediate tails by $C\int_M^\infty t^{-2}\dd t=\Oh(M^{-1})$, uniformly in $n$. Lemma~\ref{ied:lead:lem:outer} bounds $zF_n'(z)/n$ by $\Oh(n^{-2})$ on outer arcs, whose $t$-length is $\Oh(n)$, so their contribution is $\Oh(n^{-1})$. First let $n\to\infty$ and then $M\to\infty$ in \eqref{ied:lead:eq:cauchy}. Fourier inversion of the density of $u\e^{\sigma u}\nu(\dd u)$ at $u=1$ gives
\begin{equation}\label{ied:lead:eq:limitcoeff}
 n a_n r^n\longrightarrow
 \frac1{2\pi}\int_\R P'(\sigma+\ii t)\e^{-\ii t}\dd t
 =\e^\sigma h(1)>0.
\end{equation}
This also proves \eqref{ied:lead:eq:constantintegral}.

Finally, expansion of \eqref{ied:lead:eq:constants} yields
\[
 r^{-n}=(n/T)^n n^{-1/3}T^{1/3}\e^{-B-\sigma}(1+o(1)).
\]
The factor $\e^{-\sigma}$ from the contour radius cancels the factor $\e^\sigma$ in \eqref{ied:lead:eq:limitcoeff}. Division by $n$ proves \eqref{ied:lead:eq:asymptotic} with the stated positive constant.
\end{proof}

\begin{writenote}[Refined in Part~II, 6~October 2026]
Part~II's Theorem~\ref{ied:corr:thm:diagonal} adds the first relative correction, with remainder $\Oh(n^{-7/5})$, by extracting coefficients from the thirtieth derivative (Section~\ref{ied:corr:sec:extract}) instead of the first.
\end{writenote}

\section{Proportional height and fixed shifts}\label{ied:lead:sec:proportional}
\begin{theorem}\label{ied:lead:thm:proportional}
Let $0<\beta_0<\beta_1<\infty$ be fixed. Uniformly over positive integer pairs $(n,m)$ with $m\to\infty$ and $\beta=n/m\in[\beta_0,\beta_1]$,
\begin{equation}\label{ied:lead:eq:proportional}
 A(n,m)\sim\frac{h(\beta)}m
 \left(\frac mT\right)^n
 m^{-\beta/3}T^{\beta/3}\e^{-\beta B}.
\end{equation}
The exact integer height $m$ and exact ratio $\beta=n/m$ occur in this formula.
\end{theorem}
\begin{proof}
Fix real $\sigma$, let $r_m=T/(D_m-\sigma)$, and extend
\[
 Q_m(t)=\frac zmF_m'(z),\quad z=r_m\e^{\ii t/m},\quad |t|\le\pi m,
\]
by zero elsewhere. Proposition~\ref{ied:lead:prop:local} gives convergence to $Q(t)=P'(\sigma+\ii t)$ uniformly on each compact interval. Proposition~\ref{ied:lead:prop:tail} gives $|Q_m(t)|\le C/t^2$ in the intermediate range; Lemma~\ref{ied:lead:lem:outer} gives $\Oh(m^{-2})$ in the outer range. The limiting derivative also has an integrable $\Oh(t^{-2})$ tail. Splitting the integral at a fixed $M$, then letting $m\to\infty$ and $M\to\infty$, proves
\[
 \|Q_m-Q\|_{L^1(\R)}\longrightarrow0.
\]
Their Fourier transforms consequently converge uniformly for all real frequencies, by the $L^1$ norm bound. At the exact frequency $\beta=n/m$, the coefficient identity is
\[
 n A(n,m)r_m^n=\frac1{2\pi}\int_\R Q_m(t)\e^{-\ii\beta t}\dd t
 =\beta\e^{\beta\sigma}h(\beta)+o(1).
\]
Fourier inversion applies to the continuous density of $u\e^{\sigma u}\nu(\dd u)$. On the chosen compact interval, $\beta$ is bounded away from zero and $h$ has a strictly positive minimum. The absolute error therefore gives the uniform relative estimate
\[
 A(n,m)=\frac{\e^{\beta\sigma}h(\beta)}m r_m^{-n}(1+o(1)).
\]
Write $q_m=-(\log m)/3+(\log T)/3-B-\sigma$. Since $n=\beta m$ exactly,
\begin{align*}
 r_m^{-n}
 &= (m/T)^n\exp\left(\beta m\log(1+q_m/m)\right)\\
 &= (m/T)^n m^{-\beta/3}T^{\beta/3}
       \e^{-\beta B-\beta\sigma}
       \left(1+\Oh((\log m)^2/m)\right),
\end{align*}
uniformly for these $\beta$. Cancellation of the $\sigma$ factors proves the theorem.
\end{proof}

\begin{corollary}\label{ied:lead:cor:shifts}
For every fixed integer $k$, as $n\to\infty$,
\[
 \frac{[z^n]F_{n+k}(z)}{a_n}\longrightarrow\e^k.
\]
For fixed real $\lambda>0$, put $\alpha_n=\{\lambda n\}$ and $m=\lfloor\lambda n\rfloor$. Then
\begin{equation}\label{ied:lead:eq:floor}
 A(n,m)\sim C_\lambda\e^{-\alpha_n/\lambda}
       (\lambda n/T)^n n^{-1-1/(3\lambda)},
 \qquad
 C_\lambda=h(1/\lambda)\lambda^{-1-1/(3\lambda)}
               T^{1/(3\lambda)}\e^{-B/\lambda}.
\end{equation}
\end{corollary}
\begin{proof}
For $m=n+k$, continuity gives $h(n/m)\to h(1)$, while $(m/n)^n\to\e^k$. All remaining ratios of prefactors in \eqref{ied:lead:eq:proportional} and \eqref{ied:lead:eq:asymptotic} tend to one. For the second assertion, $m=\lambda n-\alpha_n$ gives
\[
 (m/(\lambda n))^n=\e^{-\alpha_n/\lambda}(1+\Oh(n^{-1})),
 \quad n/m=1/\lambda+\Oh(n^{-1}).
\]
Substitution into \eqref{ied:lead:eq:proportional}, using continuity of $h$ and $\Oh(\log n/n)$ for changes in the logarithmic powers, proves \eqref{ied:lead:eq:floor}.
\end{proof}
The oscillating floor factor need not tend to one. Dropping it changes the leading equivalent. Neither theorem asserts uniformity as $\beta\downarrow0$ or $\beta\to\infty$.

\begin{writenote}[Extended in Parts~II and~III, 6~October 2026]
Part~II's Theorem~\ref{ied:corr:thm:height} adds the first correction on the same compact ranges. Part~III's Theorem~\ref{ied:log:thm:transfer} extends the leading equivalent of Theorem~\ref{ied:lead:thm:proportional} (the same formula, since $e^{\beta d_m}=m^{-\beta/3}T^{\beta/3}e^{-\beta B}$) toward $\beta\downarrow0$, uniformly on every band $(\log m)^{-D}\le\beta\le B_0$, with an explicit relative error; this yields the crossover of Theorem~\ref{ied:log:thm:crossover}. Uniformity as $\beta\to\infty$ remains open (Section~\ref{ied:sec:further}, item~\ref{ied:q:endpoint}).
\end{writenote}

\begin{corollary}\label{ied:lead:cor:window}
For each fixed $C>0$, uniformly over integers $|k|\le C\sqrt n$,
\[
 \e^{-k}\frac{[z^n]F_{n+k}(z)}{a_n}
 =\exp\left(-\frac{k^2}{2n}\right)(1+o(1)).
\]
\end{corollary}
\begin{proof}
Put $m=n+k$ in Theorem~\ref{ied:lead:thm:proportional}. Uniformly in this range, $\beta=n/m=1+\Oh(n^{-1/2})$, so $h(\beta)/h(1)\to1$. The other prefactor ratios tend uniformly to one: their logarithms, aside from the continuous density ratio, are $\Oh((1+|k|)\log n/n)$. The exponential ratio is determined by
\[
 n\log(1+k/n)=k-\frac{k^2}{2n}+\Oh_C(n^{-1/2}).
\]
Substitution and division by the diagonal equivalent prove the claim.
\end{proof}
This is a Gaussian-shaped asymptotic profile for height shifts. It is not a probability central limit theorem; no probability distribution on heights has been introduced.


\section{An inverse consequence}\label{ied:lead:sec:inverse}
\begin{corollary}\label{ied:lead:cor:inverse}
Let $y=a_n$ with $n\to\infty$, put $Y=\log y$, and let $W$ denote the positive real Lambert function. Put
\[
 X=\frac{Y}{W(Y/T)},\qquad q=\log(X/T)+1.
\]
Then
\begin{equation}\label{ied:lead:eq:inverse}
 n=X+\frac{\frac43\log X-\log c}{q}
      +o(1/\log X).
\end{equation}
In particular $n-X\to4/3$. The explicit corrected expression in \eqref{ied:lead:eq:inverse}, and the large real root $x(y)$ of
\[
 \log y=x\log(x/T)-\tfrac43\log x+\log c,
\]
recover $n$ by nearest-integer rounding for all sufficiently large exact sequence values $y=a_n$.
\end{corollary}
\begin{proof}
Let $L(x)=x\log(x/T)$. Then $L(X)=Y$ and $L'(X)=q$. Theorem~\ref{ied:lead:thm:main} gives
$L(n)-L(X)=\frac43\log n-\log c+o(1)$. First comparison gives $n/X\to1$, and the mean-value theorem then gives $n-X=\Oh(1)$. Thus
\[
 L(n)-L(X)=q(n-X)+\Oh(X^{-1}),\qquad
 \log n=\log X+\Oh(X^{-1}).
\]
Division by $q$ proves \eqref{ied:lead:eq:inverse}. Applying the same mean-value argument to the full leading model gives $x(a_n)=n+o(1/\log n)$; the rounding assertions follow.
\end{proof}
No effective onset is asserted. This is recovery from exact sequence values, not a universal threshold-rounding rule for arbitrary $y$ between them. The leading equivalent alone does not provide a higher-order error term.

\begin{writenote}[Sharpened in Part~II; the transseries volume, 6~October 2026]
Part~II's Corollary~\ref{ied:corr:cor:inverse} sharpens \eqref{ied:lead:eq:inverse} to an error $\Oh(X^{-7/5}/\log X)$; its $\delta_0=(\frac43\log X-\log c)/q_X$ is the correction here. $X=Y/W(Y/T)$ solves $X\log(X/T)=Y$, an instance of the factorial core \texttt{p0:prop:factorial-core} of \cite{TSvol} ($\kappa=1$, its $d=-\log T$). The correction and the rounding at exact values are proved directly by the mean-value theorem; they parallel \texttt{p0:thm:perturbed-inversion} and part~(3) of \texttt{p0:thm:staircase} but are not instances of them (front matter, ``The OEIS statements, and the inverses'').
\end{writenote}

\section{Exact computation and numerical diagnostics}\label{ied:lead:sec:computation}
\subsection{The coefficient algorithm}
Logarithmically differentiating \eqref{ied:lead:eq:eulerproduct} gives
\[
 \frac{zF_{h+1}'(z)}{1+F_{h+1}(z)}=\sum_{j\ge1}b_jz^j,
 \qquad b_j=\sum_{d\mid j}d[z^d]F_h(z).
\]
Therefore, writing $1+F_{h+1}(z)=\sum_{k\ge0}c_kz^k$, one obtains
\begin{equation}\label{ied:lead:eq:exactrec}
 c_0=1,\qquad c_k=\frac1k\sum_{j=1}^k b_jc_{k-j}\quad(k\ge1).
\end{equation}
Truncation at degree $N$ is exact for every coefficient retained. The public standard-library program checks divisibility before each integer division. It stores $F_h$ with constant term zero and restores the convention $A(0,h)=1$ only when reporting the array. Degree and height are each capped at $640$; this is an implementation and testing limit, not a limitation of the theorems. The diagonal calculation uses $\Oh(N^3)$ integer arithmetic operations and $\Oh(N)$ stored integers; operation costs grow with the bit lengths.

An independent product implementation expands each factor using
\[
 (1-z^d)^{-a}=\sum_{k\ge0}\binom{a+k-1}{k}z^{dk}\quad(a\ge1)
\]
and multiplies the finite truncations. It checks every entry through degree and height $12$. A separate generated exact fixture contains diagonal values through $414$. The public tests also compare the $22$ displayed OEIS values, $n=0,\ldots,21$. The full OEIS b-file was not successfully retrieved in the source review; the $415$-entry fixture is therefore explicitly described as generated mathematical data, not as a $415$-term external comparison.

All public pass/fail checks use explicit exceptions, so Python optimization does not remove them. Local chunked decimal conversion works with Python's minimum configurable integer-string digit limit of $640$ without modifying that process-wide limit. This matters because the largest exact integers exceed that many digits.

\subsection{The amplitude computation is not certified}\label{ied:lead:sub:notcertified}
The normalized ratios $a_n(T/n)^n n^{4/3}$ decrease slowly over the tested range: approximately $4.60133$ at $n=40$, $4.55287$ at $n=80$, $4.52400$ at $n=160$, $4.51777$ at $n=200$, and $4.50437$ at $n=414$. These finite values alone neither prove a limit nor certify its decimal digits.

A separate optional numerical script evaluates the repelling parametrization. The first terms of its normalized formal coordinate are
\[
 \Phi(w)=-2/w+\tfrac13\log w-\frac w{36}
          +\frac{w^2}{540}+\frac{w^3}{7776}+\cdots.
\]
The script uses nine correction coefficients, solves the truncated inverse equation near $s-M$ by Newton iteration, applies $M$ forward iterates of $f$, and propagates the derivative by multiplication by $\e^w$. Subtracting $2/(1-s)^2$ from the derivative before quadrature improves the observed tail: its known contribution to $h(1)$ is $2/\e$ for $\sigma<1$. Tests varying the forward depth, vertical contour, cutoff, and mesh give values near $4.4923299$, consistent with the conjectured $4.4923\ldots$. Floating-point agreement is not an interval enclosure. There is no certified digit count in this report, and the analytic theorem does not depend on these calculations.

\begin{writenote}[Checks, 6~October 2026]
The shipped diagnostic \file{data/225-leading-data-amplitude_diagnostic.json} records $h(1)=1.8146095596784502$ and $c=4.492329896751906$; the write recomputed $T^{1/3}e^{-B}h(1)=4.4923298968$ from that $h(1)$. Part~II's derivative diagnostics give $h(1)\approx1.81460956$ as well (Subsection~\ref{ied:corr:sub:diagnostics}). The coefficients above agree with Part~III's $a_1=-1/36$, $a_2=1/540$, $a_3=1/7776$, and the write's rerun of \file{coordinate_series.py --order 9} also gives $a_4=-71/435456$. All values remain uncertified.
\end{writenote}

\subsection{Reproducibility and limits}
The accompanying archive contains the complete editable TeX, exact programs and tests, numerical diagnostic script, exact fixture, source ledger, and a deterministic builder. The builder verifies an input manifest, runs explicit tests, creates a private pdf\LaTeX\ format, compiles three times without shell escape, rejects layout and reference warnings, and packages sorted members with fixed timestamps. It preserves the source tree and writes outputs only to a new separate directory. Ordinary and optimized cap-$640$ runs are compared by full ZIP bytes and by each member. Byte identity is a claim about the same toolchain; arbitrary TeX, font, or numerical-library versions need not reproduce identical bytes.

\begin{writenote}[The companion here, 6~October 2026]
Report~225's companion is shipped as \file{code/225-leading-*.py}, \file{data/225-leading-*} and \file{225-leading-sources-SOURCE_LEDGER.md}; its delivered README was replaced by the report README, and its PDF and \file{MANIFEST.sha256} are not shipped. The builder and tests run only in the delivered layout; the report README gives the rerun route from the arrival commit.
\end{writenote}

\section{Further questions and quantitative extensions}\label{ied:lead:sec:further}
\subsection{Validated evaluation of the positive constant}\label{ied:lead:sub:validated}
The exact expression \eqref{ied:lead:eq:constantintegral} is a useful starting point for a certificate, but a validated computation requires more than increased precision. One must bound the omitted Fatou-coordinate terms on a stated sector, enclose the inverse branch, control the forward propagation of both value and derivative, validate the finite-interval quadrature, and give an explicit integrable bound for the remaining vertical tail. The sector proof above supplies qualitative constants; it does not extract numerical values for all of them. A practical goal is a reproducible interval for $c$ narrow enough to validate the conjectured initial digits, with each of these error sources accounted for separately.

The rational subtraction used in the diagnostic script could be extended to several terms derived from the large-$s$ inverse expansion. This can reduce cancellation and the cutoff length, but every subtracted term and its Fourier contribution must be explicit. Mesh or precision stability cannot replace bounds for truncation of the asymptotic coordinate and the infinite integral.

\begin{writenote}[Still open, 6~October 2026]
Not settled by Parts~II or~III; Section~\ref{ied:sec:further}, item~\ref{ied:q:validated}. Part~II's closed form of $\kappa$ removes one constant from the first correction, which still needs $h(1)$, $h'(1)$, $h''(1)$.
\end{writenote}

\subsection{Higher corrections to the diagonal}\label{ied:lead:sub:higher}
The logarithmic center $D_n$ is determined only to the precision needed for a leading equivalent. Higher corrections require retaining more subsidiary terms in $R_h(z)$, including their dependence on $h$, and constructing a corresponding hierarchy of nonautonomous clocks. Their matching constants interact with the inverse-coordinate expansion. It is natural to investigate corrections involving powers of $1/n$ and polynomials in $\log n$, but no complete expansion of this form is asserted here.

A formal expansion at bounded $s$ is insufficient for a coefficient theorem. The error must be controlled in an integrable norm on the growing contour, with the intermediate-angle and outer-angle arguments strengthened to the same order. The present choice $K=\lfloor\sqrt n\rfloor$ balances errors only for convergence; a higher-order theorem would need a new quantified matching analysis. This identifies a concrete obstruction to promoting formal series directly into an all-orders claim.

\begin{writenote}[Partly answered by Part~II, 6~October 2026]
Part~II proves the first correction, $(-\frac1{18}\log^2n+A\log n+C)/n$ with remainder $\Oh(n^{-7/5})$ (Theorem~\ref{ied:corr:thm:diagonal}), by two further clock terms and a quantified growing-strip matching with $K=\lfloor n^{1/2}\rfloor$; it confirms the expected shape, a polynomial in $\log n$ of degree two over $n$. Corrections beyond the first and an all-orders theorem remain open (Section~\ref{ied:sec:further}, item~\ref{ied:q:higher}).
\end{writenote}

\subsection{Effective bounds and inverse thresholds}\label{ied:lead:sub:effective}
Making the constants in the clock, stopping, and sector estimates explicit would give finite-$n$ bounds and a computable onset for the leading equivalent. Such bounds would in turn permit certified recovery of integer thresholds for arbitrary input sizes. Corollary~\ref{ied:lead:cor:inverse} currently supplies eventual rounding at exact sequence values; it does not say how large ``eventually'' is, nor does it control the location of a general input between adjacent terms. A rigorous threshold algorithm should combine exact small-index computation with explicit upper and lower bounds, rather than treat a floating-point evaluation of $c$ as exact.

One can also seek inverse statements for $A(n,\lfloor\lambda n\rfloor)$. Formula~\eqref{ied:lead:eq:floor} shows why the bounded floor phase must enter such statements. A single smooth approximation that discards this phase can have a nonvanishing multiplicative error.

\begin{writenote}[Still open, 6~October 2026]
Part~II's Corollary~\ref{ied:corr:cor:inverse} sharpens the inverse at exact sequence values, still without onset; Part~III's inverses concern ratios in the window $m\asymp n\log n$, not $A(n,\lfloor\lambda n\rfloor)$. Section~\ref{ied:sec:further}, items~\ref{ied:q:effective} and~\ref{ied:q:floorinverse}.
\end{writenote}

\subsection{Other height scales and the density}\label{ied:lead:sub:heights}
Theorem~\ref{ied:lead:thm:proportional} treats $n/m$ in a fixed compact subset of $(0,\infty)$. Extending it toward zero or infinity requires information about $h(t)$ at the corresponding endpoint and renewed uniformity estimates. In particular, the strict-positivity proof does not give an expansion of $h(t)$ at zero. The measure identity
\[
 (\e^t-1)\nu=\sum_{k\ge2}\nu^{*k}/k!
\]
is a possible starting point for endpoint analysis and stronger regularity, such as real analyticity. Smoothness away zero is already established in Corollary~\ref{ied:lead:cor:smooth}, but differentiation of this convolution identity or exchange of limits with its series needs separate justification.

Kaneiwa's fixed-leaf polynomial regime, proportional height, and fixed-height partition regimes are genuinely different limiting questions. A uniform transition theory between them would explain how the polynomial and parabolic descriptions fit together. It may also clarify which parts of this proof extend to other multiset recurrences or other initial generating functions. Such a generalization must recompute the subsidiary forcing and the clock constant; the unforced map alone does not determine the answer.

\begin{writenote}[Partly answered by Part~III, 6~October 2026]
At the endpoint zero Part~III proves the expansion asked for here: $h(0)=2$, the Hankel representation and the power-log series (Proposition~\ref{ied:log:prop:hzero}, Theorem~\ref{ied:log:thm:hankel}, Corollary~\ref{ied:log:cor:densityseries}), and it extends the height theorem toward $n/m\to0$ on logarithmic bands (Theorem~\ref{ied:log:thm:transfer}), with Kaneiwa's normalization compared in Section~\ref{ied:log:sec:kaneiwa}. Theorem~\ref{ied:log:thm:hankel} also gives real analyticity of $h$ on some interval $(0,\epsilon)$. Open: $n/m\to\infty$, real analyticity on all of $(0,\infty)$, the transition to the fixed-degree regime, other recurrences (Section~\ref{ied:sec:further}, items~\ref{ied:q:endpoint}--\ref{ied:q:other}).
\end{writenote}

\clearpage
\part{First logarithmic correction for iterated Euler diagonals}\label{ied:corr:part}
\parteqs{II.}
\partsource{Bundle Report~226, archive \file{Report226.zip}, delivered as \file{src/report226.tex} (17 pages). Delivered title block ``First logarithmic correction for iterated Euler diagonals / Report 226 / 5 October 2026''. Printed as delivered, with \textbf{[write]} notes. Delivered Section~$k$ is Section~$k+11$ here, its statement $k.j$ is $(k+11).j$, and its equation $(k)$ is $(\mathrm{II}.k)$.}
\begin{center}
\small
\begin{tabular}{@{}>{\raggedright\arraybackslash}p{0.46\textwidth} >{\raggedright\arraybackslash}p{0.46\textwidth}@{}}
\toprule
Reading conventions of Part~II & Elsewhere in the report\\
\midrule
``Report 225'', \cite{R225} & Part~I\\
$b_{n,m}$ the array; $d_\star$, $d$, $d_m$ & Parts~I, III: $A(n,m)$; Part~I: $D_n$, $d=n-h$\\
$\rho_1,\rho_2=h'(1)/h(1),h''(1)/h(1)$; $\rho_j(\beta)$ & Part~III: $\rho_j$ rational\\
$A$, $C$ the correction constants; $G$ Catalan's constant & Parts~I, III: $A(n,m)$; Part~III: $\mathcal C$ Catalan, $C_0$, $C(\beta)$\\
$M(y)$, $N(y)$ clock functions; $M=n^{1/20}$; $\mathcal C_z(w)$ & Part~III: $M=m^{1/20}$, $N$ an order\\
$q=30$ the derivative order; $\alpha=4/3$, $\ell=\log c$ & Part~I: $q$ entrance constant; Part~III: $\ell_n$ Kaneiwa\\
\bottomrule
\end{tabular}
\end{center}
\begin{partabstract}
For the iterated Euler diagonal $a_n$ in OEIS A290354, we prove
\[
 a_n=c(2n/\pi)^n n^{-4/3}
 \left\{1+\frac{-\log^2 n/18+A\log n+C}{n}
                  +\Oh(n^{-7/5})\right\}.
\]
The constants $A$ and $C$ are exact expressions in derivatives of the positive density associated with the normalized repelling parametrization of $w\mapsto\e^w-1$. A new nonautonomous clock constant is evaluated in terms of $\pi$, $\log2$, and Catalan's constant. The proof constructs two further clock terms, obtains a weighted fourth-power defect, and matches a forced Fatou correction on a growing vertical strip. Thirtieth-derivative coefficient extraction turns these estimates into the stated relative remainder. A uniform correction for compact positive degree-to-height ratios and a sharpened Lambert $W$ inverse are included. This companion uses the analytic foundation of the preceding unrefereed Report 225, whose unchanged PDF and TeX accompany the archive. It is an unrefereed research proof, not a proof-assistant formalization. Floating-point diagnostics are not certified decimal enclosures; no all-orders theorem is claimed.
\end{partabstract}
\section{Definition and main results}\label{ied:corr:sec:statement}
Define
\begin{equation}\label{ied:corr:eq:rec}
 F_0(z)=z,\qquad
 F_{h+1}(z)=\exp\left(\sum_{j\ge1}\frac{F_h(z^j)}j\right)-1.
\end{equation}
For $n\ge1$ and $m\ge0$, put $b_{n,m}=[z^n]F_m(z)$ and $a_n=b_{n,n}$; set $a_0=1$ and $b_{0,m}=1$ separately when reporting the integer arrays. Thus $a_n$ is A290354 and $b_{n,m}$ is A290353~\cite{A290354,A290353}. The analytic generating functions themselves have constant term zero. They count unlabeled rooted trees whose leaves all have the stated height: a new root carries a nonempty unordered multiset of trees at the previous height. Equivalently,
\[
 1+F_{h+1}(z)=\prod_{d\ge1}(1-z^d)^{-[z^d]F_h(z)}.
\]
There is no factorial normalization. The diagonal begins
\[
 1,1,2,6,30,170,1337,12166,133476,1676364.
\]

The exact constants and functions used throughout are
\begin{align}\label{ied:corr:eq:constants}
 T&=\frac\pi2,&
 B&=\frac12-\frac\pi6-\frac{\pi^2}{16}-\frac{\log2}{6},&
 d_\star&=\frac13\log T-B,\\
 \kappa&=\frac{\pi^3}{96}+\frac{\pi^2}{12}-\frac{13\pi}{48}
        -\frac{11}{36}-\frac{G}{6}-\frac{\pi\log2}{24},&&&\label{ied:corr:eq:kappa}
\end{align}
where $G=\sum_{j\ge0}(-1)^j/(2j+1)^2$ is Catalan's constant. Let $P$ denote the entire repelling parametrization normalized by
\begin{equation}\label{ied:corr:eq:Pphi}
 P(s+1)=\e^{P(s)}-1,\qquad
 \Phi(w)=-\frac2w+\frac13\log w+\Oh(w),\qquad \Phi(P(s))=s
\end{equation}
in the inverse sector near negative real $s$. Report 225~\cite{R225} constructs this normalization and proves
\begin{equation}\label{ied:corr:eq:density}
 P(s)=\int_0^\infty \e^{sx}h(x)\dd x,\qquad
 h\in C^\infty(0,\infty),\qquad h(x)>0\quad(x>0).
\end{equation}
The density has every exponential moment. Set
\begin{equation}\label{ied:corr:eq:rhos}
 c=T^{1/3}\e^{-B}h(1),\qquad
 \rho_1=\frac{h'(1)}{h(1)},\qquad \rho_2=\frac{h''(1)}{h(1)}.
\end{equation}
The zero additive constant in \eqref{ied:corr:eq:Pphi} is essential: a translation of the Fatou coordinate would change $P$ and $h$.

\begin{theorem}\label{ied:corr:thm:diagonal}
As $n\to\infty$,
\begin{equation}\label{ied:corr:eq:main}
 a_n=c(n/T)^n n^{-4/3}
 \left(1+\frac{-\log^2n/18+A\log n+C}{n}
                  +\Oh(n^{-7/5})\right),
\end{equation}
where $c>0$ and
\begin{align}\label{ied:corr:eq:AC}
 A&=\frac{d_\star}3+\frac49+\frac{\rho_1}3,\\
 C&=-\frac{d_\star^2}2-d_\star\left(\frac43+\rho_1\right)
       -T\kappa-\frac13-\frac{4\rho_1}3-\frac{\rho_2}2.\notag
\end{align}
The error is relative to the displayed leading expression. Its exponent is a convenient nonoptimal choice.
\end{theorem}
With $d=-\frac13\log n+d_\star$, the numerator of the correction can equivalently be written
\begin{equation}\label{ied:corr:eq:Qd}
 -\frac{d^2}2-d\left(\frac43+\rho_1\right)-T\kappa
             -\frac13-\frac{4\rho_1}3-\frac{\rho_2}2.
\end{equation}
In particular the squared-logarithm coefficient is exactly $-1/18$. Retaining only a more accurate contour radius would miss some terms in $A$ and $C$.

\begin{theorem}\label{ied:corr:thm:height}
Let $I=[\beta_0,\beta_1]\subset(0,\infty)$ be fixed. Uniformly for integers $m\to\infty$ and $n\ge1$ with $\beta=n/m\in I$, put $d_m=-\frac13\log m+d_\star$ and
$\rho_j(\beta)=h^{(j)}(\beta)/h(\beta)$ for $j=1,2$. Then
\begin{align}\label{ied:corr:eq:height}
 b_{n,m}={}&h(\beta)T^{\beta/3}\e^{-\beta B}
 (m/T)^n m^{-1-\beta/3}
 \left(1+\frac{\mathcal Q_\beta(d_m)}m+\Oh_I(m^{-7/5})\right),\\
 \mathcal Q_\beta(d)={}&-\frac{\beta d^2}2
 -d\left(1+\frac\beta3+\beta\rho_1(\beta)\right)-\beta\kappa T\notag\\
 &-\frac13-\left(1+\frac\beta3\right)\rho_1(\beta)
 -\frac\beta2\rho_2(\beta).\notag
\end{align}
The exact height and exact ratio occur in the formula; no additional factor of $\beta$ belongs in its leading amplitude.
\end{theorem}
No uniform conclusion is asserted as $\beta$ approaches zero or infinity. In particular replacing an integer height by an unfloored real height without retaining its phase is not justified.

\begin{writenote}[Toward $\beta\to0$ in Part~III, 6~October 2026]
Part~III extends the \emph{leading} term of \eqref{ied:corr:eq:height} to the bands $(\log m)^{-D}\le\beta\le B_0$ (Theorem~\ref{ied:log:thm:transfer}); no Part extends the correction $\mathcal Q_\beta$ there, and $\beta\to\infty$ is open (Section~\ref{ied:sec:further}, item~\ref{ied:q:endpoint}).
\end{writenote}

\subsection{Dependence on the preceding report and prior work}\label{ied:corr:sub:dependence}
Report 225 proves the leading equivalent and its analytic foundation. The present report extends its clock and coefficient analysis; it does not independently reprove all entrance and angular-tail estimates. Section~\ref{ied:corr:sec:interface} records the full interface used, and the accompanying archive includes the unchanged complete proof of Report 225. Neither report has undergone journal refereeing.

The OEIS attributes the leading equivalent with numerical amplitude $4.4923\ldots$ to V\'aclav Kot\v{e}\v{s}ovec, dated 14 August 2017~\cite{A290354}. The fixed-degree polynomiality and tangent leading profile have earlier antecedents explained and sourced in Report 225. The correction proved here is not inferred merely from that fixed-degree information. We make no worldwide priority claim for the result or for parabolic-coordinate methods.

\begin{writenote}[Reading Part~II against Part~I, 6~October 2026]
``Report 225'' and \cite{R225} are Part~\ref{ied:lead:part} of this report; its Sections~2--7 are Part~I's Sections~\ref{ied:lead:sec:subsidiary}--\ref{ied:lead:sec:extraction}, with the same numbers. The included copy is byte-identical to Part~I and is not shipped. At leading order Theorem~\ref{ied:corr:thm:diagonal} is Part~I's Theorem~\ref{ied:lead:thm:main} and Theorem~\ref{ied:corr:thm:height} is Part~I's Theorem~\ref{ied:lead:thm:proportional}. Here $b_{n,m}$ is Part~I's $A(n,m)$, and $d_\star=\frac13\log T-B$, so $D_n$ is Part~I's $D_n$. Part~II's $\rho_j$ are not Part~III's rational $\rho_j$.
\end{writenote}

\section{Analytic foundation from Report 225}\label{ied:corr:sec:interface}
All branches agree with real values on the positive axis. Constants in local analytic estimates may depend on fixed sectors or fixed real-width bounds, but not on the growing indices. The following are the precise results imported from Report 225, Sections 2--7~\cite{R225}.

\begin{enumerate}[label=\textup{(\roman*)}]
\item Each $F_h$ is analytic in $|z|<1$, with nonnegative coefficients and linear coefficient one. If $(h+1)|q|$ is sufficiently small, then $|F_h(q)|\le C|q|$ uniformly. This follows by positive-coefficient comparison on the real radius and a reciprocal recurrence bound; it does not assume that the main value $F_h(z)$ is small.
\item Define
\begin{align}\label{ied:corr:eq:JLH}
 J(y)&=2\arctan y-T,\qquad L(1)=0,\qquad
 L'(y)=\frac{y^3/3-y-4/3-2J(y)}{(1+y^2)^2},\notag\\
 H(z)&=\frac Tz+\frac13\log\frac1z+B.
\end{align}
The old clock $J(w/z)/z+L(w/z)$ gives a noncircular positive-real entrance bound and its Cauchy extension. More precisely, let $D_n=n-\frac13\log n+d_\star$ and $r=T/D_n$. For suitable fixed $q>0$, $K_0$, and $\epsilon>0$, if $K_0\le K\le\epsilon n$, then on $|z-r|\le qK/n^2$ the orbit through $n-K$ is small and
\begin{equation}\label{ied:corr:eq:entrance}
 |F_h(z)|\asymp(n-h)^{-1}\quad(0\le h\le n-K),
 \qquad F_h(z)/z\ \hbox{lies in a fixed narrow positive cone}.
\end{equation}
The cone has a slightly larger fixed cone available for Taylor segments. At entrance the old phase estimate is
\begin{equation}\label{ied:corr:eq:oldphase}
 \Phi(F_{n-K}(z))=n-K-H(z)
       +\Oh(K^{-1}+n^{-1}+K^3/n^2).
\end{equation}
\item The sector coordinate $\Phi$ in \eqref{ied:corr:eq:Pphi} satisfies the Abel equation $\Phi(f(w))=\Phi(w)+1$, $f(w)=\e^w-1$. Its inverse is uniquely selected in fixed interior sectors about the closed left half-plane for $s$ of large modulus, and
\begin{equation}\label{ied:corr:eq:Pinv}
 P(s)=-\frac2s+\Oh\left(\frac{\log|s|}{|s|^2}\right).
\end{equation}
A small point $w$ in a smaller coordinate sector satisfying $\Phi(w)=s$ is the same inverse value $P(s)$; global injectivity is unnecessary. The inverse iterates $g^{\circ j}(w)$, $g(w)=\log(1+w)$, obey
$|g^{\circ j}(w)|\le C/(|w|^{-1}+j)$ in nested sectors. Extension by finitely many iterates of $f$ gives the entire $P$.
\item There are fixed $b,\delta,M_0,C>0$ such that for $M_0\le |t|\le\delta n$,
\begin{equation}\label{ied:corr:eq:oldtaildisk}
 |z-r\e^{\ii t/n}|\le b|t|/n^2
       \quad\Longrightarrow\quad |F_n(z)|\le C/|t|.
\end{equation}
For the outer angles $\delta\le|\theta|\le\pi$, the compact arc $\zeta\approx T\e^{\ii\theta}$ has fixed complex neighborhoods on which $F_n(\zeta/n)=\Oh(n^{-1})$. These neighborhoods follow from the bounded ODE profile
$\zeta\tan(\zeta u/2+\pi/4)$ for $0\le u\le1$ away from its positive pole. Both are bounds on two-dimensional analytic neighborhoods, not merely on one-dimensional arcs.
\item The positive-measure construction yields \eqref{ied:corr:eq:density}, with no atom at zero, support $[0,\infty)$, and strict positivity of $h$ on $(0,\infty)$. Its leading coefficient consequence is $a_n\sim c(n/T)^n n^{-4/3}$.
\end{enumerate}

\begin{writenote}[Where Part~I proves each item, 6~October 2026]
(i)~Lemma~\ref{ied:lead:lem:subsidiary} and the first step of its proof. (ii)~\eqref{ied:lead:eq:JL}--\eqref{ied:lead:eq:clock}, Lemmas~\ref{ied:lead:lem:clock}--\ref{ied:lead:lem:entrance} and \eqref{ied:lead:eq:entryphase} (Part~I writes $K\le\kappa n$ for the $K\le\epsilon n$ here; this $\kappa$ is not the clock constant). (iii)~Lemma~\ref{ied:lead:lem:fatou}, \eqref{ied:lead:eq:fatou}--\eqref{ied:lead:eq:Pinverse}; the bound on inverse iterates is the estimate $|G^j(u)|\ge c(|u|+j)$ of its proof, read in $w=2/u$. (iv)~Proposition~\ref{ied:lead:prop:tail}, whose proof (not its statement) bounds $F_n$ on the whole disk about $z_0$, and Lemma~\ref{ied:lead:lem:outer}, whose proof (not its statement, which concerns the moving circle) gives the bound on complex neighbourhoods of the arc. \textup{(Extended after the independent check of 6~October 2026; the first wording read ``Proposition~\ref{ied:lead:prop:tail}, whose proof bounds $F_n$ on the whole disk about $z_0$, and Lemma~\ref{ied:lead:lem:outer}''.)} (v)~Proposition~\ref{ied:lead:prop:measure} and Theorem~\ref{ied:lead:thm:main}. The paragraph below repeats Part~I's Corollary~\ref{ied:lead:cor:smooth} with its proof. The section is printed in full because its displays are cited below.
\end{writenote}

For clarity, the derivative regularity needed below also follows directly from this interface. Sector Cauchy estimates in \eqref{ied:corr:eq:Pinv} give, for each fixed $q\ge1$,
\begin{equation}\label{ied:corr:eq:Pderivative}
 P^{(q)}(\sigma+\ii t)=\Oh_q(|t|^{-q-1})
\end{equation}
uniformly for $\sigma$ in a bounded interval. This is the Fourier transform of $x^q\e^{\sigma x}h(x)$. If $q>k$, its product with $t^k$ is integrable, so Fourier differentiation gives a $C^k$ weighted density. Dividing by the nonzero smooth weight on compact positive intervals proves $h\in C^\infty(0,\infty)$. No boundary behavior of $h$ or its derivatives at zero is used.

\section{Subsidiary forcing through fifth scaled order}\label{ied:corr:sec:forcing}
Write
\[
 R_h(z)=\sum_{j\ge2}F_h(z^j)/j,\qquad w=F_h(z),\qquad y=w/z,\qquad t=hz.
\]
The map is $w_{h+1}=f(w_h+R_h(z))$.
\begin{lemma}\label{ied:corr:lem:forcing}
Uniformly for $h\le C_0/|z|$, as $z\to0$,
\begin{align}\label{ied:corr:eq:forcing}
 R_h(z)={}&\frac{z^2}2+z^3\left(\frac13+\frac t2\right)
 +z^4\left(\frac14+\frac{t^2}4\right)\\
 &+z^5\left(\frac15+\frac{7t}{12}+\frac{t^3}6\right)+\Oh(|z|^6).\notag
\end{align}
\end{lemma}
\begin{proof}
Coefficient extraction in one Euler step gives
\[
 [q^2]F_h=h,\qquad [q^3]F_h=\frac{h(h+1)}2,\qquad
 [q^4]F_h=\frac{h(h+1)(2h+1)}6.
\]
Indeed, if the current coefficients are $a,b,c$, their increments are respectively $1$, $a+1$, and $b+a^2/2+3a/2+1$; the latter becomes $(h+1)^2$. All three start at zero.

Take $\varrho=c_0/(h+1)$ with $c_0$ small. The small-subsidiary bound gives $\sup_{|q|\le\varrho}|F_h(q)|\le C\varrho$. Cauchy and a geometric sum therefore give, for $|q|\le\varrho/2$,
\begin{equation}\label{ied:corr:eq:subsidiarytail}
 F_h(q)=q+hq^2+\frac{h(h+1)}2q^3+
 \frac{h(h+1)(2h+1)}6q^4+\Oh((h+1)^4|q|^5).
\end{equation}
For small $z$, every $q=z^j$, $j\ge2$, lies in this disk. The $j=2$ remainder is $\Oh((h+1)^4|z|^{10})=\Oh(|z|^6)$, and $j\ge6$ contributes $\Oh(|z|^6)$. At order $z^5$, the baseline is $1/5$; the linear-$h$ part of $[q^3]F_h$ at $j=2$ contributes $t/4$, $[q^2]F_h$ at $j=3$ contributes $t/3$, and the cubic-$h$ part of $[q^4]F_h$ at $j=2$ contributes $t^3/6$. This proves \eqref{ied:corr:eq:forcing}, including the otherwise easily lost $7t/12$ term.
\end{proof}

\begin{writenote}[Extends Part~I, 6~October 2026]
This extends Part~I's \eqref{ied:lead:eq:R}, whose $hz^4/2$ is the term $z^3t/2$ here. The three coefficients $[q^2]F_h$, $[q^3]F_h$, $[q^4]F_h$ are the rows $n=2,3,4$ of A290353 (its cross-references name A001477, A000217, A000330); the write checked them against the exact recurrence for $h\le12$.
\end{writenote}

For the clock algebra write $d_y=1+y^2$ and $a(y)=d_y/2$. Expansion of the exponential gives
\begin{equation}\label{ied:corr:eq:ystep}
 \Delta y=za+z^2b(y,t)+z^3v(y,t)+z^4u(y,t)
           +\Oh\bigl(|z|^5(1+|y|)^6\bigr),
\end{equation}
where
\begin{align}\label{ied:corr:eq:stepcoeffs}
 b(y,t)&=y^3/6+y/2+1/3+t/2,\notag\\
 v(y,t)&=y^4/24+y^2/4+y/3+yt/2+3/8+t^2/4,\\
 u(y,t)&=y^5/120+y^3/12+y^2/6+y^2t/4+3y/8+yt^2/4\notag\\
       &\hspace{1.5em}+11/30+5t/6+t^3/6.\notag
\end{align}
The estimate is uniform for bounded $t$ when $|w|+|z|$ is small. These are finite Taylor identities with analytic remainders, not extrapolations from numerical values.

\section{Two further nonautonomous clock terms}\label{ied:corr:sec:clocks}
Define $M(1)=N(1)=0$ on the same simply connected cone as $J,L$. First set
\begin{equation}\label{ied:corr:eq:Q2}
 Q_2(y,t)=J'v+J''ab+\frac{J'''}6a^3+L'b+\frac{L''}2a^2.
\end{equation}
The required equation for $M$ is
\begin{equation}\label{ied:corr:eq:Mrecursive}
 aM'=-Q_2(y,J)-\frac{L}{d_y}.
\end{equation}
Its expanded form is
\begin{align}\label{ied:corr:eq:M}
 M'=-\frac1{36d_y^3}\bigl\{&36J^2(y^2-1)+12J(5y^3-3y-8)+72Ld_y\\
 &+y^6+9y^4+40y^3-21y^2-24y-29\bigr\}.\notag
\end{align}
The feedback term $L/d_y$ is necessary. Setting $hz=J$ before retaining the self-consistency terms would give a wrong next clock.

Next define
\begin{align}\label{ied:corr:eq:Q3}
 Q_3(y,t)={}&J'u+J''(av+b^2/2)+\frac{J'''}2a^2b+\frac{J''''}{24}a^4\notag\\
 &+L'v+L''ab+\frac{L'''}6a^3+M'b+\frac{M''}2a^2,\\
 G_3(y)={}&Q_3(y,J)+L\,\partial_t Q_2(y,J)+M/d_y,
 \qquad N'=-G_3/a.\label{ied:corr:eq:N}
\end{align}
Here $\partial_tQ_2(y,J)=J/d_y+L'/2$. Equations \eqref{ied:corr:eq:Mrecursive} and \eqref{ied:corr:eq:N}, with the zero values at one, are exact recursive definitions involving only differentiation, integration, and rational operations. They specify the functions without a fitted constant.

\subsection{Large ratio expansions}
Uniformly on each closed smaller cone,
\begin{align}\label{ied:corr:eq:JLexp}
 J&=T-2/y+2/(3y^3)+\Oh(y^{-5}),\notag\\
 L&=\tfrac13\log y+B+\frac5{6y^2}
      +\frac{4/9+2T/3}{y^3}-\frac7{4y^4}+\Oh(y^{-5}).
\end{align}
Substitution in \eqref{ied:corr:eq:M} gives
\begin{align}\label{ied:corr:eq:Mprimeexp}
 M'={}&-\frac1{36}-\frac1{6y^2}-\frac{5T/3+10/9}{y^3}\\
 &+\frac{-(2/3)\log y-2B-T^2+9/2}{y^4}+\Oh(y^{-5}),\notag
\end{align}
and integration yields
\begin{align}\label{ied:corr:eq:Mexp}
 M={}&-\frac y{36}+\kappa+\frac1{6y}+
       \frac{5T/6+5/9}{y^2}\\
 &+\frac{(2/9)\log y+2B/3+T^2/3-77/54}{y^3}+\Oh(y^{-4}).\notag
\end{align}
Thus $M$ has neither a divergent $\log y$ nor a divergent $\log^2y$ term. The constant is the convergent integral
\begin{equation}\label{ied:corr:eq:kappaint}
 \kappa=\frac1{36}+\int_1^\infty\left(M'(y)+\frac1{36}\right)\dd y.
\end{equation}

For the next clock, direct expansion in \eqref{ied:corr:eq:N} gives
\begin{align}\label{ied:corr:eq:G3exp}
 G_3={}&-\frac{y^3}{540}-\frac y{360}+\frac T{12}+\frac1{18}\\
 &+\frac{(5/18)\log y+5B/6+5T^2/12-3/20}{y}
       +\Oh(\log y/y^2),\notag\\
 N'={}&\frac y{270}+\frac1{540y}-\frac{T/6+1/9}{y^2}\notag\\
 &+\frac{-(5/9)\log y-5B/3-5T^2/6+161/540}{y^3}
       +\Oh(\log y/y^4),\notag\\
 N={}&\frac{y^2}{540}+\frac{\log y}{540}+n_0+\frac{T/6+1/9}{y}\notag\\
 &+\frac{(5/18)\log y+5B/6+5T^2/12-11/1080}{y^2}
       +\Oh(\log y/y^3).\notag
\end{align}
The finite $n_0$ is fixed by $N(1)=0$ and is not needed at the present order. Unlike $M$, $N$ contains a divergent logarithm. Its clock contribution is $z^2\log y$, below the first correction but relevant to future orders.

\subsection{Exact evaluation of the new clock constant}\label{ied:corr:sub:kappa}
Put $u=\arctan y$. An exact primitive, checked by differentiation and its value at $u=T/2$, is
\begin{align}\label{ied:corr:eq:Lprimitive}
 L(\tan u)={}&-\tfrac13\log\cos u+\tfrac13-\tfrac16\cos2u
 +(T-2/3)u-u^2\notag\\
 &+(T/2-1/3-u)\sin2u-\tfrac16\log2-\tfrac14T^2+T/3.
\end{align}
Substituting it and $J=2u-T$ in \eqref{ied:corr:eq:M} gives
\[
 (M'+1/36)\dd y=\left(\tfrac23\cos^2u\log\cos u+\mathcal P(u)\right)\dd u,
\]
where the finite trigonometric polynomial is
\[
 \mathcal P=p_0+p_2\cos2u+p_4\cos4u+q_2\sin2u+q_4\sin4u
\]
and
\begin{align*}
 p_0&=2u^2-2Tu+T^2/2+8u/3-4T/3+1/12+\log2/6,\\
 p_2&=3u^2-3Tu+3T^2/4+10u/3-5T/3+1/3+\log2/6,\\
 p_4&=u^2-Tu+T^2/4+2u/3-T/3+1/12,\\
 q_2&=2u/3-T/3+2/9,\qquad q_4=7u/6-7T/12+7/18.
\end{align*}
Elementary termwise integration yields
\begin{align*}
 \int_{T/2}^T\mathcal P(u)\dd u={}&-\frac{5\pi}{16}-\frac5{12}-\frac{\log2}{12}
 +\frac{\pi\log2}{24}+\frac{\pi^3}{96}+\frac{\pi^2}{12}.
\end{align*}
The two logarithmic integrals are
\begin{align*}
 \int_{\pi/4}^{\pi/2}\log\cos u\dd u&=-\frac\pi4\log2-\frac G2,\\
 \int_{\pi/4}^{\pi/2}\cos2u\log\cos u\dd u&=\frac{\log2}{4}+\frac\pi8+\frac14.
\end{align*}
The first follows by combining the half-period log-cosine integral with
$G=-\int_0^{\pi/4}\log\tan u\dd u$; the latter identity also follows from the Fourier series of the logarithm. For the second, integration by parts leaves the boundary contribution $\log2/4$ and $\int_{\pi/4}^{\pi/2}\sin^2u\dd u=\pi/8+1/4$. The other endpoint term tends to zero. Using $\cos^2u=(1+\cos2u)/2$ in \eqref{ied:corr:eq:kappaint}, including its $1/36$ normalization, proves exactly \eqref{ied:corr:eq:kappa}.

\begin{writenote}[Checks, 6~October 2026]
At intake the closed form was compared with a direct quadrature of \eqref{ied:corr:eq:kappaint}, using the primitive \eqref{ied:corr:eq:Lprimitive} (itself checked against quadrature of \eqref{ied:lead:eq:JL}): $-0.2543481698511$ against $-0.2543481698560$, the difference coming from the truncated tail. The delivered symbolic suite (\file{code/226-firstcorr-code-clock_checks.py}) reproduced \file{data/226-firstcorr-data-symbolic_checks.json} on a copy at the write.
\end{writenote}

\section{Weighted defect and accumulated clock error}\label{ied:corr:sec:defect}
Define
\begin{equation}\label{ied:corr:eq:C3}
 \mathcal C_z(w)=J(w/z)/z+L(w/z)+zM(w/z)+z^2N(w/z),
 \qquad E_h=\mathcal C_z(F_h(z))-h.
\end{equation}
All four functions vanish at ratio one, so $E_0=0$.
\begin{lemma}\label{ied:corr:lem:defect}
Suppose $h\le C_0/|z|$, the orbit is in the small-clock range, and $y_h=F_h(z)/z$ remains in the fixed cone. Then
\begin{equation}\label{ied:corr:eq:weighted}
 |E_{h+1}-E_h|\le
 \frac{C|z|^2|E_h|}{(1+|y_h|)^2}+C\bigl(|F_h(z)|^4+|z|^4\bigr).
\end{equation}
Consequently, on the entrance disks of \eqref{ied:corr:eq:entrance},
\begin{equation}\label{ied:corr:eq:Cerror}
 \mathcal C_z(F_{n-K}(z))=n-K+\Oh(K^{-3}+n^{-3}).
\end{equation}
\end{lemma}
\begin{proof}
Put $Y=1+|y|$ on a slightly wider cone. There $|d_y|\asymp Y^2$ and $|w|+|z|\asymp|z|Y$. The definitions and sector expansions give, for the fixed derivative orders required below,
\begin{align*}
 J^{(r)}&=\Oh(Y^{-r-1}),& L^{(r)}&=\Oh(Y^{-r}),\\
 M^{(r)}&=\Oh(Y^{1-r}),& N^{(r)}&=\Oh(Y^{2-r})\qquad(r\ge1),
\end{align*}
and $J=\Oh(1)$, $L=\Oh(1+\log Y)$, $M=\Oh(Y)$, $N=\Oh(Y^2)$. Since $\Delta y=\Oh(|z|Y^2)$, its relative size is $\Oh(|w|+|z|)$; all Taylor segments stay in the larger cone for a sufficiently small clock radius.

Temporarily replace $R_h$ by its polynomial through degree five, keeping a bounded independent $t$. The discarded $\Oh(z^6)$ forcing changes the clock by $\Oh(z^6/(|w|+|z|)^2)=\Oh(z^4)$. Expand $J/z$ to fourth order in the displacement, $L$ to third, $zM$ to second, and $z^2N$ to first. Their omitted terms are bounded respectively by
\[
 |z|^{-1}Y^{-6}(|z|Y^2)^5,\quad
 Y^{-4}(|z|Y^2)^4,\quad
 |z|Y^{-2}(|z|Y^2)^3,\quad
 |z|^2(|z|Y^2)^2.
\]
Each is $\Oh(|z|^4Y^4)$. The step remainder $\Oh(|z|^5Y^6)$ in \eqref{ied:corr:eq:ystep}, multiplied by the leading ratio derivative $\Oh(|z|^{-1}Y^{-2})$, has the same bound.

The retained clock increment is exactly
\begin{align*}
 1+z\frac{t-J}{d_y}
 +z^2\{Q_2(y,t)+aM'\}
 +z^3\{Q_3(y,t)+aN'\}+\Oh(|z|^4Y^4).
\end{align*}
Set $t_0=z\mathcal C_z(w)=J+zL+z^2M+z^3N$. It is bounded, and
$t_0-J=\Oh(|z|(1+\log Y))$ because $|z|Y$ is small. Expansion in $t$ about $J$ gives the terms
\begin{align*}
 z^2\{L/d_y+Q_2(y,J)+aM'\}
 +z^3\{M/d_y+L\partial_tQ_2(y,J)+Q_3(y,J)+aN'\}.
\end{align*}
They vanish by \eqref{ied:corr:eq:Mrecursive} and \eqref{ied:corr:eq:N}. All discarded $t$-Taylor terms are $\Oh(|z|^4Y^4)$: representative largest bounds are $|z|^4Y^3$ and $|z|^4(1+\log Y)^2Y^2$, both admissible.

The actual variable satisfies $t=hz=t_0-zE_h$ exactly. Both $t$ and $t_0$ are bounded before any bound on $E_h$ is used, so their connecting segment is bounded. On it $\partial_tR_{\rm polynomial}=\Oh(z^3)$ and $\partial_w\mathcal C_z=\Oh((|w|+|z|)^{-2})$. The derivative of the increment with respect to $t$ is therefore $\Oh(|z|/Y^2)$. Multiplication by $-zE_h$ gives the feedback term in \eqref{ied:corr:eq:weighted}; $|z|^4Y^4\asymp|w|^4+|z|^4$ gives the remainder.

Discrete Gronwall has multiplier at most $\exp(Cn|z|^2)=\Oh(1)$. The already established entrance orbit gives
\[
 \sum_{h<n-K}|F_h(z)|^4=\Oh(K^{-3}),\qquad n|z|^4=\Oh(n^{-3}).
\]
Together with $E_0=0$ this proves \eqref{ied:corr:eq:Cerror}. The orbit bounds come from the old clock, so the improved defect is not being used to bootstrap its own hypotheses.
\end{proof}

\section{Finite Fatou expansion and forced matching}\label{ied:corr:sec:fatou}
\begin{lemma}\label{ied:corr:lem:fatoufinite}
In fixed smaller coordinate sectors,
\begin{equation}\label{ied:corr:eq:phifinite}
 \Phi(w)=-\frac2w+\frac13\log w-\frac w{36}+\frac{w^2}{540}+\Oh(w^3).
\end{equation}
For
\begin{equation}\label{ied:corr:eq:U}
 U_0(w)=\frac2{3w^3}+\frac5{6w^2}+\frac1{6w},
\end{equation}
one has $U_0(f(w))-U_0(w)=-\Phi'(w)/2+\Oh(w)$.
\end{lemma}
\begin{proof}
Let $\Phi_2$ be the four displayed terms in \eqref{ied:corr:eq:phifinite}. Direct series division gives
\[
 \Phi_2(f(w))-\Phi_2(w)-1=-w^4/5184+\Oh(w^5).
\]
Sum this defect along the inverse orbit from Section~\ref{ied:corr:sec:interface}. Its absolute total is bounded by
$C\sum_{j\ge0}(|w|^{-1}+j)^{-4}=\Oh(w^3)$. This constructs a normalized Abel coordinate with \eqref{ied:corr:eq:phifinite}. Its difference from the original normalized coordinate is inverse-map invariant and tends to zero along inverse iterates, hence vanishes. Sector Cauchy estimates justify differentiated remainders. A second finite calculation gives
\[
 U_0(f(w))-U_0(w)+\Phi_2'(w)/2=-w/1080+\Oh(w^2).
\]
Replacing $\Phi_2'$ by $\Phi'$ changes this by $\Oh(w^2)$. Both lower pole terms in $U_0$ are needed, including for cancellation of the constant term.
\end{proof}

\begin{writenote}[Generalized in Part~III, 6~October 2026]
Part~III's Lemma~\ref{ied:log:lem:finitefatou} proves the expansion to every order, with rational $a_j$ and differentiated remainders, by the same summation of the Abel defect along inverse orbits; its $a_1=-1/36$, $a_2=1/540$ are the coefficients here.
\end{writenote}

Combining the expansions of $J,L,M,N$ gives, when $|w/z|$ is large in the clock cone,
\begin{align}\label{ied:corr:eq:newmatch}
 \mathcal C_z(w)={}&H(z)+\kappa z+\Phi(w)+z^2U_0(w)\notag\\
 &+\Oh\left(|w|^3+|z|^2(1+\log|w/z|)
                +\frac{|z|^3}{|w|^3}+\frac{|z|^4}{|w|^5}\right).
\end{align}
For example, the $z^2/w^3$, $z^2/w^2$, and $z^2/w$ terms come respectively from $J$, $L$, and $zM$ and assemble exactly into $z^2U_0$. Lower $M,N$ tails are absorbed by the displayed remainder in the small-clock range.

Take
\begin{equation}\label{ied:corr:eq:scales}
 K=\lfloor n^{1/2}\rfloor,\qquad z\asymp n^{-1},\qquad
 w=F_{n-K}(z)\asymp K^{-1}.
\end{equation}
The four matching errors in \eqref{ied:corr:eq:newmatch} have orders at most
$K^{-3}$, $n^{-2}\log n$, $n^{-3}K^3$, and $n^{-4}K^5$; each is $\Oh(n^{-3/2})$. The clock error in \eqref{ied:corr:eq:Cerror} has the same order. Thus, for $\Theta_z=\Phi+z^2U_0$,
\begin{equation}\label{ied:corr:eq:newentry}
 \Theta_z(F_{n-K}(z))=n-K-H(z)-\kappa z+\Oh(n^{-3/2}).
\end{equation}
The sign of the new phase constant is negative because $+\kappa z$ is moved to the other side of the clock identity.

For an actual late step, use the Abel equation before Taylor expansion:
$\Phi(w_{h+1})=\Phi(w_h+R_h)+1$. Since $R_h=z^2/2+\Oh(z^3)$, Lemma~\ref{ied:corr:lem:fatoufinite} yields
\begin{equation}\label{ied:corr:eq:cocycle}
 \Theta_z(w_{h+1})-\Theta_z(w_h)-1
 =\Oh\left(|z|^2|w_h|+\frac{|z|^3}{|w_h|^2}
                             +\frac{|z|^4}{|w_h|^4}\right).
\end{equation}
Indeed the $z^2\Phi'/2$ term cancels the leading $U_0$ difference. The second-order Taylor error for $\Phi$ is $\Oh(z^4/w_h^3)$, absorbed by the last term, and perturbing the argument in $U_0$ gives $\Oh(z^4/w_h^4)$.

\section{Uniform local limit on a growing strip}\label{ied:corr:sec:strip}
\begin{proposition}\label{ied:corr:prop:correctedlocal}
For every fixed real-width bound $C_0$, uniformly at
\[
 z=\frac T{D_n-s},\qquad |\Rea s|\le C_0,\qquad
 |\Ima s|\le n^{1/20}+2,
\]
one has
\begin{equation}\label{ied:corr:eq:localcorrected}
 F_n(z)=P\bigl(n-H(z)-\kappa z\bigr)+\Oh(n^{-3/2}).
\end{equation}
\end{proposition}
\begin{proof}
Write $M_n=n^{1/20}+2$ and use $K$ in \eqref{ied:corr:eq:scales}. The distance to $r=T/D_n$ is $\Oh(M_n/n^2)$, so the whole parameter set lies in the old entrance disk of radius $qK/n^2$. All entrance and new matching estimates are uniform there.

Let $\alpha=n-H(z)-\kappa z$. The exact identity
\[
 H\left(\frac T{D_n-s}\right)=n-s+\frac13\log\frac{D_n-s}{n}
\]
gives $\alpha=s+\Oh((\log n+M_n)/n)$; its real part stays bounded and its imaginary part differs from $\Ima s$ by $o(1)$.

Choose one fixed integer $L$ large relative to $C_0$ and the inverse-sector threshold. It does not depend on $M_n$. For $0\le j\le K-L$, set $m=K-j$ and use the prototype phase $\alpha-m$. Its real part is at most $C_0+1-m<0$ and
\[
 |\alpha-m|\asymp m+|\Ima s|.
\]
It lies in a fixed interior inverse sector about the closed left half-plane. A bootstrap phase error at most one preserves fixed angular and radial margins. The corresponding inverse values satisfy
\begin{equation}\label{ied:corr:eq:latevalues}
 |w_j|\asymp(m+|\Ima s|)^{-1}.
\end{equation}
The actual $\Oh(n^{-2})$ forcing changes the uncorrected $\Phi$ phase by at most $Cn^{-2}(m+M_n)^2$. Its total is
\[
 \Oh\left(n^{-2}\sum_{m=L}^K(m+M_n)^2\right)
       =\Oh(K^3/n^2)=o(1),
\]
because $M_n=o(K)$. The old entry error in \eqref{ied:corr:eq:oldphase} is also $o(1)$. These strictly improve the bootstrap for large $n$. The inverse-identification clause ensures every perturbed point uses the same branch. There is no new injectivity assumption.

The bounds just obtained imply the sums needed for the corrected phase:
\[
 \sum|w_j|=\Oh(\log K),\qquad
 \sum|w_j|^{-2}=\Oh(K^3),\qquad
 \sum|w_j|^{-4}=\Oh(K^5).
\]
Summing \eqref{ied:corr:eq:cocycle} to time $n-L$ gives
\[
 \Oh(n^{-2}\log n+n^{-3}K^3+n^{-4}K^5)=\Oh(n^{-3/2}),
\]
of the same order as the error in \eqref{ied:corr:eq:newentry}. Removing $z^2U_0$ at time $n-L$ costs
$\Oh(n^{-2}(1+M_n)^3)=\Oh(n^{-37/20})$. Hence
\[
 \Phi(F_{n-L}(z))=\alpha-L+\Oh(n^{-3/2}).
\]
The inverse derivative is uniformly bounded on the shifted strip, by analyticity on its bounded part and the sector estimate on its unbounded part. We obtain
$F_{n-L}(z)=P(\alpha-L)+\Oh(n^{-3/2})$.

The values $P(\alpha-L)$ form a bounded set and tend to zero at large imaginary height. Its closure is compact. The fixed entire composition $f^{\circ L}$ and its derivatives are bounded on a compact enlargement, so the remaining $L$ iterates preserve the error; their subsidiary forcing contributes only $\Oh(n^{-2})$. This proves \eqref{ied:corr:eq:localcorrected}. Large imaginary parts remain in the inverse sector, which is why one fixed $L$ suffices.
\end{proof}

\subsection{Exponential parameter and explicit remainders}
Set $d=-\frac13\log n+d_\star$, $D_n=n+d$, $r=T/D_n$, and
\[
 G_n(s)=F_n(r\e^{s/n}),\qquad
 Q_n(s)=(d+1/3)s-s^2/2-d/3-\kappa T.
\]
The exact phase identity is
\begin{align}\label{ied:corr:eq:phaseexact}
 n-H(r\e^{s/n})-\kappa r\e^{s/n}
 ={}&D_n(1-\e^{-s/n})-\tfrac13\log(1+d/n)\notag\\
    &+s/(3n)-\kappa T\e^{s/n}/(n+d).
\end{align}
It implies
\begin{equation}\label{ied:corr:eq:phaseerror}
 n-H(r\e^{s/n})-\kappa r\e^{s/n}
 =s+Q_n(s)/n+R_n(s),\qquad
 |R_n(s)|\le Cn^{-2}(1+|d|+|s|)^3.
\end{equation}
The rational parameter for $r\e^{s/n}$ is $D_n(1-\e^{-s/n})$; this preserves a slightly enlarged fixed-width strip and increases its height by $o(1)$. Both $P'$ and $P''$ are uniformly bounded on such strips. Taylor expansion in \eqref{ied:corr:eq:phaseerror} has additional error bounded by
$Cn^{-2}(1+|d|+|s|)^4$. On $|\Ima s|\le n^{1/20}+2$ it is smaller than $n^{-3/2}$. Proposition~\ref{ied:corr:prop:correctedlocal} therefore gives
\begin{equation}\label{ied:corr:eq:Gexpand}
 G_n(s)=P(s)+\frac{Q_n(s)P'(s)}n+\Oh(n^{-3/2}).
\end{equation}
Cauchy's estimate on fixed-radius disks inside the enlarged strip gives the same error for every fixed derivative, in particular for derivative thirty. These degree-three and degree-four bounds avoid an unspecified polynomial whose degree could compromise the growing-strip argument.

\section{Derivative extraction and classical Fourier inversion}\label{ied:corr:sec:extract}
For every nonnegative integer $q$, termwise coefficient extraction gives exactly
\begin{equation}\label{ied:corr:eq:extract}
 n a_n r^n=\frac1{2\pi}\int_{-\pi n}^{\pi n}
                  G_n^{(q)}(\ii t)\e^{-\ii t}\dd t.
\end{equation}
Indeed a term of degree $k$ acquires $(k/n)^q$, and orthogonality over this full period selects $k=n$, whose factor is one. Analyticity in a neighborhood of the circle justifies the termwise step. No integration by parts or endpoint assumption is needed.

\subsection{Transport of analytic disks}
The estimates imported from Report 225 are disk estimates, which permit arbitrary fixed derivative order. Choose a sufficiently small fixed $a>0$. For $|t|\le\delta n$,
\[
 |s-\ii t|\le a|t|\quad\Longrightarrow\quad
 |r\e^{s/n}-r\e^{\ii t/n}|
 \le r\e^{a\delta}a|t|/n\le b|t|/n^2.
\]
Thus \eqref{ied:corr:eq:oldtaildisk} and Cauchy give
\begin{equation}\label{ied:corr:eq:highderivtail}
 |G_n^{(30)}(\ii t)|\le C_{30}|t|^{-31}
       \qquad(M_0\le|t|\le\delta n).
\end{equation}
For the outer angles, the fixed neighborhoods in $\zeta=nz$ admit $s$ disks of radius $a_0n$ for one fixed small $a_0>0$. The bound $G_n=\Oh(n^{-1})$ there gives $G_n^{(30)}=\Oh(n^{-31})$. Disks crossing the endpoints $\pm\pi n$ cause no problem because the exponential parametrization is periodic.

Put $M=n^{1/20}$. The finite-height intermediate and outer integrals outside $[-M,M]$ are
\[
 \Oh(M^{-30}+n^{-30})=\Oh(n^{-3/2}).
\]
The central integral of the differentiated error in \eqref{ied:corr:eq:Gexpand} is
\[
 \Oh(Mn^{-3/2})=\Oh(n^{-29/20}).
\]
For the limiting transforms, \eqref{ied:corr:eq:Pderivative} and Leibniz give
\[
 (Q_nP')^{(30)}(\ii t)
     =\Oh(|t|^{-30}+\log n\,|t|^{-31})
\]
for large $|t|$. Its omitted tail, divided by $n$, is smaller still, and that of $P^{(30)}$ is $\Oh(M^{-30})$. Both central integrals can consequently be replaced by full-line integrals with error $\Oh(n^{-7/5})$. The margin between $29/20$ and $7/5$ is deliberate.

\subsection{Signs and derivative weights in inversion}
Write $\mathcal I[V](x)=(2\pi)^{-1}\int_\R V(\ii t)\e^{-\ii tx}\dd t$ whenever the integral is absolutely convergent. We only apply it to the differentiated transforms with $q=30$. At positive $x$, Fourier inversion and its justified ordinary derivatives give
\begin{align}\label{ied:corr:eq:inversions}
 \mathcal I[P^{(q)}](x)&=x^qh(x),\notag\\
 \mathcal I[(P')^{(q)}](x)&=x^{q+1}h(x),\notag\\
 \mathcal I[(sP')^{(q)}](x)&=-(x^{q+1}h)'(x)+qx^qh(x),\\
 \mathcal I[(s^2P')^{(q)}](x)&=(x^{q+1}h)''(x)
 -2q(x^qh)'(x)+q(q-1)x^{q-1}h(x).\notag
\end{align}
For example,
$(sP')^{(q)}=sP^{(q+1)}+qP^{(q)}$ and
$(s^2P')^{(q)}=s^2P^{(q+1)}+2qsP^{(q)}+q(q-1)P^{(q-1)}$.
The $t$- and $t^2$-weighted integrals are absolutely convergent by \eqref{ied:corr:eq:Pderivative}; multiplication by $s$ corresponds to minus ordinary differentiation in $x$ for this Laplace sign convention. These identities invoke no endpoint distributions and no regularity at zero.

At $x=1$, the four inversions in \eqref{ied:corr:eq:inversions} are respectively
\[
 h(1),\quad h(1),\quad -h(1)-h'(1),\quad 2h'(1)+h''(1).
\]
In particular the $h(1)$ coefficient in the last expression cancels. Substituting into \eqref{ied:corr:eq:extract} and \eqref{ied:corr:eq:Gexpand} gives
\begin{align}\label{ied:corr:eq:coefficientcorrected}
 n a_n r^n={}&h(1)+\frac1n\bigl\{(-d/3-\kappa T)h(1)\notag\\
 &\hspace{1em}-(d+1/3)(h(1)+h'(1))-h'(1)-h''(1)/2\bigr\}
 +\Oh(n^{-7/5}).
\end{align}
Finally,
\begin{align*}
 r^{-n}&=(n/T)^n\e^d\left(1-\frac{d^2}{2n}
                         +\Oh(\log^4n/n^2)\right),\\
 \e^d&=n^{-1/3}T^{1/3}\e^{-B}.
\end{align*}
Multiplication, division by the positive $h(1)$, and substitution of $d$ give \eqref{ied:corr:eq:Qd}, hence \eqref{ied:corr:eq:main} and \eqref{ied:corr:eq:AC}. Cross-products and the radius remainder are $\Oh(\log^4n/n^2)$ and are smaller than the stated error. This completes the proof of Theorem~\ref{ied:corr:thm:diagonal}.

\section{Uniform proportional height correction}\label{ied:corr:sec:height}
\begin{proof}[Proof of Theorem~\ref{ied:corr:thm:height}]
Use the height-scaled variables $r=T/(m+d_m)$ and $G_m(s)=F_m(r\e^{s/m})$. The exact period identity at $\beta=n/m$ is
\begin{equation}\label{ied:corr:eq:betaextract}
 m\beta^q b_{n,m}r^n=\frac1{2\pi}\int_{-\pi m}^{\pi m}
          G_m^{(q)}(\ii t)\e^{-\ii\beta t}\dd t,\qquad q=30.
\end{equation}
A degree-$k$ term acquires $(k/m)^q$, and the full period selects $k=n$, yielding exactly $\beta^q$. All growing-strip and tail estimates depend on height $m$, not on the target degree. Changing the oscillatory factor has modulus one for real $\beta$; division by $\beta^q$ is uniformly harmless on $I$.

Evaluate \eqref{ied:corr:eq:inversions} at $x=\beta$ and divide by $\beta^q$. The first result is $h(\beta)$; the correction terms become
\[
 \beta h(\beta),\qquad -h(\beta)-\beta h'(\beta),\qquad
 2h'(\beta)+\beta h''(\beta).
\]
For example the third identity simplifies to $-h-xh'$ after division by $x^q$, and in the fourth the coefficient of $h$ cancels identically, leaving $2h'+xh''$. Hence
\begin{align*}
 m b_{n,m}r^n={}&h(\beta)+\frac{J_\beta(d_m)}m+\Oh_I(m^{-7/5}),\\
 J_\beta(d)={}&(-d/3-\kappa T)\beta h(\beta)
 -(d+1/3)(h(\beta)+\beta h'(\beta))\\
 &-h'(\beta)-\beta h''(\beta)/2.
\end{align*}
Since $n=\beta m$ exactly,
\[
 r^{-n}=(m/T)^n\e^{\beta d_m}
 \left(1-\frac{\beta d_m^2}{2m}+\Oh_I(\log^4m/m^2)\right),
 \quad \e^{\beta d_m}=m^{-\beta/3}T^{\beta/3}\e^{-\beta B}.
\]
Multiplication gives $\mathcal Q_\beta=J_\beta/h-\beta d_m^2/2$ in \eqref{ied:corr:eq:height}. Continuity and positivity imply $\min_Ih>0$; the derivatives are bounded on $I$. This converts the uniform absolute coefficient error into the asserted relative error. At $\beta=1$ the expression reduces exactly to Theorem~\ref{ied:corr:thm:diagonal}.
\end{proof}

\begin{writenote}[Checked, 6~October 2026]
The write checked the reduction: $\mathcal Q_1(d)$ is \eqref{ied:corr:eq:Qd}. The theorem refines Part~I's Theorem~\ref{ied:lead:thm:proportional}, whose leading term it contains.
\end{writenote}

The exact-height formula can be used directly for $m=n+k$, for square-root height windows, or for $m=\lfloor\lambda n\rfloor$ with fixed $\lambda>0$. In the last case the leading floor phase remains $\exp(-\{\lambda n\}/\lambda)$, as in Report 225. Formula \eqref{ied:corr:eq:height} retains all such dependence by keeping the exact $m$ and $n/m$; replacing them prematurely by their limits can discard contributions as large as the correction itself.

\section{A sharpened inverse formula}\label{ied:corr:sec:inverse}
Set $\alpha=4/3$, $\ell=\log c$, and
\[
 Q(u)=-u^2/18+Au+C.
\]
Theorem~\ref{ied:corr:thm:diagonal} implies
\begin{equation}\label{ied:corr:eq:logmodel}
 \log a_n=n\log(n/T)-\alpha\log n+\ell
               +\frac{Q(\log n)}n+\Oh(n^{-7/5}),
\end{equation}
because the square of the first relative correction is $\Oh(\log^4n/n^2)$.

\begin{corollary}\label{ied:corr:cor:inverse}
For an exact sequence value $y=a_n$, let $Y=\log y$ and, for large $Y$, define
\begin{equation}\label{ied:corr:eq:inversevars}
 X=\frac{Y}{W(Y/T)},\qquad U=\log X,\qquad q_X=\log(X/T)+1,
\end{equation}
where $W$ is the positive real Lambert branch. Put
\begin{equation}\label{ied:corr:eq:inversecoeffs}
 \delta_0=\frac{\alpha U-\ell}{q_X},\qquad
 \delta_1=-\frac{\delta_0^2/2-\alpha\delta_0+Q(U)}{q_X}.
\end{equation}
Then
\begin{equation}\label{ied:corr:eq:inverse}
 n=X+\delta_0+\frac{\delta_1}{X}
          +\Oh\left(\frac{X^{-7/5}}{\log X}\right).
\end{equation}
\end{corollary}
\begin{proof}
The definition of $W$ gives $X\log(X/T)=Y$. The leading equivalent first gives $n=X+\Oh(1)$. Define the explicit smooth model
\[
 F(x)=x\log(x/T)-\alpha\log x+\ell+Q(\log x)/x.
\]
For bounded $\delta$, Taylor expansion about $X$ gives uniformly
\begin{align*}
 F(X+\delta)-Y={}&q_X\delta-\alpha U+\ell\\
 &+\frac{\delta^2/2-\alpha\delta+Q(U)}X
 +\Oh((1+U^2)/X^2).
\end{align*}
Now $\delta_0=\Oh(1)$ and $\delta_1=\Oh(U)$, so substituting
$\delta=\delta_0+\delta_1/X$ cancels both retained orders. The residual is
$\Oh((1+U^2)/X^2)=\Oh(X^{-7/5})$. At the true integer $n$, \eqref{ied:corr:eq:logmodel} gives $F(n)-Y=\Oh(n^{-7/5})$. On the bounded comparison interval $F'(x)\sim\log X$. The mean value theorem proves \eqref{ied:corr:eq:inverse}. Only $F$ is differentiated; the asymptotic remainder is used solely at the true integer.
\end{proof}
The constants $\delta_0,\delta_1$ in this inverse depend on $X$ and are distinct from the fixed clock constant $d_\star$. The large real solution of $F(x)=\log y$ has the same error at $y=a_n$. With exact constants this yields eventual nearest-integer recovery, but neither an effective onset nor a certificate for arbitrary thresholds between consecutive sequence values. A fixed numerical error in $\log c$ creates an index error of order $1/\log X$, much larger than the fine rate in \eqref{ied:corr:eq:inverse}; the diagnostic decimal must not be substituted as if exact.

\begin{writenote}[Relation to Part~I and to the transseries volume, 6~October 2026]
$\delta_0$ is the correction of Part~I's \eqref{ied:lead:eq:inverse}; the write rechecked the two orders \eqref{ied:corr:eq:inversecoeffs} by expanding $F(X+\delta_0+\delta_1/X)-Y$. As in Part~I, $X$ is an instance of \texttt{p0:prop:factorial-core} of \cite{TSvol}; the two Newton-type corrections are proved directly and are not instances of \texttt{p0:thm:perturbed-inversion} (front matter).
\end{writenote}

\section{Exact computation and numerical diagnostics}\label{ied:corr:sec:computation}
\subsection{Exact recurrence and finite checks}
Writing $1+F_{h+1}(z)=\sum_{k\ge0}c_kz^k$ and logarithmically differentiating the Euler product gives
\begin{equation}\label{ied:corr:eq:exactrecurrence}
 b_j=\sum_{d\mid j}d[z^d]F_h(z),\qquad
 c_0=1,\qquad c_k=\frac1k\sum_{j=1}^k b_jc_{k-j}.
\end{equation}
Degree truncation is exact. The public code checks divisibility before every integer division. An independent finite-product implementation uses the binomial expansion of $(1-z^d)^{-a}$ and agrees through degree and height twelve. A separately generated exact fixture contains all diagonal values through $n=414$; only the first twenty-two displayed OEIS terms are an external sequence comparison. The fixture is not presented as a retrieved OEIS b-file.

The code also verifies finite formal-coordinate identities and the new clock algebra exactly over rational symbolic expressions. Such checks protect against algebraic transcription errors; they do not prove analytic remainder estimates, sector entrance, or the all-orders convergence of any formal series. All pass/fail checks use explicit exceptions and remain active under Python optimization. Local chunked decimal serialization handles long exact integers under \texttt{PYTHONINTMAXSTRDIGITS=640} without changing the global limit. The implementation caps and tested ranges are stated in the archive README.

\subsection{Uncertified amplitude and derivative experiments}\label{ied:corr:sub:diagnostics}
The accompanying optional floating-point programs solve a truncated Fatou inverse, take finitely many forward iterates, propagate derivatives, and numerically invert the transform. For example, $P'''$ is the transform of $x^3h(x)$. Its absolutely convergent inversion and its first two ordinary $x$ derivatives give $h(1)$, $h'(1)$, and $h''(1)$. Subtracting $12/(1-s)^4$ improves the observed tail; its inverse is $2x^3\e^{-x}$, with values and derivatives at one equal to $2/\e$, $4/\e$, and $2/\e$. After subtraction and quadrature, the derivative weights are undone by
\[
 h'= (x^3h)'|_1-3h(1),\qquad
 h''=(x^3h)''|_1-6h'(1)-6h(1).
\]
These are diagnostic computations with no interval error control for the truncated coordinate, inverse Newton step, forward propagation, finite contour, or quadrature.

Representative numerical values, not certified digit claims, are
\begin{align*}
 h(1)&\approx1.81460956,& \rho_1&\approx-0.43688195,&
 \rho_2&\approx-0.15328008,\\
 \kappa&\approx-0.25434817,& A&\approx0.60098418,&
 C&\approx-0.49816124.
\end{align*}
The associated leading amplitude is approximately $4.4923299$. At $n=414$, floating-point normalization of the exact integer gives
\[
 a_n(T/n)^n n^{4/3}\approx4.50436880933,
\]
while the unfitted first-correction prediction gives approximately $4.50433122939$, a difference about $3.76\times10^{-5}$. No coefficient was fitted to these finite sequence values. Agreement and parameter sensitivity are useful diagnostics, not proofs of a remainder or of displayed digits.

\begin{writenote}[Checks and a comparison, 6~October 2026]
At intake, with these constants the corrected prediction matched the exact $a_n$ with residual $1.4\times10^{-4}$ at $n=200$, against $2.5\times10^{-2}$ for the leading term. With Part~I's diagnostic $h(1)$, $2^{1/3}h(1)=2.2862647816$ agrees to eight significant digits with the iterated Bell report's exploratory amplitude $I(1)$ (Remark~\ref{ied:rem:ibdnumbers}); both values are uncertified.
\end{writenote}

\subsection{Reproducible archive}
The archive contains this complete editable TeX, exact programs and tests, optional diagnostic programs and recorded examples, the exact fixture, source provenance, and a deterministic source-preserving builder. It also includes the unchanged Report 225 PDF and TeX as a clearly labeled supporting reference. The new report does not modify that earlier deliverable.

The builder validates its source manifest, runs the exact checks, creates a private pdf\LaTeX\ format, compiles without shell escape, rejects overfull boxes, missing glyphs, and unresolved references, checks that source bytes were preserved, and writes a new PDF and sorted fixed-timestamp ZIP to a separate new directory. Ordinary and optimized runs under the restricted digit limit are compared as actual complete ZIPs and member bytes. Identical-toolchain byte reproducibility is distinct from mathematical correctness; different TeX or font versions may change PDF bytes. Every rendered page is also inspected before release.

\begin{writenote}[The companion here, 6~October 2026]
Report~226's own files are shipped as \file{code/226-firstcorr-*.py}, \file{data/226-firstcorr-*} and \file{226-firstcorr-sources-SOURCE_LEDGER.md}. Its eight byte copies of Report~225's files and its \file{supporting/} copies of Report~225 are shipped once, as Report~225's files and as Part~I. Its builder and guard tests expect those files beside them; the report README gives the rerun route from the arrival commit.
\end{writenote}

\section{Further questions and the boundary of the theorem}\label{ied:corr:sec:further}
\subsection{A genuine all orders expansion}\label{ied:corr:sub:allorders}
The natural formal expectation is
\[
 \frac{a_n}{c(n/T)^nn^{-4/3}}
 \sim1+\sum_{r\ge1}\frac{P_r(\log n)}{n^r},\qquad \deg P_r\le2r.
\]
Only $P_1(X)=-X^2/18+AX+C$ is proved here. If an all-orders construction establishes the stated degree bound and proves that no further contribution reaches degree $2r$, the leading coefficient would be
\[
 [X^{2r}]P_r(X)=\frac{(-1/18)^r}{r!}.
\]
In particular the formal second-order highest term would be $\log^4n/(648n^2)$. This conditional inference is not an all-orders theorem and does not follow from the first correction alone.

A proof needs an induction controlling arbitrary fixed scaled subsidiary order, large-ratio clock expansions and resonant logarithms, forced Fatou cohomological corrections, and growing-strip errors with extraction derivative order chosen for each target precision. The $\log y/540$ in $N$ already shows why clocks cannot be treated as a plain power series in $z$: a $z^2\log(1/z)$ contribution can enter the global phase. Exponentiating the radius alone misses this structure.

\begin{writenote}[Still open, 6~October 2026]
Not settled by Part~III, which works at one polynomial rate in $m$ (Subsection~\ref{ied:log:sub:other}). The conditional coefficient $(-1/18)^r/r!$ is recorded as unproved in Section~\ref{ied:sec:further}, item~\ref{ied:q:higher}.
\end{writenote}

\subsection{Effective constants and certified inversion}\label{ied:corr:sub:effective}
Validated values of $h(1)$ and its derivatives require explicit sector remainder bounds, an enclosed inverse branch, controlled forward derivative propagation, interval quadrature, and an integrable certified tail. The present proof supplies qualitative constants and no finite onset. The exact formula for $\kappa$ is elementary apart from Catalan's constant, but certifying the full correction still requires the density derivative calculations.

An effective relative-error bound could turn the scalar inverse model into a certified threshold algorithm by combining exact small-index counts with rigorous upper and lower bounds. Theorem~\ref{ied:corr:thm:diagonal} alone does not decide where a general input lies between $a_n$ and $a_{n+1}$. One should separate this computational threshold problem from eventual rounding at exact sequence values.

\begin{writenote}[Still open, 6~October 2026]
Section~\ref{ied:sec:further}, items~\ref{ied:q:validated} and~\ref{ied:q:effective}.
\end{writenote}

\subsection{Endpoint height regimes and density structure}\label{ied:corr:sub:endpoint}
The compact-positive ratio theorem leaves open $n/m\to0$ and $n/m\to\infty$. These require endpoint information about $h$ and renewed uniform contour estimates. Smooth strict positivity away from zero supplies neither an expansion at zero nor a large-argument asymptotic. The positive-measure identity from Report 225,
\[
 (\e^x-1)\nu=\sum_{k\ge2}\nu^{*k}/k!,\qquad \nu(\dd x)=h(x)\dd x,
\]
is a possible starting point, but exchanging endpoint limits or many derivatives with this convolution series requires separate arguments.

A transition theory connecting fixed-degree polynomiality, proportional height, and fixed-height partition asymptotics would also be useful. The forcing calculations here demonstrate that another initial generating function or multiset recurrence cannot simply inherit the same clock constants from the autonomous map. Its subsidiary terms and matching normalization must be recomputed.

\begin{writenote}[Partly answered by Part~III, 6~October 2026]
Part~III treats $n/m\to0$ at leading order on logarithmic bands (Theorem~\ref{ied:log:thm:transfer}) and proves an expansion of $h$ at zero (Proposition~\ref{ied:log:prop:hzero}, Theorem~\ref{ied:log:thm:germ}, Corollary~\ref{ied:log:cor:densityseries}), so ``an expansion at zero'' is no longer missing. Open: $n/m\to\infty$ and a large-argument asymptotic of $h$, the correction on shrinking bands, the transition to fixed degree, other recurrences (Section~\ref{ied:sec:further}, items~\ref{ied:q:endpoint}, \ref{ied:q:polysmall}, \ref{ied:q:other}).
\end{writenote}

\clearpage
\part{Logarithmic height crossover for iterated Euler transforms}\label{ied:log:part}
\parteqs{III.}
\partsource{Bundle Report~228, archive \file{Report228.zip}, delivered as \file{src/report228.tex} (17 pages). Delivered title block ``Logarithmic height crossover for iterated Euler transforms / Report 228 / 5 October 2026''. Printed as delivered, with \textbf{[write]} notes. Delivered Section~$k$ is Section~$k+23$ here, its statement $k.j$ is $(k+23).j$, and its equation $(k)$ is $(\mathrm{III}.k)$. Section~\ref{ied:sec:further}, which closes the report, is added in the write.}
\begin{center}
\small
\begin{tabular}{@{}>{\raggedright\arraybackslash}p{0.46\textwidth} >{\raggedright\arraybackslash}p{0.46\textwidth}@{}}
\toprule
Reading conventions of Part~III & Elsewhere in the report\\
\midrule
``Reports 225 and 226'', \cite{R225}, \cite{R226} & Parts~I and~II\\
$m/(n\log n)\to\lambda$; $L=\log n$; $R(n,m)$ the ratio & Part~I: $m=\lfloor\lambda n\rfloor$; $L(y)$ the clock; $R_h(z)$\\
$K$ a constant; $C_0=1-\gamma-\log2$; $\mathcal C$ Catalan; $C(\beta)$ the germ & Parts~I, II: $K$ an index; $C_0$ bounds; Part~II: $G$ Catalan\\
$\rho_j\in\mathbb Q$, $c_j$ log-density coefficients; $a_j$ Fatou coefficients; $b_j$ germ coefficients & Part~II: $\rho_j$ ratios of derivatives of $h$; all: $c_k$, $b_j$ in the exact recurrence\\
$g(t)=t^{t/3}h(t)$; $W(x)=P(-x)$; $\mathcal G(u)$ & Parts~I, II: $g(w)=\log(1+w)$; $W$ Lambert; Part~I: $G(u)$\\
$\tau=-\log r$; $\chi_n$; $y=m/n$ & Part~I: $\tau$ a time; Parts~I, II: $y=w/z$\\
\bottomrule
\end{tabular}
\end{center}
\begin{partabstract}
For the iterated Euler-transform array $A(n,m)$ in OEIS A290353, we prove
\[
 \frac{A(n,m)}{2(2/\pi)^n m^{n-1}}
 \longrightarrow \exp\!\left(-\frac1{3\lambda}\right)
 \quad\text{if}\quad \frac{m}{n\log n}\longrightarrow\lambda\in(0,\infty).
\]
The result extends the positive fixed-slope analysis of Reports 225 and 226 to degree-to-height ratios decreasing as any fixed inverse power of the logarithm. We prove an exact small-argument Hankel representation for their normalized positive density $h$, show that $t^{t/3}h(t)$ has a one-sided smooth germ, and obtain every fixed logarithmic endpoint order with differentiated remainders. A rational coefficient generator isolates a single zeta value at each order of the logarithmic density. Exact cancellation of two logarithms yields every fixed inverse-logarithm order in the crossover window, without $\log\log n$ terms. We also invert the smooth suppression model and quantify rounded integer-height prescriptions and observed-ratio inverses. The exact Kaneiwa leading-monomial normalization changes the first correction constant by $+1$.

This unrefereed research proof uses explicitly stated analytic interfaces from the two preceding reports, included unchanged in the archive. It does not assert convergence of an infinite endpoint series, an all-orders expansion in $1/m$, an effective onset, discrete monotonicity, or exact integer-height recovery. Numerical comparisons are unvalidated floating diagnostics built from exact finite coefficients.
\end{partabstract}
\section{The array and the crossover}
Define
\begin{equation}\label{ied:log:eq:recurrence}
 F_0(z)=z,\qquad
 F_{m+1}(z)=\exp\left(\sum_{j\ge1}\frac{F_m(z^j)}j\right)-1,
 \qquad A(n,m)=[z^n]F_m(z)\quad(n\ge1).
\end{equation}
This is the height-indexed array in OEIS A290353~\cite{A290353}. It counts unlabeled rooted trees whose leaves all have the prescribed height. Equivalently, the new root carries a nonempty unordered multiset of trees from the preceding height, and
\[
 1+F_{m+1}(z)=\prod_{d\ge1}(1-z^d)^{-A(d,m)}.
\]
The generating functions have constant term zero; the tabular convention $A(0,m)=1$ is imposed separately. There is no factorial normalization.

Put
\begin{align}\label{ied:log:eq:constants}
 T&=\frac\pi2,&
 B&=\frac12-\frac\pi6-\frac{\pi^2}{16}-\frac{\log2}{6},\\
 C_0&=1-\gamma-\log2,&
 K&=\frac{1-\gamma+\log(T/2)}3-B,\notag
\end{align}
where $\gamma$ is Euler's constant. Our primary normalization is
\begin{equation}\label{ied:log:eq:R}
 R(n,m)=\frac{A(n,m)}{2T^{-n}m^{n-1}}.
\end{equation}
All asymptotics with two indices below take $n,m\to\infty$ in their specified ranges.

\begin{theorem}\label{ied:log:thm:crossover}
If $m/(n\log n)\to\lambda\in(0,\infty)$, then
\begin{equation}\label{ied:log:eq:crossover}
 R(n,m)\longrightarrow\e^{-1/(3\lambda)}.
\end{equation}
More precisely, let $\Lambda\subset(0,\infty)$ be compact, put $L=\log n$, and set $m=\lfloor\lambda nL\rfloor$. For every fixed $N\ge1$, uniformly for $\lambda\in\Lambda$,
\begin{equation}\label{ied:log:eq:logcross}
 \log R(n,m)=-\frac1{3\lambda}+\frac{K}{\lambda L}
       +\sum_{j=2}^N\frac{c_j}{(\lambda L)^j}
       +\Oh_{\Lambda,N}(L^{-N-1}),
\end{equation}
where the sum is empty for $N=1$, and
\begin{equation}\label{ied:log:eq:cs}
 c_j=\rho_j-\frac{\zeta(j)}{j3^j},\qquad \rho_j\in\mathbb Q\quad(j\ge2).
\end{equation}
A finite coefficient algorithm determines every $\rho_j$. In particular,
\[
 c_2=\frac{9-\pi^2}{108},\qquad
 c_3=\frac{17}{1080}-\frac{\zeta(3)}{81}.
\]
\end{theorem}
For example, exponentiation gives
\begin{equation}\label{ied:log:eq:Rcross}
 R(n,m)=\e^{-1/(3\lambda)}
 \left\{1+\frac{K}{\lambda L}
 +\frac{K^2/2+(9-\pi^2)/108}{\lambda^2L^2}
 +\Oh_\Lambda(L^{-3})\right\}.
\end{equation}
Every order here means a fixed finite truncation followed by its asymptotic limit. It says neither that the infinite series converges nor that there is an all-orders expansion in powers of $m^{-1}$. One polynomial rate in the coefficient-transfer error will suffice for all fixed inverse-logarithm orders.

The smooth suppression-ratio inverses are proved in Section~\ref{ied:log:sec:inverse}. With $\tau=-\log r$ and $\chi_n=(\log n)/3-K$, their first height prescription is
\begin{equation}\label{ied:log:eq:inversepreview}
 \frac mn\approx\frac{\chi_n}{\tau}-\frac{c_2}{\chi_n}
 -\frac{c_3\tau}{\chi_n^2}
 -\frac{c_2^2\tau+c_4\tau^2}{\chi_n^3}.
\end{equation}
The inverse concerns a smooth asymptotic ratio; no unique discrete solution, least-height threshold, or exact recovery by rounding is claimed.

\begin{writenote}[Reading Part~III against Parts~I and~II, 6~October 2026]
``Report 225'', \cite{R225} and ``Report 226'', \cite{R226} are Parts~\ref{ied:lead:part} and~\ref{ied:corr:part}; their included copies are byte-identical and not shipped. Here $\lambda$ measures $m/(n\log n)$, not Part~I's slope $m/n$ of Corollary~\ref{ied:lead:cor:shifts}; $K$ is a constant, not Parts~I and~II's entrance index; $\rho_j$ are rational numbers, not Part~II's $h^{(j)}(1)/h(1)$; $L=\log n$, not the clock function.
\end{writenote}

\section{Imported analytic interfaces and the scope of the result}\label{ied:log:sec:interfaces}
The analytic foundation comes from Reports 225 and 226~\cite{R225,R226}. Their unchanged complete TeX and PDF files accompany this report, with hashes in the source ledger. We state exactly what is used; the present report supplies the endpoint, shrinking-ratio, and inverse arguments in full.

\subsection{The normalized coordinate and density}
Report 225, Sections 4 and 6, constructs an entire function $P$ with
\begin{equation}\label{ied:log:eq:coordinate}
 P(s+1)=\e^{P(s)}-1,\qquad
 \Phi(w)=-\frac2w+\frac13\Log w+\Oh(w).
\end{equation}
The coordinate $\Phi$ is analytic on every sufficiently small closed proper subsector of $|\arg w|<\pi$, with compatible branches on overlaps. It has zero additive constant in \eqref{ied:log:eq:coordinate}. The local inverse agrees with $P$ on every sufficiently distant proper subsector of $|\arg(-s)|<\pi$; in particular
\begin{equation}\label{ied:log:eq:inverseidentity}
 P(\Phi(w))=w,\qquad
 P'(s)=\frac2{s^2}+\Oh\left(\frac{\log|s|}{|s|^3}\right)
\end{equation}
there. For the reciprocal inverse iteration
\begin{equation}\label{ied:log:eq:Greciprocal}
 \mathcal G(u)=\frac2{\Log(1+2/u)}
             =u+1-\frac1{3u}+\Oh(u^{-2}),
\end{equation}
the normalized coordinate satisfies
\[
 \Psi(\mathcal G(u))=\Psi(u)+1,\quad
 \Psi(u)=u+\tfrac13\Log u+\Oh(u^{-1}),\quad
 \Phi(w)=-\Psi(2/w)+\tfrac13\log2.
\]
On nested proper sectors its inverse orbits satisfy
$|\mathcal G^{\circ j}(u)|\ge c(|u|+j)$, and the normalized orbit limit is available on these sectors. These sector estimates, rather than a solely real-axis expansion, are used below.

The same report proves the positive Laplace representation
\begin{equation}\label{ied:log:eq:density}
 P(s)=\int_0^\infty\e^{st}h(t)\dd t,
 \qquad h\in C^\infty(0,\infty),\qquad h(t)>0\ (t>0).
\end{equation}
The density has every exponential moment. Fourier inversion applies to $P'$ on vertical lines and to the fixed higher derivatives used below. We do not import endpoint smoothness or an endpoint expansion for $h$.

\begin{writenote}[Where Part~I proves this, 6~October 2026]
``Report 225, Sections 4 and 6'' are Part~I's Sections~\ref{ied:lead:sec:fatou} and~\ref{ied:lead:sec:measure}: Lemma~\ref{ied:lead:lem:fatou} (whose proof constructs $G=\mathcal G$, $\Psi$ and $\Phi=-\Psi(2/w)+\frac13\log2$, with the orbit bound $|G^j(u)|\ge c(|u|+j)$), Proposition~\ref{ied:lead:prop:measure} and Corollary~\ref{ied:lead:cor:smooth}; Fourier inversion of $P'$ is \eqref{ied:lead:eq:limitcoeff}.
\end{writenote}

\subsection{The absolute derivative estimates before division}
For large $m$, set
\begin{align}\label{ied:log:eq:scaling}
 d_m&=-\tfrac13\log m+\tfrac13\log T-B,
 &r_m&=\frac{T}{m+d_m},\notag\\
 G_m(s)&=F_m(r_m\e^{s/m}),
 &Q_m(s)&=(d_m+\tfrac13)s-\tfrac12s^2-\tfrac13d_m-\kappa T.
\end{align}
Here $\kappa$ is the fixed real constant of Report 226. Its value is immaterial for the transfer in this report; for a fully specified interface it is
\[
 \kappa=\frac{\pi^3}{96}+\frac{\pi^2}{12}-\frac{13\pi}{48}
 -\frac{11}{36}-\frac{\mathcal C}{6}-\frac{\pi\log2}{24},
 \qquad \mathcal C=\sum_{j\ge0}\frac{(-1)^j}{(2j+1)^2}.
\]
Report 226, Sections ``Derivative extraction and classical Fourier inversion'' and ``Uniform proportional height correction,'' proves the following absolute estimates, all depending on height $m$, not on the target degree. With $M=m^{1/20}$, the thirtieth derivative of
\[
 G_m(s)-P(s)-m^{-1}Q_m(s)P'(s)
\]
is $\Oh(m^{-3/2})$ on $s=\ii u$, $|u|\le M$, by analytic disk estimates on an enlarged strip. The finite-height intermediate derivative is $\Oh(|u|^{-31})$, and the outer derivative is $\Oh(m^{-31})$ on an interval of length $\Oh(m)$. Their absolute tails are
\begin{equation}\label{ied:log:eq:tails}
 \Oh(M^{-30}+m^{-30})=\Oh(m^{-3/2}).
\end{equation}
The limiting derivative has $P^{(30)}(\ii u)=\Oh(|u|^{-31})$, and
\begin{equation}\label{ied:log:eq:Qtail}
 (Q_mP')^{(30)}(\ii u)
   =\Oh(|u|^{-30}+\log m\,|u|^{-31})
\end{equation}
for large $|u|$. No degree-to-height division has yet occurred.

Write $\widetilde G_{m,30}(u)$ for $G_m^{(30)}(\ii u)$ on
$[-\pi m,\pi m]$ and zero elsewhere. This notation means extension of the derivative's values, not differentiation of a zero extension. Taking absolute values before integrating the preceding bounds proves
\begin{equation}\label{ied:log:eq:L1input}
 \left\|\widetilde G_{m,30}
 -P^{(30)}(\ii\,\cdot)
 -\frac{(Q_mP')^{(30)}(\ii\,\cdot)}m\right\|_{L^1(\R)}
       =\Oh(m^{-7/5}).
\end{equation}
Indeed the central error is $\Oh(Mm^{-3/2})=\Oh(m^{-29/20})$;
\eqref{ied:log:eq:tails} and the limiting $P$ tail are $\Oh(m^{-3/2})$; the correction tail after division by $m$ is
\[
 \Oh\bigl(m^{-1}(M^{-29}+\log m\,M^{-30})\bigr),
\]
which is smaller still. These estimates even give the exponent $29/20$ here; the weaker $7/5$ is convenient and already sufficient. Equation~\eqref{ied:log:eq:L1input} is an absolute $L^1$ consequence of the proved estimates, not an extrapolation of Report 226's compact-slope relative theorem.

\begin{writenote}[Where Part~II proves this, 6~October 2026]
Report~226's sections ``Derivative extraction and classical Fourier inversion'' and ``Uniform proportional height correction'' are Part~II's Sections~\ref{ied:corr:sec:extract} and~\ref{ied:corr:sec:height}: \eqref{ied:corr:eq:Gexpand}, \eqref{ied:corr:eq:highderivtail}, the outer-angle bound and the tail estimates following it, there at height $n$ and here at height $m$ (Part~II's proof of Theorem~\ref{ied:corr:thm:height} makes the same change). $d_m$ and $Q_m$ are Part~II's $d_m$ and $Q_n$ with $n$ replaced by $m$. The $L^1$ bound \eqref{ied:log:eq:L1input} itself is derived here.
\end{writenote}

\subsection{Classical background and attribution}\label{ied:log:sub:classical}
The endpoint calculation belongs to the universal unforced parabolic map $w\mapsto\e^w-1$. Finite sectorial Fatou expansions, their rational formal recursion, and Hankel inversion are classical methods. Dudko and Sauzin~\cite{DS} give general sectorial asymptotics, including nonzero resiter; the present reciprocal germ has resiter $1/3$. We give a direct finite-defect proof adapted to the precise normalization above.

The public iterated Bell report~\cite{IBD} already contains all-finite-order inverse-coordinate machinery and a related contour amplitude. Its inspected proportional-depth transfer fixes $n/m$ in a positive compact interval and explicitly does not assert endpoint uniformity. In its notation, renamed here to avoid collisions, the small-variable coordinate is
\[
 \Psi_{\mathrm{IBD}}(w)=-\Phi(w)+\tfrac13\log2.
\]
This local identity follows from the common normalized reciprocal orbit limit. A formal contour match suggests an amplitude relation
$I(\beta)=\beta\,2^{\beta/3}h(\beta)$, but that identification requires its own analytic contour correspondence and is not used here. We neither attribute an endpoint theorem to that source nor claim a new array-specific origin for the universal coefficient calculus.

Prellberg's 2002 seminar, as summarized by Mishna~\cite{Prellberg}, already presents the parabolic logarithmic mechanism and its $1/3$ specialization. Nagaev and Vakhtel~\cite{NV} give an important unforced antecedent: their critical branching large-deviation correction has the form
$\exp[-\gamma u\log u/m]$, with $\gamma=1/3$ in the Poisson case, in an appropriate moderate-deviation range. The heuristic high-moment saddle $u\asymp n$ places the logarithmic crossover in that range when $m/n\to\infty$. Passing from those probability estimates to a full moment equivalent needs tail domination, which is not supplied by this comparison; they do not furnish the forced Euler coefficient transfer proved here. The new scope relative to Report 225 is the quantitative transfer as $n/m$ decreases to logarithmic scales, using the forced-array contour estimates and its constants $T,B$. No worldwide priority claim is made for the parabolic mechanism or universal endpoint calculus.

\begin{writenote}[The iterated Bell report, 6~October 2026]
\cite{IBD} is \file{a139383-iterated-bell-diagonals} in this repository; its article is unchanged since the pin, so the sections named are its current Sections~2, 4 (Part~I there) and 12, 13, 17 (Part~II there). Its $\Psi$ is $\Psi_{\mathrm{IBD}}$ here and its $I(\beta)$ is the amplitude of its Theorem \texttt{ibd:pd:thm:main}. The relation $I(\beta)=\beta2^{\beta/3}h(\beta)$ is unproved; at $\beta=1$ the two sides agree to eight significant digits (Remark~\ref{ied:rem:ibdnumbers}), and Section~\ref{ied:sec:further}, item~\ref{ied:q:ibd}, records it as a question shared with that report, with a conditional consequence for its open endpoint (Remark~\ref{ied:rem:labelled}). That report's own literature section places the Nagaev--Vakhtel range outside its compact-slope regime $m\asymp n$; here $m/n\to\infty$, which is the range named above.
\end{writenote}

\section{Finite Fatou expansions with sectorial remainders}\label{ied:log:sec:fatou}
\begin{lemma}\label{ied:log:lem:finitefatou}
For every integer $N\ge0$ there are unique rational numbers $a_1,\ldots,a_N$ such that
\begin{equation}\label{ied:log:eq:fatoufinite}
 \Phi(w)=-\frac2w+\frac13\Log w+\sum_{j=1}^N a_jw^j
                    +\Oh(w^{N+1})
\end{equation}
uniformly on each sufficiently small closed proper sector. For every fixed derivative order $k$, the $k$th derivative of the remainder is $\Oh(|w|^{N+1-k})$ on smaller closed sectors. The first values are
\[
 a_1=-\frac1{36},\quad a_2=\frac1{540},\quad
 a_3=\frac1{7776},\quad a_4=-\frac{71}{435456}.
\]
\end{lemma}
\begin{proof}
Choose a finite reciprocal model
\[
 V_N(u)=u+\tfrac13\Log u+\sum_{j=1}^N v_ju^{-j}.
\]
The expansion of \eqref{ied:log:eq:Greciprocal} has rational coefficients. In
$V_N(\mathcal G(u))-V_N(u)-1$, the new coefficient $v_j$ first contributes
$-jv_ju^{-j-1}$. Cancellation is therefore triangular, unique and rational, and gives
\[
 D_N(u):=V_N(\mathcal G(u))-V_N(u)-1=\Oh(u^{-N-2}).
\]
Along an inverse orbit, $V_N(u)-u-(\Log u)/3\to0$, so the normalized orbit limit is the same $\Psi$ as in Section~\ref{ied:log:sec:interfaces}. Telescoping gives the sign-sensitive identity
\begin{equation}\label{ied:log:eq:defectsum}
 \Psi(u)-V_N(u)=\sum_{j\ge0}D_N(\mathcal G^{\circ j}(u)).
\end{equation}
The sign is positive: the finite sum equals
$V_N(\mathcal G^{\circ J}(u))-J-V_N(u)$ before taking $J\to\infty$.
The orbit bound gives absolute convergence and
\[
 \sum_{j\ge0}|D_N(\mathcal G^{\circ j}(u))|
 \le C\sum_{j\ge0}(|u|+j)^{-N-2}=\Oh(|u|^{-N-1}).
\]
Transforming by $\Phi(w)=-\Psi(2/w)+(\log2)/3$ proves
\eqref{ied:log:eq:fatoufinite}; thus $a_j=-v_j/2^j$. Cauchy's derivative estimates on disks of radius proportional to $|w|$ inside nested sectors give every fixed differentiated remainder. Direct triangular substitution gives the stated coefficients. No convergence of the infinite formal coordinate series is required.
\end{proof}

\section{The endpoint value by direct Fourier inversion}\label{ied:log:sec:endpointvalue}
The leading endpoint value follows already from the first-order interface in Report 225.
\begin{proposition}\label{ied:log:prop:hzero}
The density extends continuously to zero with
\begin{equation}\label{ied:log:eq:hzero}
 h(t)=2+\Oh(t\log(1/t)),\qquad h(0)=2.
\end{equation}
\end{proposition}
\begin{proof}
Fourier inversion of $P'$ on the line $\Rea s=-1/t$ gives
\[
 t h(t)=\frac{\e}{2\pi}\int_\R
 P'(-1/t+\ii u)\e^{-\ii ut}\dd u.
\]
The entire line is in the fixed inverse sector $|\arg(-s)|\le\pi/2$.
The leading term $2/s^2$ in \eqref{ied:log:eq:inverseidentity} inverts to $2t$. The absolute error is bounded by a constant times
\[
 \int_\R
 \frac{\log\sqrt{t^{-2}+u^2}}{(t^{-2}+u^2)^{3/2}}\dd u
 =\Oh(t^2\log(1/t)),
\]
by $u=v/t$. Division by $t$ proves the assertion.
\end{proof}

\begin{writenote}[Answers Parts~I and~II, 6~October 2026]
This answers, at the endpoint zero, the remark closing the proof of Part~I's Proposition~\ref{ied:lead:prop:measure} and the questions of Subsections~\ref{ied:lead:sub:heights} and~\ref{ied:corr:sub:endpoint} that ask for an expansion of $h$ at zero.
\end{writenote}

\section{An exact Hankel representation}\label{ied:log:sec:hankel}
Set
\begin{equation}\label{ied:log:eq:Acoord}
 \mathcal A(w)=\Phi(w)+2/w-\tfrac13\Log w=\Oh(w).
\end{equation}
Fix $\pi/2<\theta<\pi$ and $a>0$. The contour $H=H_{\theta,a}$ starts at infinity on argument $-\theta$, travels inward to $a\e^{-\ii\theta}$, follows the circle $|v|=a$ counterclockwise through the positive axis to $a\e^{\ii\theta}$, and leaves along argument $\theta$. All logarithms in the contour formula use the principal branch.

\begin{theorem}\label{ied:log:thm:hankel}
For sufficiently small positive $t$,
\begin{equation}\label{ied:log:eq:hankel}
 h(t)=\frac{2(2t)^{-t/3}}{2\pi\ii}
 \int_H\e^v v^{t/3-2}\exp[-t\mathcal A(2t/v)]\dd v.
\end{equation}
The integral is absolutely convergent and is an exact identity. It supplies analytic continuation to a sufficiently small punctured sector about positive $t$.
\end{theorem}
\begin{proof}
Let $W(x)=P(-x)$. With an upward vertical contour, the absolutely convergent inverse of $P'$ is
\begin{equation}\label{ied:log:eq:bromwich}
 t h(t)=-\frac1{2\pi\ii}
 \int_{c-\ii\infty}^{c+\ii\infty}W'(x)\e^{tx}\dd x,
 \qquad c>0.
\end{equation}
To check the sign, replacing $s$ by $-x$ reverses the vertical orientation and $P'(-x)=-W'(x)$. Equivalently, change $u$ to $-u$ in the absolutely convergent Fourier formula. We take $c=1/t$.

Choose $\theta<\theta_1<\pi$. For small $t$, the line and scaled contour $x=v/t$ lie beyond the inverse threshold in the sector $|\arg x|\le\theta_1$, apart from no problematic compact portion since $|v|\ge a$. Truncate both at common large $x$-radius $R$. The circular connectors from the vertical endpoints to the rays lie to the left of the vertical line, so $|\e^{tx}|\le\e$ there. The estimate $W'(x)=\Oh(R^{-2})$ and connector length $\Oh(R)$ give integral $\Oh(R^{-1})$. Thus these connectors vanish as $R\to\infty$. The vertical integral converges absolutely by $W'=\Oh(|x|^{-2})$, and the ray integrals by $\cos\theta<0$. Cauchy's theorem deforms the upward line to the stated lower-inward, upper-outward Hankel orientation.

For the second deformation put
\begin{align*}
 D_t(v)&=\tfrac13(\Log v-\log(2t))-\mathcal A(2t/v),\\
 x_\ell(v)&=v/t+\ell D_t(v),\qquad 0\le\ell\le1.
\end{align*}
Uniformly for $|v|\ge a$ on $H$,
\[
 |D_t(v)|\le C(1+|\log|v||+|\log t|),\qquad
 \sup_{v\in H}|tD_t(v)/v|\longrightarrow0.
\]
Consequently this homotopy stays in the larger proper inverse sector and satisfies $|x_\ell(v)|\ge|v|/(2t)$ for small $t$. Its endpoint at $\ell=1$ is exactly
\begin{equation}\label{ied:log:eq:xv}
 x(v)=-\Phi(2t/v)
 =v/t+\tfrac13(\Log v-\log(2t))-\mathcal A(2t/v).
\end{equation}
The compatible principal branches are valid since $\arg(2t/v)=-\arg v$ stays strictly inside $(-\pi,\pi)$. In particular no negative-real branch cut is crossed.

At ray endpoints $v=R\e^{\pm\ii\theta}$ the homotopy connectors have length $\Oh(\log R+|\log t|)$ and
\[
 \Rea(t x_\ell(v))=R\cos\theta+\Oh(t\log R+t|\log t|).
\]
Their integrals tend to zero exponentially as $R\to\infty$, for fixed small $t$. On each finite truncated homotopy the integrand $W'(x)\e^{tx}$ is entire, so Cauchy's theorem applies. This is a deformation of paths; it does not require global univalence of the map $v\mapsto x(v)$.

By the imported inverse identification,
\[
 W(x(v))=P(\Phi(2t/v))=2t/v,
 \qquad W'(x(v))\dd x=-2t v^{-2}\dd v.
\]
Substitution into \eqref{ied:log:eq:bromwich}, retaining its minus sign, gives
\eqref{ied:log:eq:hankel}, including the factor $2$. As an orientation check,
$(2\pi\ii)^{-1}\int_H\e^vv^{-2}\dd v=1$, which gives the positive endpoint value.

Finally choose $\eta>0$ with $\eta+\theta<\pi$. For complex $t$ with
$0<|t|<\epsilon$ and $|\arg t|<\eta$, every $2t/v$ stays in a fixed coordinate sector. Use the corresponding branch of $\Log t$ in \eqref{ied:log:eq:hankel}. Uniform convergence on compact subsets follows because $\e^{\Rea v}$ on the rays dominates all fixed polynomial and logarithmic factors. The resulting holomorphic function agrees with the real density for positive $t$ by the identity already proved.
\end{proof}

\begin{writenote}[A consequence, 6~October 2026]
Since the continuation is holomorphic on a punctured sector containing $(0,\epsilon)$ and agrees with $h$ there, $h$ is real analytic on $(0,\epsilon)$. Real analyticity on all of $(0,\infty)$, asked in Part~I's Subsection~\ref{ied:lead:sub:heights}, is not shown (Section~\ref{ied:sec:further}, item~\ref{ied:q:analytic}).
\end{writenote}

\section{A smooth endpoint germ and the logarithmic transseries}\label{ied:log:sec:germ}
Here ``endpoint transseries'' refers only to finite asymptotic expansions in powers $t^j(\log t)^r$ and their differentiated remainders. No summability, exponentially small completion, or convergence assertion is attached to the term.

\begin{theorem}\label{ied:log:thm:germ}
Let $g(t)=t^{t/3}h(t)$ for $t>0$. There are real coefficients $b_j$, with $b_0=2$, such that for every fixed $N,k\ge0$,
\begin{equation}\label{ied:log:eq:germ}
 \frac{\dd^k}{\dd t^k}
 \left(g(t)-\sum_{j=0}^N b_jt^j\right)
       =\Oh(t^{N+1-k})\qquad(t\downarrow0).
\end{equation}
Consequently $g$ has a one-sided $C^\infty$ extension to $[0,\epsilon)$ and
$b_j=g^{(j)}(0)/j!$. This does not assert analyticity at zero.
\end{theorem}
\begin{proof}
Cancel the explicit $t^{-t/3}$ in \eqref{ied:log:eq:hankel} to obtain
\begin{equation}\label{ied:log:eq:germcontour}
 g(t)=\frac{2\,2^{-t/3}}{2\pi\ii}
 \int_H\e^vv^{-2}
 \exp\left[\tfrac t3\Log v-t\mathcal A(2t/v)\right]\dd v.
\end{equation}
Use Lemma~\ref{ied:log:lem:finitefatou} to any required finite order. On $H$, $|v|\ge a$, and the remainder after retaining $\mathcal A$ through degree $q$ is uniformly
$\Oh(|t|^{q+1}|v|^{-q-1})$. Multiplication by $t$ improves its order once more. The remaining factors $2^{-t/3}$ and $v^{t/3}$ are entire in $t$. Their finite Taylor remainders are bounded by fixed powers of $1+|\Log v|$ times a small fixed power of $|v|$. Along the rays these are integrably dominated by $\e^{\Rea v}$; the inner arc is compact and avoids zero. Finite expansion and integration thus yield, uniformly in a smaller punctured $t$-sector,
\[
 g(t)=\sum_{j=0}^N b_jt^j+\Oh(|t|^{N+1}).
\]
Apply Cauchy's estimate to the actual holomorphic remainder on disks of radius $\delta t$ centered at positive $t$. This proves \eqref{ied:log:eq:germ}, rather than formally differentiating an unsupported pointwise expansion. For a prescribed derivative order choose $N\ge k$; its limit is $k!b_k$. Choosing one additional order establishes compatibility of these endpoint derivatives. The constant follows from Proposition~\ref{ied:log:prop:hzero}, or directly from the orientation check in Theorem~\ref{ied:log:thm:hankel}.
\end{proof}

The classical reciprocal-Gamma Hankel identity, on the orientation specified above, is
\[
 \frac1{2\pi\ii}\int_H\e^vv^{-z}\dd v=\frac1{\Gamma(z)}.
\]
For a finite nonnegative multi-index $\mathbf r=(r_1,r_2,\ldots)$, let
$R_{\mathbf r}=\sum r_j$ and $J_{\mathbf r}=\sum jr_j$. Expanding the last exponential in \eqref{ied:log:eq:germcontour} gives the finite-order generator
\begin{equation}\label{ied:log:eq:generator}
 g(t)\sim2\,2^{-t/3}
 \sum_{\mathbf r}
 \frac{(-1)^{R_{\mathbf r}}2^{J_{\mathbf r}}}
      {\Gamma(2+J_{\mathbf r}-t/3)}
 \left(\prod_j\frac{a_j^{r_j}}{r_j!}\right)
 t^{R_{\mathbf r}+J_{\mathbf r}}.
\end{equation}
At any specified degree only finitely many multi-indices occur. Equation~\eqref{ied:log:eq:generator} is interpreted by truncation; Theorem~\ref{ied:log:thm:germ} proves its remainder, and no convergence of the full multi-index sum is used.

\begin{corollary}\label{ied:log:cor:densityseries}
Writing $X=\log t$, for every fixed $N\ge0$,
\begin{equation}\label{ied:log:eq:densityseries}
 h(t)=\sum_{j=0}^N t^jH_j(X)
 +\Oh\bigl(t^{N+1}(1+|X|)^{N+1}\bigr),\quad
 H_j(X)=\sum_{r=0}^j b_{j-r}\frac{(-X/3)^r}{r!}.
\end{equation}
The $k$th derivative of its remainder is
$\Oh(t^{N+1-k}(1+|\log t|)^{N+1})$ for every fixed $k$.
In particular $H_j$ has degree exactly $j$, with leading coefficient
$2(-1/3)^j/j!$.
\end{corollary}
\begin{proof}
Use $h(t)=\exp[-t\Log t/3]g(t)$ in the punctured sector and multiply finite expansions. The exponential remainder obeys the stated power-log bound. Cauchy estimates on the same proportional disks give the differentiated bound; equivalently one may differentiate the finite product with its controlled remainder. Since $b_0=2$, its highest logarithmic coefficient is the stated nonzero number.
\end{proof}
In particular,
\begin{align}\label{ied:log:eq:hfirst}
 h(t)={}&2+\frac23t(C_0-\log t)\notag\\
 &+\frac{t^2}{9}\left((\log t-C_0)^2+\frac32-\frac{\pi^2}{6}\right)
 +\Oh\bigl(t^3(1+|\log t|)^3\bigr).
\end{align}
The differentiated estimates imply
\[
 h'(t)=-\tfrac23(\log t+\gamma+\log2)+\Oh(t\log^2t),
 \qquad h''(t)=-\frac2{3t}+\Oh(\log^2t).
\]
Thus the original density $h$ is not $C^1$ at zero, although the explicitly renormalized germ $g$ is smooth.

\section{The rational and zeta coefficient structure}\label{ied:log:sec:coefficients}
Because $g(t)=2+\Oh(t)$ on a sector, a branch of $\log(g(t)/2)$ exists near the endpoint, approaching zero. Theorem~\ref{ied:log:thm:germ} gives ordinary finite Taylor remainders for this logarithm and each fixed derivative. Therefore
\begin{equation}\label{ied:log:eq:logdensity}
 \log\frac{h(t)}2=-\frac t3\log t+\frac{C_0}3t
          +\sum_{j=2}^N c_jt^j+\Oh(t^{N+1}).
\end{equation}
After removing the single explicit $t\log t$ term, every $k$th derivative of the remainder is $\Oh(t^{N+1-k})$.

\begin{proposition}\label{ied:log:prop:rhozeta}
The coefficients in \eqref{ied:log:eq:logdensity} have the form \eqref{ied:log:eq:cs}. Their first rational parts are
\begin{align*}
 \rho_2&=\frac1{12},& \rho_3&=\frac{17}{1080},&
 \rho_4&=\frac{13}{3888},\\
 \rho_5&=\frac{2071}{2449440},&
 \rho_6&=\frac{5131261}{22044960000},&
 \rho_7&=\frac{261791507}{3968092800000}.
\end{align*}
\end{proposition}
\begin{proof}
Factor $1/\Gamma(2-t/3)$ out of \eqref{ied:log:eq:generator}. The remaining formal ordinary series $S(t)$ has rational coefficients and constant term one, because each Gamma ratio is exactly
\[
 \frac{\Gamma(2-t/3)}{\Gamma(2+J-t/3)}
       =\prod_{k=2}^{J+1}(k-t/3)^{-1},
\]
with empty product one when $J=0$. Thus
\begin{equation}\label{ied:log:eq:factorS}
 \frac{g(t)}2\sim \frac{2^{-t/3}}{\Gamma(2-t/3)}S(t),
 \qquad \log S(t)\in t^2\mathbb Q[[t]].
\end{equation}
At each finite order this is justified by the already proved expansion. The familiar local Gamma expansion yields
\[
 -\log\Gamma(2-t/3)=\frac{1-\gamma}3t
 +\sum_{j\ge2}\frac{1-\zeta(j)}{j3^j}t^j.
\]
Combining it with \eqref{ied:log:eq:factorS} proves the rational/zeta form and the linear coefficient $C_0/3$. For example,
\[
 S(t)=1+\frac{t^2}{36}+\frac{11t^3}{3240}+\Oh(t^4),
\]
which immediately gives $c_2$ and $c_3$. Finite multi-index enumeration with the rational $a_j$ gives the further displayed $\rho_j$. The archive supplies exact-arithmetic reproduction of these identities, including a separate finite reversion check. The analytic remainders follow from the preceding proofs, not from symbolic checking.
\end{proof}

\section{Transfer to shrinking degree to height ratios}\label{ied:log:sec:transfer}
\begin{theorem}\label{ied:log:thm:transfer}
Uniformly over positive integer degrees $n$ as $m\to\infty$, writing $\beta=n/m$, one has the absolute estimate
\begin{equation}\label{ied:log:eq:absolute}
 m A(n,m)r_m^n=h(\beta)
                 +\Oh\left(\frac{\log m}{m\beta^{30}}\right).
\end{equation}
For every fixed $B_0>0$, uniformly for $0<\beta\le B_0$,
\begin{align}\label{ied:log:eq:relative}
 A(n,m)&=\frac{h(\beta)}m(m/T)^n\e^{\beta d_m}
           \{1+\Oh(E(m,\beta))\},\\
 E(m,\beta)&=\frac{\log m}{m\beta^{30}}
                   +\frac{\beta\log^2m}{m}.\notag
\end{align}
The relative expression is an asymptotic equivalent whenever its stated error tends to zero. In particular this holds uniformly on every fixed logarithmic band
\begin{equation}\label{ied:log:eq:logband}
 (\log m)^{-D}\le\beta\le B_0,\qquad D>0\ \text{fixed}.
\end{equation}
\end{theorem}
\begin{proof}
Cauchy's estimates for the inverse-sector expansion of $P$ give
$P^{(q)}(\ii u)=\Oh(|u|^{-q-1})$ at infinity for every fixed $q\ge1$.
Leibniz' rule, the quadratic form of $Q_m$, and $d_m=\Oh(\log m)$ imply
\[
 \|(Q_mP')^{(30)}(\ii\,\cdot)\|_1=\Oh(\log m).
\]
The bounded interval uses analyticity; on the tails \eqref{ied:log:eq:Qtail} is integrable. Hence the pre-division estimate \eqref{ied:log:eq:L1input} implies
\begin{equation}\label{ied:log:eq:L1drop}
 \|\widetilde G_{m,30}-P^{(30)}(\ii\,\cdot)\|_1
                       =\Oh(\log m/m).
\end{equation}
For every positive integer $n$, termwise full-period extraction is exact:
\begin{equation}\label{ied:log:eq:extract}
 m\beta^{30}A(n,m)r_m^n
 =\frac1{2\pi}\int_{-\pi m}^{\pi m}
 G_m^{(30)}(\ii u)\e^{-\ii\beta u}\dd u.
\end{equation}
Indeed a monomial of degree $k$ acquires $(k/m)^{30}$; orthogonality over a period selects $k=n$, leaving $\beta^{30}$. The contour is inside the disk of analyticity for large $m$. No integration by parts or endpoint matching assumption is needed.

Fourier inversion of $P^{(30)}$ gives $\beta^{30}h(\beta)$ for each $\beta>0$. The oscillatory factor has absolute value one, so \eqref{ied:log:eq:L1drop} and division by $\beta^{30}$ give \eqref{ied:log:eq:absolute}. The constant before that division is independent of the target degree. The displayed $\beta^{-30}$ loss is therefore genuine and has not been discarded. Only this absolute assertion is globally uniform in $\beta$.

By Proposition~\ref{ied:log:prop:hzero}, strict positivity, and continuity, $\min_{[0,B_0]}h>0$. Thus division by $h(\beta)$ is uniform on this interval. Furthermore
\begin{align*}
 r_m^{-n}&=(m/T)^n\exp\{n\log(1+d_m/m)\}\\
 &=(m/T)^n\e^{\beta d_m}
 \left\{1+\Oh\left(\frac{\beta\log^2m}{m}\right)\right\}
\end{align*}
uniformly for $\beta\le B_0$. This proves \eqref{ied:log:eq:relative}. On \eqref{ied:log:eq:logband},
\[
 E(m,\beta)=\Oh\left(\frac{(\log m)^{30D+1}}m
                      +\frac{\log^2m}{m}\right)=o(1).
\]
The constants may depend on $B_0,D$, and the fixed derivative order. There is no uniform relative theorem for arbitrary $\beta\downarrow0$ here.
\end{proof}

\begin{writenote}[Relation to Parts~I and~II, 6~October 2026]
On a fixed compact $\beta$-range \eqref{ied:log:eq:relative} is Part~I's \eqref{ied:lead:eq:proportional} with the rate $\Oh(\log m/m+\log^2m/m)$ in place of $o(1)$; Part~II's Theorem~\ref{ied:corr:thm:height} is sharper there. The new content is the uniformity down to $(\log m)^{-D}$.
\end{writenote}

\section{Exact logarithmic cancellation and the crossover proof}\label{ied:log:sec:crossoverproof}
Combining \eqref{ied:log:eq:R} and \eqref{ied:log:eq:relative} gives
\[
 R(n,m)=\frac{h(\beta)}2\e^{\beta d_m}\{1+\Oh(E(m,\beta))\}.
\]
In a logarithmic band with sufficiently small upper endpoint, its logarithm is legitimate because $E=o(1)$ and $h>0$. Equation~\eqref{ied:log:eq:logdensity} then gives, for every fixed $N\ge1$,
\begin{equation}\label{ied:log:eq:logmaster}
 \log R(n,m)=-\frac\beta3\log n+K\beta
      +\sum_{j=2}^N c_j\beta^j
      +\Oh\bigl(\beta^{N+1}+E(m,\beta)\bigr).
\end{equation}
The identity behind this formula is exact:
\[
 -\frac\beta3\log\beta-\frac\beta3\log m
       =-\frac\beta3\log n,
 \qquad \log\beta+\log m=\log n.
\]
It removes potential $\log\log n$ terms before any asymptotic substitution.

\begin{proof}[Proof of Theorem~\ref{ied:log:thm:crossover}]
If $m/(n\log n)\to\lambda\in(0,\infty)$, then
$\beta\to0$, $\beta\log n\to1/\lambda$, and $\log m\sim\log n$. The ratio belongs, for example, to a fixed logarithmic band with $D=2$ for all sufficiently large $n$, and $E\to0$. Equation~\eqref{ied:log:eq:logmaster} gives \eqref{ied:log:eq:crossover}.

For the uniform fixed-order assertion let $m=\lfloor\lambda nL\rfloor$, $L=\log n$, and $\lambda\in\Lambda$. Then
\[
 \beta=\frac1{\lambda L}+\Oh_\Lambda\left(\frac1{nL^2}\right),
 \qquad E(m,\beta)=\Oh_\Lambda(L^{30}/n).
\]
Integer rounding changes $\beta L$ by $\Oh_\Lambda(1/(nL))$. These errors are smaller than every fixed inverse power of $L$. Substitute the finite expansion \eqref{ied:log:eq:logmaster}; its endpoint remainder becomes $\Oh_{\Lambda,N}(L^{-N-1})$. This proves \eqref{ied:log:eq:logcross} and, by ordinary finite exponentiation, \eqref{ied:log:eq:Rcross}. No estimate of arbitrarily high order in $m^{-1}$ was used.
\end{proof}

More generally, whenever $\beta\to0$ inside a fixed logarithmic band,
\begin{equation}\label{ied:log:eq:generalregime}
 R(n,m)=\exp\left[-\frac{n}{3m}\log n\right]
                \{1+\Oh(\beta+E(m,\beta))\}.
\end{equation}
Thus $R\to1$ if $m/(n\log n)\to\infty$ and $m/n$ grows at most as a fixed power of $\log n$. Also $R\to0$ if $m/n\to\infty$ while $m/(n\log n)\to0$; this range lies in a logarithmic band because $m/n=o(\log n)$. Fixed degree followed by height tending to infinity is a different limit and is not covered by this transfer.

\section{Kaneiwa normalization and its first correction}\label{ied:log:sec:kaneiwa}
Kaneiwa~\cite{Kaneiwa} proves that his $p(r;n)$ is a degree-$(n-1)$ polynomial in $r$ with leading coefficient
\begin{equation}\label{ied:log:eq:ell}
 \ell_n=\frac{E_{n-1}}{(n-1)!}
       =[x^{n-1}](\tan x+\sec x).
\end{equation}
His indexing is $p(r;n)=A(n,r+1)$, so $r=m-1$. His fixed-degree polynomial theorem alone provides no uniform error as $n$ and $r$ grow together.

The nearest pole of $\tan x+\sec x$ is $T=\pi/2$, with principal part $2/(T-x)$. The apparent singularity at $-T$ is removable; the next pole is $-3T$. Subtract the principal part and use Cauchy's coefficient estimate on any fixed circle of radius strictly between $T$ and $3T$. For every fixed $1<c<3$ this gives
\begin{equation}\label{ied:log:eq:ellestimate}
 \frac{\ell_n}{2T^{-n}}=1+\Oh(c^{-n}).
\end{equation}
If $n\to\infty$ and $\beta=n/m\to0$, then
\[
 (n-1)\log(1-1/m)=-\beta+\Oh(1/m+n/m^2)\longrightarrow0.
\]
Consequently the crossover limit is unchanged when \eqref{ied:log:eq:R} is replaced by
\begin{equation}\label{ied:log:eq:Rk}
 R_{\mathrm K}(n,m)=\frac{A(n,m)}{\ell_n(m-1)^{n-1}}.
\end{equation}
For higher corrections the height shift must be retained. Directly,
\begin{equation}\label{ied:log:eq:Kshift}
 \log R_{\mathrm K}(n,m)=\log R(n,m)+\beta
       +\Oh(1/m+n/m^2+c^{-n}).
\end{equation}
In the crossover window these residual errors are smaller than every fixed inverse power of $\log n$. Thus Theorem~\ref{ied:log:thm:crossover} and its exponential version hold for $R_{\mathrm K}$ with $K$ replaced everywhere by $K+1$, while every $c_j$ for $j\ge2$ is unchanged. In particular the second-order exponential coefficient becomes $((K+1)^2/2+c_2)/(\lambda^2L^2)$. If one instead normalizes by $\ell_n m^{n-1}$, $K$ stays unchanged.

\begin{writenote}[Kaneiwa in Part~I, 6~October 2026]
The indexing $p(r;n)=A(n,r+1)$ is Part~I's $p(r;n)=[z^n]F_{r+1}$ (Subsection~\ref{ied:lead:sub:prior}), where the same Propositions~1 and~2 of Kaneiwa are credited.
\end{writenote}

\section{Smooth suppression ratio inverses}\label{ied:log:sec:inverse}
Fix a compact interval $\mathcal R=[r_0,r_1]\subset(0,1)$ and write
\begin{equation}\label{ied:log:eq:inversevariables}
 \tau=-\log r\in[\tau_0,\tau_1]=[-\log r_1,-\log r_0],
 \qquad \chi_n=\frac{\log n}{3}-K.
\end{equation}
The suppression logarithm $\tau$ is distinct from the density argument.
Define the actual smooth germ
\begin{equation}\label{ied:log:eq:Cgerm}
 C(\beta)=\log\frac{g(\beta)}2-\frac{C_0}{3}\beta.
\end{equation}
It has $C(0)=C'(0)=0$ and finite Taylor expansions
$C(\beta)=\sum_{j=2}^Nc_j\beta^j+\Oh(\beta^{N+1})$, with all fixed differentiated remainders. This definition uses the actual density, not an assumed convergent formal sum.

\subsection{The unique small model root}
\begin{proposition}\label{ied:log:prop:inverseroot}
There is $\beta_0>0$ such that for every sufficiently large $n$ and every $r\in\mathcal R$ the equation
\begin{equation}\label{ied:log:eq:modelequation}
 \tau=\chi_n\beta_*-C(\beta_*)
\end{equation}
has exactly one solution $\beta_*(n,r)\in(0,\beta_0)$. Uniformly in $r$,
$\beta_*\sim\tau/\chi_n$.
\end{proposition}
\begin{proof}
Choose $\beta_0$ so that $g>0$ and $C$ is smooth on $[0,\beta_0]$.
For large $n$, the derivative of the right side is
$\chi_n-C'(\beta)\ge \chi_n/2$ throughout that interval. Its values at the endpoints are zero and a quantity tending to infinity. The intermediate value theorem and strict increase prove existence and uniqueness in this endpoint interval. With $C(\beta)=\Oh(\beta^2)$ and $\tau$ in a fixed positive compact interval, first $\beta_*=\Oh(\chi_n^{-1})$, and then
$\beta_* =\tau/\chi_n+\Oh(\chi_n^{-3})$. No uniqueness outside the stated interval is asserted.
\end{proof}

\subsection{Every fixed reversion order}
Put $x=1/\chi_n$ and $u=\beta_*/x$. Equation~\eqref{ied:log:eq:modelequation} becomes
\begin{equation}\label{ied:log:eq:ureversion}
 u=\tau+C(xu).
\end{equation}
Finite substitution determines every coefficient recursively. The first terms are
\begin{align}\label{ied:log:eq:betainverse}
 \beta_*={}&\frac\tau{\chi_n}+\frac{c_2\tau^2}{\chi_n^3}
 +\frac{c_3\tau^3}{\chi_n^4}
 +\frac{2c_2^2\tau^3+c_4\tau^4}{\chi_n^5}
 +\Oh_{\mathcal R}(\chi_n^{-6}),\\
 y_*:=\frac1{\beta_*}={}&\frac{\chi_n}\tau-\frac{c_2}{\chi_n}
 -\frac{c_3\tau}{\chi_n^2}
 -\frac{c_2^2\tau+c_4\tau^2}{\chi_n^3}
 +\Oh_{\mathcal R}(\chi_n^{-4}).\label{ied:log:eq:yinverse}
\end{align}
The absent $\chi_n^{-2}$ term in \eqref{ied:log:eq:betainverse} and absent constant term in \eqref{ied:log:eq:yinverse} are consequences of $C'(0)=0$. Since $c_2<0$, the first displayed correction to the height model is positive.

\begin{theorem}\label{ied:log:thm:finiteinverse}
For every fixed $J\ge0$ there are explicitly generated polynomials $d_j(\tau)$ such that
\begin{align}\label{ied:log:eq:yJ}
 y_J(n,r)&=\frac{\chi_n}\tau+\sum_{j=1}^Jd_j(\tau)\chi_n^{-j},\\
 y_*(n,r)-y_J(n,r)&=\Oh_{\mathcal R,J}(\chi_n^{-J-1}).\notag
\end{align}
Here $d_1=-c_2$, $d_2=-c_3\tau$, and
$d_3=-c_2^2\tau-c_4\tau^2$. The coefficient $d_J$ requires only
$c_2,\ldots,c_{J+1}$.
\end{theorem}
\begin{proof}
Take a finite formal truncation $U_M(\tau,x)$ of \eqref{ied:log:eq:ureversion} through degree $M$. Taylor's theorem for the smooth $C$ gives a residual
\[
 U_M-\tau-C(xU_M)=\Oh_{\mathcal R,M}(x^{M+1}).
\]
Both the true $u$ and this truncation stay in a common positive compact interval, uniformly in $\tau$. There
\[
 \frac\partial{\partial u}\{u-\tau-C(xu)\}
       =1-xC'(xu)\ge\tfrac12
\]
for small $x$. The real mean-value theorem bounds $u-U_M$ by twice the residual. Since $u$ stays bounded away from zero, finite reciprocal expansion then gives every claimed order for $y_*=1/(xu)$. Write $\beta_*=\tau xV$. The equation becomes $V=1+\sum_{j\ge2}c_j\tau^{j-1}x^jV^j$ at each finite order. Every nonconstant coefficient of $V$ is therefore a polynomial divisible by $\tau$. The same holds for $1/V-1$, so division by $\tau x$ introduces no negative powers of $\tau$ beyond the leading $1/(\tau x)$. Formal recursion shows that the coefficient of $x^j$ in $y_*$ first involves $c_{j+1}$, proving polynomiality and the stated dependency. In particular direct substitution gives \eqref{ied:log:eq:betainverse} and \eqref{ied:log:eq:yinverse} with the displayed signs.

Equivalently one can solve the truncated polynomial model
$\tau=\chi_n\beta-\sum_{j=2}^{J+1}c_j\beta^j$ in its unique small-root interval. Its omitted $C$ remainder is $\Oh(\beta^{J+2})$; division by the model derivative costs $\chi_n^{-1}$ and reciprocation costs $\chi_n^2$. Thus its height differs from $y_*$ by $\Oh(\chi_n^{-J-1})$, and it has the same $y_J$. Other polynomial roots are irrelevant. No convergence of an infinite reversion series is invoked.
\end{proof}

\subsection{Rounded height prescriptions and observed ratios}
\begin{theorem}\label{ied:log:thm:integerprescription}
Choose any integer $m_J(n,r)$ with
$|m_J-ny_J(n,r)|\le1$. Then, uniformly for $r\in\mathcal R$,
\begin{equation}\label{ied:log:eq:prescriptionerror}
 \log\frac{R(n,m_J)}r
   =\Oh_{\mathcal R,J}\left(\chi_n^{-J-2}+\frac{\chi_n^{30}}n\right).
\end{equation}
Rounding the exact smooth model $ny_*$ instead gives error
$\Oh_{\mathcal R}(\chi_n^{30}/n)$ in this logarithmic ratio.

Conversely, suppose integer heights have $m/n\asymp\log n$ and the actual observed ratio $r_n=R(n,m)$ lies in $\mathcal R$. Then
\begin{equation}\label{ied:log:eq:observedinverse}
 \frac mn=y_J(n,r_n)
 +\Oh_{\mathcal R,J}\left(\chi_n^{-J-1}+\frac{\chi_n^{31}}n\right),
\end{equation}
with constants also depending on the fixed comparison bounds in $m/n\asymp\log n$.
\end{theorem}
\begin{proof}
Using the exact smooth germ in \eqref{ied:log:eq:relative}, rather than a Taylor truncation, gives on this height window
\begin{equation}\label{ied:log:eq:exactmodelratio}
 \log R(n,m)=-\chi_n\beta+C(\beta)
                         +\Oh(\chi_n^{30}/n).
\end{equation}
Indeed $m\asymp n \chi_n$, $\beta\asymp \chi_n^{-1}$, and
$E(m,\beta)=\Oh(\chi_n^{30}/n)$. As a function of $y=m/n$, the smooth part is
\[
 F_n(y)=-\chi_n/y+C(1/y),\qquad
 F_n'(y)=\frac{\chi_n-C'(1/y)}{y^2}\asymp \chi_n^{-1}
\]
uniformly on any common crossover-height window. The error
$y_J-y_*=\Oh(\chi_n^{-J-1})$ therefore costs $\Oh(\chi_n^{-J-2})$ in the log ratio. Integer rounding costs $\Oh(1/(n\chi_n))$, which is absorbed by $\chi_n^{30}/n$. Since $F_n(y_*)=-\tau=\log r$, this proves \eqref{ied:log:eq:prescriptionerror} and the exact-model statement.

For the converse, take $\tau_n=-\log r_n$ at the actual observed ratio. Equation~\eqref{ied:log:eq:exactmodelratio} and the lower derivative bound show that $m/n-y_*(n,r_n)=\Oh(\chi_n^{31}/n)$. Add Theorem~\ref{ied:log:thm:finiteinverse} to obtain \eqref{ied:log:eq:observedinverse}. Inverting the ratio error costs precisely one extra factor $\chi_n$.
\end{proof}
These formulas use a specified target $r$, or the actual $r_n$ at each $n$. Replacing $r_n$ by its mere limiting value generally loses higher-order information unless an appropriate convergence rate is known separately. The available coefficient theorem gives no discrete monotonicity, exact ratio solvability, least integer height, or effective onset. Thus the estimates are smooth prescriptions and asymptotic observed-ratio inverses, not threshold or exact-integer-recovery theorems.

For the Kaneiwa normalization \eqref{ied:log:eq:Rk}, replace $K$ by $K+1$ throughout this section, hence replace $\chi_n$ by $(\log n)/3-K-1$. The residual shift in \eqref{ied:log:eq:Kshift} is smaller than $\chi_n^{30}/n$ on the window. The coefficients $c_j$, reversion signs, and stated orders remain unchanged.

\begin{writenote}[No Lambert function here, 6~October 2026]
Unlike Parts~I and~II, these inverses involve no Lambert~$W$: the core of \eqref{ied:log:eq:modelequation} is linear (the case $b=0$ of \texttt{p0:thm:lambert-core} in \cite{TSvol}), and the reversion with the smooth germ $C$ is proved by Taylor's theorem; it is an analogue, not an instance, of \texttt{p0:thm:perturbed-inversion}, whose hypotheses ask for analyticity, and the rounded prescriptions are not instances of \texttt{p0:thm:staircase}. The write rechecked \eqref{ied:log:eq:betainverse}--\eqref{ied:log:eq:yinverse} by substitution.
\end{writenote}

\section{Exact finite comparisons and reproducibility}\label{ied:log:sec:computation}
\subsection{An exact finite coefficient algorithm}
Logarithmic differentiation of the Euler product gives, with the preceding-height coefficients $f_d=A(d,m)$,
\begin{equation}\label{ied:log:eq:integerrecurrence}
 b_j=\sum_{d\mid j}d f_d,\qquad c_0=1,\qquad
 c_k=\frac1k\sum_{j=1}^kb_jc_{k-j},\qquad
 A(k,m+1)=c_k\quad(k\ge1).
\end{equation}
Truncation at degree $N$ is exact for all coefficients retained. The public implementation uses integers and checks divisibility at each division. A separate finite-product implementation checks small tables using
$(1-z^d)^{-a}=\sum_{k\ge0}\binom{a+k-1}{k}z^{dk}$.
The working limits, degree and height at most $640$, are explicit implementation caps, not theorem restrictions. Computing height $M$ through degree $N$ uses $\Oh(MN^2)$ integer operations and $\Oh(N)$ live coefficients apart from saved outputs; bit costs grow with coefficient size.

The code uses explicit exceptions for all pass/fail conditions, so checks remain active under Python's \texttt{-O}. Chunked integer decimal conversion works with the minimum configurable integer-string digit limit $640$ and does not change any process-wide limit. A guard test exercises a value exceeding that digit count. The exact coefficient generation is independent of any decimal normalization or asymptotic theorem.

\subsection{Finite crossover diagnostics}\label{ied:log:sub:finitediag}
Table~\ref{ied:log:tab:finite} uses exact recurrence values at
$n=20,40,80$ and $m=\lfloor\lambda n\log n\rfloor$ for
$\lambda=1/2,1,3/2$. All heights are below $640$. For finite-height comparison define
\begin{equation}\label{ied:log:eq:finiteapprox}
 R_2(n,\lambda)=\e^{-1/(3\lambda)}
 \left(1+\frac{K}{\lambda\log n}
 +\frac{K^2/2+c_2}{(\lambda\log n)^2}\right).
\end{equation}
This is the explicit second-order ratio expansion, rather than the exponential of a truncated logarithm. The exact coefficient uses the displayed integer height. The nominal floor choices agree at 80 and 160 decimal digits; this is a computational check, not an interval certificate for the transcendental floor. Values of $R$ and $R_2$ are rounded floating evaluations, not interval enclosures; only $A(n,m)$ is computed exactly. The column $\e^{-1/(3\lambda)}$ is the limiting ratio. These finite examples illustrate the size of the correction and do not establish an effective error bound or a certified asymptotic onset.
\begin{table}[htbp]
\centering\small
% [write] The text of crossover_table.tex (shipped as data/228-logheight-results-crossover_table.tex), inlined:
% Generated by code/reproduce.py; exact counts, unvalidated decimal diagnostics.
% R2 is the explicit lambda-based second-order ratio expansion, not exp of a truncated logarithm.
\begin{tabular}{rrrrrr}
\hline
$n$ & $\lambda$ & $m$ & $R(n,m)$ & $e^{-1/(3\lambda)}$ & $R_2(n,\lambda)$ \\
\hline
20 & 0.5 & 29 & 0.89828117 & 0.51341712 & 0.86765860 \\
20 & 1 & 59 & 0.94696428 & 0.71653131 & 0.93775982 \\
20 & 1.5 & 89 & 0.96417635 & 0.80073740 & 0.95910776 \\
40 & 0.5 & 73 & 0.80413195 & 0.51341712 & 0.78974074 \\
40 & 1 & 147 & 0.89724904 & 0.71653131 & 0.89222870 \\
40 & 1.5 & 221 & 0.93039388 & 0.80073740 & 0.92738192 \\
80 & 0.5 & 175 & 0.74764498 & 0.51341712 & 0.73949701 \\
80 & 1 & 350 & 0.86485251 & 0.71653131 & 0.86215689 \\
80 & 1.5 & 525 & 0.90778645 & 0.80073740 & 0.90621689 \\
\hline
\end{tabular}
\caption{Unvalidated floating comparisons built from exact finite coefficients. The two displayed approximations have different purposes: the limit records the crossover theorem, and $R_2$ retains the first two crossover ratio corrections.}\label{ied:log:tab:finite}
\end{table}
The accompanying JSON records exact decimal coefficients, the chosen integer heights, normalization formulas, and diagnostic precision. The coefficient generator separately verifies the rational values through $\rho_7$ and the finite inverse identities. Symbolic and exact-integer checks are reproducible algebra; floating agreement is not an analytic remainder certificate.

\begin{writenote}[The labelled counterpart, 6~October 2026]
Remark~\ref{ied:rem:labelled} gives the same comparison for the labelled array $H(n,m)$ at the same heights, against a conditional prediction.
\end{writenote}

\subsection{Reproducing the archive}
The archive includes editable TeX, exact and symbolic programs, generated comparisons, the source ledger, and unchanged Reports 225 and 226. Its builder verifies a source manifest, reruns the public checks, copies inputs to a new external output directory, creates a private pdf\LaTeX\ format, compiles without shell escape, and rejects unresolved references, missing characters, or overfull boxes. Sources are preserved. Sorted ZIP members have fixed timestamps and permissions. Normal and optimized runs, including actual extraction and rebuilding of the ZIP, are compared by complete archive bytes and member hashes. Byte identity is a same-toolchain claim; different TeX, font, Python, or numerical-library versions can change output bytes.

\begin{writenote}[The companion here, 6~October 2026]
Report~228's companion is shipped as \file{code/228-logheight-*.py}, \file{data/228-logheight-*}, \file{228-logheight-code-README.md} and \file{228-logheight-sources-SOURCE_LEDGER.md}; its \file{supporting/} copies are Parts~I and~II. The delivered \file{code/reproduce.py --check} regenerated both result files byte for byte on a copy at the write; the report README gives the rerun route.
\end{writenote}

\section{Further questions and limits}\label{ied:log:sec:questions}
\subsection{Polynomially small ratios}\label{ied:log:sub:polysmall}
The explicit $\beta^{-30}$ loss identifies an obstacle to much smaller ratios. The present theorem is organized around fixed logarithmic bands and the height window $m\asymp n\log n$. A broad transfer theory for polynomially small $\beta$, especially one approaching the fixed-degree polynomial regime, is not proved here. A narrow range may be tested against the displayed absolute error, but that observation alone does not supply a sharp uniform transition theorem, optimal range, or a high-height error matching Kaneiwa's polynomial expansion. A useful next step is to exploit lower derivative orders or the retained first correction in the absolute Fourier norm, with explicitly balanced endpoint errors. Such a result would need its own statement and proof rather than a claim of arbitrary-small-$\beta$ uniformity.

\begin{writenote}[Still open, 6~October 2026]
Section~\ref{ied:sec:further}, item~\ref{ied:q:polysmall}.
\end{writenote}

\subsection{Effective bounds and validated evaluation}\label{ied:log:sub:effective}
Our remainder constants and onset are qualitative. A numerical certificate would require explicit bounds for the sectorial coordinate remainder, its inverse identification, contour tails, the finite-height derivative errors, and floating evaluation of the chosen model. Higher numerical precision or repeated mesh agreement cannot replace these bounds. An effective counterpart of \eqref{ied:log:eq:prescriptionerror} could produce rigorously enclosed height intervals for target ratios. Determining a least height would additionally require suitable discrete information about $R(n,m)$; no monotonicity result is available in this report.

\begin{writenote}[Still open, 6~October 2026]
Section~\ref{ied:sec:further}, items~\ref{ied:q:validated} and~\ref{ied:q:effective}.
\end{writenote}

\subsection{Stronger endpoint structure}\label{ied:log:sub:endpoint}
The smooth germ $g$ has an all-fixed-order Taylor asymptotic expansion on a sector, but analyticity at zero and convergence of its Taylor series have not been shown. Classical resurgence theory suggests studying sharper coefficient growth and summability questions, with attention to the nonzero resiter. These would concern the universal map, not a new Euler-specific density. Likewise a complete analytic comparison with the public IBD contour amplitude would clarify the global contour normalization beyond the local coordinate identity recorded above.

\begin{writenote}[Still open; shared with the iterated Bell report, 6~October 2026]
Section~\ref{ied:sec:further}, items~\ref{ied:q:analytic} and~\ref{ied:q:ibd}.
\end{writenote}

\subsection{Other arrays and other height windows}\label{ied:log:sub:other}
The constants $T$ and $B$ arise from the subsidiary Euler forcing. For another multiset recurrence, or another initial function, the universal endpoint calculation may persist while the clock and transfer change. Extending this method requires proving the new forced entrance and angular-tail estimates. At the opposite endpoint $\beta\to\infty$, a different density regime is needed. Finally, every fixed inverse-logarithm order here coexists with only one polynomial-in-height transfer rate; an all-orders $1/m$ theorem would require a further hierarchy of quantitative forced corrections and is not a consequence of the present proof.

\begin{writenote}[Still open, 6~October 2026]
Section~\ref{ied:sec:further}, items~\ref{ied:q:higher}, \ref{ied:q:endpoint} and~\ref{ied:q:other}.
\end{writenote}

\clearpage
\section{Further questions and research (all three Parts)}\label{ied:sec:further}
\begin{writenote}[Added in the write, 6~October 2026]
This section gathers the open questions of the three Parts under Vladimir's standing rule of 4~October 2026: every claim that a manuscript states without proof, every limitation it records as a question, and the question shared with the iterated Bell report. Each item names its sources, the sketch they give, and what is missing. Questions answered inside the merge are marked at the questions themselves. No claim of the three manuscripts was found to be wrong, and none is refuted here.
\end{writenote}
\begin{enumerate}[label=\arabic*.,ref=\arabic*,leftmargin=*,itemsep=4pt]
\item\label{ied:q:validated} \emph{Validated constants.} Sources: Part~I, Subsection~\ref{ied:lead:sub:validated}; Part~II, Subsection~\ref{ied:corr:sub:effective}; Part~III, Subsection~\ref{ied:log:sub:effective}. The amplitude $c$ has the absolutely convergent formula \eqref{ied:lead:eq:constantintegral}; the first correction needs $h(1)$, $h'(1)$, $h''(1)$, since $\kappa$ is now in closed form \eqref{ied:corr:eq:kappa}. Sketch: interval enclosures of the truncated Fatou coordinate on a stated sector (Lemma~\ref{ied:log:lem:finitefatou} gives the rational coefficients to every order), of the inverse branch, of the forward propagation of values and derivatives, of the finite quadrature and of the vertical tail. Missing: explicit numerical values for every constant of the sector proofs. Current values ($c\approx4.4923299$, $h(1)\approx1.81460956$) are floating diagnostics only.
\item\label{ied:q:higher} \emph{Corrections beyond the first, and an all-orders theorem.} Sources: Part~I, Subsection~\ref{ied:lead:sub:higher}; Part~II, Subsection~\ref{ied:corr:sub:allorders}; Part~III, Subsection~\ref{ied:log:sub:other} (an all-orders theorem in $1/m$). The expected form is $1+\sum_rP_r(\log n)/n^r$ with $\deg P_r\le2r$; Part~II proves $P_1$. Part~II's conditional inference, that $[X^{2r}]P_r=(-1/18)^r/r!$ (so $\log^4n/(648n^2)$ at second order) \emph{if} an all-orders construction gives the degree bound and no other contribution reaches degree $2r$, is unproved. Missing: an induction over subsidiary orders, clock hierarchies with resonant logarithms (Part~II's $N$ already has $\frac1{540}\log y$), forced Fatou cohomological corrections, and growing-strip errors with an extraction order chosen per precision.
\item\label{ied:q:effective} \emph{Effective constants, onsets and certified thresholds.} Sources: Part~I, Subsection~\ref{ied:lead:sub:effective}; Part~II, Subsection~\ref{ied:corr:sub:effective}; Part~III, Subsection~\ref{ied:log:sub:effective}. All three Parts give qualitative constants only; their inverses recover $n$ at exact sequence values eventually, without onset, and Part~III's prescriptions are smooth. Missing: explicit clock, stopping and sector constants; for least heights in Part~III, discrete information (monotonicity) about $R(n,m)$, which no Part supplies.
\item\label{ied:q:floorinverse} \emph{Inverse statements for $A(n,\lfloor\lambda n\rfloor)$.} Source: Part~I, Subsection~\ref{ied:lead:sub:effective}, second paragraph. The floor phase $e^{-\{\lambda n\}/\lambda}$ of \eqref{ied:lead:eq:floor} must enter; a smooth approximation that drops it has a non-vanishing multiplicative error. Not treated by any Part (Part~III's inverses concern $m\asymp n\log n$).
\item\label{ied:q:endpoint} \emph{The regime $n/m\to\infty$, and corrections near $n/m\to0$.} Sources: Part~I, Subsection~\ref{ied:lead:sub:heights}; Part~II, Subsection~\ref{ied:corr:sub:endpoint}; Part~III, Subsection~\ref{ied:log:sub:other}. The side $n/m\to0$ is answered at leading order for logarithmic bands by Part~III (Theorem~\ref{ied:log:thm:transfer}); the large-argument behaviour of $h$ and any transfer as $\beta\to\infty$ are open, and no Part proves a first correction of Part~II's kind on a shrinking band.
\item\label{ied:q:analytic} \emph{Analytic structure of $h$.} Sources: Part~I, Subsection~\ref{ied:lead:sub:heights} (``stronger regularity, such as real analyticity''); Part~III, Subsection~\ref{ied:log:sub:endpoint} (analyticity of the germ $g$ at zero, convergence of its Taylor series, coefficient growth and summability, with the nonzero resiter $1/3$). Known: $h\in C^\infty(0,\infty)$ (Corollary~\ref{ied:lead:cor:smooth}); $h$ continues holomorphically to a small punctured sector about $(0,\epsilon)$, so $h$ is real analytic on $(0,\epsilon)$ (Theorem~\ref{ied:log:thm:hankel}); $g$ has a one-sided $C^\infty$ germ (Theorem~\ref{ied:log:thm:germ}). Missing: real analyticity on all of $(0,\infty)$ (the measure identity \eqref{ied:lead:eq:measureidentity} is the suggested route; item~\ref{ied:q:ibd} would give it), and any statement on the Taylor series of $g$.
\item\label{ied:q:polysmall} \emph{Polynomially small ratios and Kaneiwa's regime.} Sources: Part~III, Subsection~\ref{ied:log:sub:polysmall}; Part~I, Subsection~\ref{ied:lead:sub:heights}, second paragraph; Part~II, Subsection~\ref{ied:corr:sub:endpoint}, second paragraph. Part~III's transfer loses $\beta^{-30}$ and stops at logarithmic bands; a uniform theory linking proportional height, the crossover and Kaneiwa's fixed-degree polynomials (Section~\ref{ied:log:sec:kaneiwa}) is open. Part~III's suggested step: lower derivative orders or the retained first correction in the absolute Fourier norm, with balanced endpoint errors.
\item\label{ied:q:other} \emph{Other arrays and initial functions.} Sources: Part~I, Subsection~\ref{ied:lead:sub:heights}; Part~II, Subsection~\ref{ied:corr:sub:endpoint}; Part~III, Subsection~\ref{ied:log:sub:other}. The forcing and clock constants ($T$, $B$, $\kappa$) must be recomputed for another multiset recurrence or initial function; the endpoint calculus of the unforced map may persist. Missing: the forced entrance and tail estimates in each new case.
\item\label{ied:q:ibd} \emph{The amplitude identity with the iterated Bell report} (a question shared with \file{a139383-iterated-bell-diagonals}). Source: Part~III, Section~\ref{ied:log:sec:interfaces} (``a formal contour match suggests an amplitude relation $I(\beta)=\beta\,2^{\beta/3}h(\beta)$, but that identification requires its own analytic contour correspondence and is not used here'') and Subsection~\ref{ied:log:sub:endpoint} (``a complete analytic comparison with the public IBD contour amplitude''). Here $I$ is the iterated Bell report's Part~II amplitude, $H(n,m)\sim(n-1)!\,2^{-n}(\lambda n)^n(\lambda n)^{-1/(3\lambda)}I(1/\lambda)$ for $m=\lambda n$ (its \texttt{ibd:pd:eq:factorial}), defined there by a fixed-radius contour functional. Sketch: both reports use the same normalized reciprocal orbit, and $\Psi_{\mathrm{IBD}}=-\Phi+\frac13\log2$ holds locally; the iterated Bell report's own literature section derives $I(\beta)=\beta J_P(\beta)$ from Prellberg's contour integral under the same kind of unproved contour correspondence. Evidence: agreement to eight significant digits at $\beta=1$ (Remark~\ref{ied:rem:ibdnumbers}). Missing: an analytic identification of that report's contour functional with the Laplace inversion of $P$ (orientation, normalization, vanishing endpoint terms). Consequences if proved: (a)~$h$ is real analytic on $(0,\infty)$, because that report proves $I$ real analytic there (item~\ref{ied:q:analytic}); (b)~Part~III's endpoint results ($h(0)=2$, Corollary~\ref{ied:log:cor:densityseries}, \eqref{ied:log:eq:logdensity}) give the behaviour of $I(\beta)$ as $\beta\to0$, that is, the amplitude half of the endpoint $\lambda\to\infty$ that the iterated Bell report leaves open; a labelled analogue of Theorem~\ref{ied:log:thm:transfer} would still be needed for a coefficient theorem (Remark~\ref{ied:rem:labelled}).
\end{enumerate}

\begin{remark}[\wtag\ The identity at $\beta=1$, numerically, 6~October 2026]\label{ied:rem:ibdnumbers}
The iterated Bell report's shipped exploratory runs (\file{data/53-iterated-bell-exploratory-amplitude.json} there) give $I(1)=2.2862647982$, $2.2862647847$, $2.2862647863$ (its $C=\sqrt{2\pi}I(1)/2$ between $2.865407976$ and $2.865407993$). Part~I's shipped diagnostic $h(1)=1.8146095596784502$ (\file{data/225-leading-data-amplitude_diagnostic.json}) gives $2^{1/3}h(1)=2.2862647816$. The two sides agree to eight significant digits ($2.2862648$); their differences, $3\times10^{-9}$ to $1.7\times10^{-8}$, are of the size of the spread of the three iterated Bell runs ($1.4\times10^{-8}$). \textup{(Added after the independent check of 6~October 2026: an independent evaluation of \eqref{ied:lead:eq:constantintegral} by the verification (four values of $\sigma$, Gauss--Legendre quadrature to $|u|=4\cdot10^4$ with an integration-by-parts tail; uncertified) gives $h(1)=1.814609559835$ and $2^{1/3}h(1)=2.2862647818$, so the shipped diagnostic is low by $1.6\times10^{-10}$, from its truncation of the vertical tail at $|u|=3000$; the comparison is unchanged, with differences $2.9\times10^{-9}$ to $1.65\times10^{-8}$ and agreement to eight significant digits.)} Both sides are floating values without error control; this is evidence for item~\ref{ied:q:ibd}, not a proof of it, and it says nothing about $\beta\ne1$.
\end{remark}

\begin{remark}[\wtag\ A conditional consequence for the labelled array, 6~October 2026]\label{ied:rem:labelled}
Let $H(n,m)=n![z^n](e^z-1)^{\circ m}$ and $R_{\mathrm{lab}}(n,m)=H(n,m)/(n!\,m^{n-1}2^{1-n})$. Suppose, for some $\beta_0>0$ and $D>1$, (a)~$I(\beta)=\beta2^{\beta/3}h(\beta)$ for $0<\beta\le\beta_0$, and (b)~$H(n,m)=(n-1)!\,2^{-n}m^nm^{-\beta/3}I(\beta)\{1+\epsilon(n,m)\}$ uniformly for $(\log m)^{-D}\le\beta=n/m\le\beta_0$, with $\epsilon\to0$ (the iterated Bell report's leading form \texttt{ibd:pd:eq:factorial} with $\lambda n=m$, extended to logarithmic bands). If $m/(n\log n)\to\lambda\in(0,\infty)$, then $R_{\mathrm{lab}}(n,m)\to e^{-1/(3\lambda)}$, the same limit as Part~III's Theorem~\ref{ied:log:thm:crossover}; and if $\epsilon=O(L^{-N-1})$ with $L=\log n$ and $m=\lfloor\lambda nL\rfloor$, then
\[
 \log R_{\mathrm{lab}}(n,m)=-\frac1{3\lambda}+\frac{K_{\mathrm{lab}}}{\lambda L}+\sum_{j=2}^N\frac{c_j}{(\lambda L)^j}+\Oh(L^{-N-1}),\qquad K_{\mathrm{lab}}=\frac{1-\gamma}3,
\]
with the same $c_j$ as \eqref{ied:log:eq:cs}.
\end{remark}
\begin{proof}
By (a) and (b), since $(n-1)!\,m^n\beta=n!\,m^{n-1}$,
\[
 R_{\mathrm{lab}}(n,m)=2^{\beta/3}m^{-\beta/3}\frac{h(\beta)}2\{1+\epsilon(n,m)\}.
\]
By \eqref{ied:log:eq:logdensity}, $\log(h(\beta)/2)=-\frac\beta3\log\beta+\frac{C_0}3\beta+\sum_{j=2}^Nc_j\beta^j+\Oh(\beta^{N+1})$ with $C_0=1-\gamma-\log2$. Hence, using $\log\beta+\log m=\log n$ exactly as in \eqref{ied:log:eq:logmaster},
\[
 \log R_{\mathrm{lab}}=-\frac\beta3\log n+\frac{\log2+C_0}3\beta+\sum_{j=2}^Nc_j\beta^j+\Oh(\beta^{N+1})+\log(1+\epsilon),
\]
and $\log2+C_0=1-\gamma$. If $m/(n\log n)\to\lambda$, then $\beta\to0$, $\beta\log n\to1/\lambda$, and $\beta$ eventually lies in the band of (b): $\beta\le\beta_0$, and $\beta\ge1/(2\lambda\log n)\ge(\log m)^{-D}$ because $D>1$ and $\log m\sim\log n$. So the right side tends to $-1/(3\lambda)$. For $m=\lfloor\lambda nL\rfloor$, $\beta=1/(\lambda L)+\Oh(1/(nL^2))$, whose effect is smaller than every power of $L$, as in the proof of Theorem~\ref{ied:log:thm:crossover}; substituting gives the display.
\end{proof}
Hypothesis (b) is not proved anywhere: the iterated Bell report's Theorem (its \texttt{ibd:pd:thm:main}) is uniform only for $\lambda=m/n$ in compact subsets of $(0,\infty)$, and it records the endpoints $\lambda\to0,\infty$ as open. A consistency check, not a proof: with $h(0)=2$ the right side of (b) is $n!\,m^{n-1}2^{1-n}(1+o(1))$ when $\beta\log m\to0$, the leading term of the fixed-$n$ asymptotic $H(n,r+1)\sim n!\,r^{n-1}/2^{n-1}$ of Skau and Kristensen that the iterated Bell report cites. Exact values of $H(n,m)$, computed by the write from $H(n,m+1)=\sum_jS(n,j)H(j,m)$ with Stirling numbers $S(n,j)$ and evaluated in 30-digit floating arithmetic (uncertified), behave as the remark predicts, with the heights $m=\lfloor\lambda n\log n\rfloor$ of Part~III's Table~\ref{ied:log:tab:finite}:
\begin{center}
\small
\begin{tabular}{rrrrrr}
\toprule
$n$ & $\lambda$ & $m$ & $R_{\mathrm{lab}}(n,m)$ & $e^{-1/(3\lambda)}$ & $e^{-1/(3\lambda)}\bigl(1+K_{\mathrm{lab}}/(\lambda\log n)\bigr)$\\
\midrule
20 & 0.5 & 29 & 0.56666331 & 0.51341712 & 0.56172244\\
20 & 1 & 59 & 0.75860855 & 0.71653131 & 0.75023906\\
20 & 1.5 & 89 & 0.83316752 & 0.80073740 & 0.82585011\\
40 & 0.5 & 73 & 0.55659049 & 0.51341712 & 0.55264578\\
40 & 1 & 147 & 0.74875697 & 0.71653131 & 0.74390532\\
40 & 1.5 & 221 & 0.82522543 & 0.80073740 & 0.82113138\\
80 & 0.5 & 175 & 0.54955621 & 0.51341712 & 0.54644061\\
80 & 1 & 350 & 0.74202367 & 0.71653131 & 0.73957531\\
80 & 1.5 & 525 & 0.81979276 & 0.80073740 & 0.81790547\\
\bottomrule
\end{tabular}
\end{center}
At $\lambda=1$ the distance to the limit is $0.042$, $0.032$, $0.025$ at $n=20,40,80$, and after the $K_{\mathrm{lab}}$ term $0.0084$, $0.0049$, $0.0024$. The Euler constant $K\approx0.816$ of Part~III and $K_{\mathrm{lab}}\approx0.141$ differ, as the subsidiary forcing leads one to expect. These finite values neither prove the limit nor certify any digit.

\begin{writenote}[Independent check of the write, 6~October 2026]
An adversarial check made by the intake after the write (commit \file{84fc1aaed}) re-pulled A290354 (live revision \#26 of 21~June 2018) and A290353 (\#28) and found every quoted line character for character, including the formula line in the front matter, in the note at the start of Part~I and in the bibliography. The entry's history places the conjecture in Kot\v{e}\v{s}ovec's revision \#23 of 14~August 2017, so the date ``Aug 14 2017'' that the write gives is the date of the conjecture, not of the revision read, and it is correct. The check recomputed $a_n$ as the $n$-th Euler transform of $1+x$ in integer arithmetic (equal to the b-file for $n\le40$; $a_{40}(T/40)^{40}40^{4/3}=4.6013347$), confirmed that Theorem~\ref{ied:lead:thm:main} is exactly the conjectured form, and the rows $A(2,m)$, $A(3,m)$, $A(4,m)$ for $m\le12$. It re-derived Remark~\ref{ied:rem:labelled} from (a), (b) and \eqref{ied:log:eq:logdensity}, found it labelled conditional wherever it appears (title, hypotheses, the paragraph after the proof, the README) and not listed as proved in the status table, and reproduced its table to every printed digit from an own exact computation of $H(n,m)$; extended to $n=160$ it gives $m=812$, $R_{\mathrm{lab}}=0.73782182$ at $\lambda=1$ (distance $0.0213$ to the limit, $0.0014$ after the $K_{\mathrm{lab}}$ term) and $m=406$, $R_{\mathrm{lab}}=0.54374893$ at $\lambda=0.5$ ($0.0303$ and $0.0018$). It evaluated the closed form \eqref{ied:corr:eq:kappa} to 35 digits, $\kappa=-0.25434816985107480385506489650373339$, and matched it against an independent quadrature of the defining equation \eqref{ied:corr:eq:M} to $5\times10^{-28}$. It recomputed $h(1)$ independently (Remark~\ref{ied:rem:ibdnumbers}, dated note) and confirmed the classification of the inverses in the front matter: the factorial core and the Lambert core with $a=b=1$ are instances; the corrections, the rounding at exact values and Part~III's reversion and prescriptions are not. Two wordings are corrected, each with a dated note keeping the first wording: the front matter's justification of the printed restatements (twice) overstated Part~II's novelty, since its bound on inverse orbits is in Part~I's proof of Lemma~\ref{ied:lead:lem:fatou}; and the note after Part~II's Section~\ref{ied:corr:sec:interface} now flags the neighbourhood form of the outer bound, as it already flagged the disk form of the tail bound, as proved only inside Part~I's proofs. Added: the name $L=\log(Y/T)$ in the front matter, and the independent values of $h(1)$ and $c=4.4923298971$. No mathematical error was found.
\end{writenote}
\begin{writenote}[Part~IV and these questions, 9 October 2026]
Part~IV (Section~\ref{ied:par:sec:front}) answers part of these questions;
the earlier notes, written before it arrived, are kept.
\emph{Item~\ref{ied:q:ibd}} (the amplitude identity): \textbf{answered},
conditionally on Part~IV's four inputs (Section~\ref{ied:par:sec:inputs}):
$I(t)=t\,2^{t/3}h(t)$ for every $t>0$ (Theorem~\ref{ied:par:thm:identity}).
The proof identifies the limits of the positive measures
$\mu_m=\sum_n[z^n]f^{\circ m}(z)P(-m)^n\delta_{n/m}$, whose masses are exactly
$P(0)$; it needs no global comparison of contours, the route item~\ref{ied:q:ibd}
calls missing. Both consequences listed there follow: (a)~$h$ is real
analytic on $(0,\infty)$ (below); (b)~the endpoint behaviour of $I$ as
$\beta\to0$ is Corollary~\ref{ied:par:cor:endpoint}, the amplitude half of the
iterated Bell report's endpoint $\lambda\to\infty$; a uniform coefficient
transfer toward that endpoint stays open (Part~IV's
Question~\ref{ied:par:q:uniform}). \emph{Item~\ref{ied:q:analytic}}:
\textbf{answered on the positive axis}: $h$ extends holomorphically to
$\Re t>0$ \eqref{ied:par:eq:hhol}. Analyticity of the germ $g$ at zero and its
Taylor series stay open (Part~IV's Question~\ref{ied:par:q:germ}).
\emph{Remark~\ref{ied:rem:labelled}}: its hypothesis~(a) is now
Theorem~\ref{ied:par:thm:identity} (for all $\beta>0$, under the same inputs);
hypothesis~(b) is still unproved, so the remark remains conditional.
\emph{Remark~\ref{ied:rem:ibdnumbers}}: the identity now forces
$I(1)=2^{1/3}h(1)$. With Part~IV's exploratory density
($h(1)=1.81460955983\ldots$, within $10^{-12}$ of the verification's
independent $1.814609559835$) the forced value is
$2.2862647817\ldots$; the iterated Bell report's three exploratory runs exceed
it by $2.9\times10^{-9}$, $4.5\times10^{-9}$ and $1.6\times10^{-8}$. Every value
involved is an uncertified floating evaluation, and the spread of the three
runs is $1.4\times10^{-8}$, so this records an offset of the runs, not a
refutation of anything; the remark's ``eight significant digits'' stays true.
\emph{Items~\ref{ied:q:validated} and~\ref{ied:q:endpoint}} reappear as Part~IV's
Questions~\ref{ied:par:q:certified} (enclosures of $h(\log2)$, $h'(\log2)$,
$h''(\log2)$) and~\ref{ied:par:q:large} (large-argument asymptotics of $h$).
\emph{An observation} [write], uncertified: from exact chain counts (the Newton
form $\sum_me^{-Lm}H(n,m)=\sum_k(N^k)_{n1}e^{-Lk}(1-e^{-L})^{-k-1}$, with $N$ the
strictly triangular part of the Stirling matrix) and Part~IV's exploratory
$h$, the ratio of $S_n(L)$ to the main term of
Theorem~\ref{ied:par:thm:aggregate} satisfies $n(\text{ratio}-1)=0.11556$ at
$L=\log2$ and $0.16669$ at $L=1$ for $n=320$ (monotone in $n$ from $n=25$), against
$L/6=0.11552$ and $0.16667$: the first correction looks like $L/(6n)$ with no
logarithmic term, as the chain report proves at $q=1$
($z_n/f_n=C(1+L/(6n)+O(n^{-2}))$). At $L=1$ ($q=1/(e-1)$, outside that
report's disk) this is evidence on Part~IV's Question~\ref{ied:par:q:summation},
not a proof.
\end{writenote}

\clearpage
\part{A common density for parabolic iteration}\label{ied:par:part}
% Part IV numbers equations within sections, as delivered.
\counterwithin*{equation}{section}
\renewcommand{\theequation}{\thesection.\arabic{equation}}
\renewcommand{\theHequation}{\theHsection.\arabic{equation}}
% Part IV's \Oh is a calligraphic O (Parts I-III use an upright O).
\renewcommand{\Oh}{\mathcal O}
\partsource{Manuscript~02 of the archive \texttt{ProveIt\_Research\_2026-10-08\ (1).zip},
delivered as \file{02_parabolic_amplitudes/article.tex} (16~pages). Delivered
title block ``A Common Density for Parabolic Iteration / Exact Bell--Euler
amplitude identification and weighted partition-chain constants / Research
companion prepared for Vladimir Reshetnikov / Developed with OpenAI ChatGPT /
8 October 2026''. Printed as delivered after Section~\ref{ied:par:sec:front},
which the write added, with \textbf{[write]} notes; the manuscript's
Section~$k$ is Section~$k+38$ here, its statements and equations ($k.m$) are
($k{+}38.m$), and its Appendix~A keeps its letter. Its labels carry the prefix
\file{ied:par:}. ``\cite{Euler}'' in this Part is Parts~I--III of this report,
``\cite{IBD}'' the iterated Bell report, ``\cite{Chains}'' the strict
partition-chain report.}
\begin{partabstract}
Two recent ProveIt reports study different arrays governed by the same
parabolic map $f(z)=e^z-1$. One constructs a positive density $h$ whose
Laplace transform is the normalized entire Fatou orbit $P$; the other
obtains an iterated-Bell amplitude $I$ from a fixed contour.
Their precise relation was left as an explicit open question. We prove
\[
 I(t)=t\,2^{t/3}h(t)\qquad(t>0)
\]
by identifying two limits of positive coefficient measures. This avoids
a global comparison of the two contours and proves that $h$ extends
holomorphically to $\Re t>0$. A coefficientwise rational majorant then
controls every iteration depth in a geometric aggregate. It gives the
exact Lengyel-constant identity and, more generally, a holomorphic
formula on $\Re q>0$ for the leading amplitude of strict partition chains
marked by $q$. Complete proofs, source pins, exact finite checks and
independent numerical diagnostics are supplied. The analytic inputs
from the earlier reports are identified explicitly; no worldwide
priority claim or new leading-equivalent claim is made.
\end{partabstract}

\section{Provenance, inputs and notation of Part~IV}\label{ied:par:sec:front}

[Added 9~October 2026.] Part~IV was added in a second write of this report,
after Parts~I--III had been merged (\file{84fc1aaed}) and independently
checked (\file{a3a2d5582}). It is the manuscript \emph{A Common Density for
Parabolic Iteration}, which starts from item~\ref{ied:q:ibd} of
Section~\ref{ied:sec:further} (it cites it as ``Section~37, item~9'') and answers
it. Everything of the manuscript is printed in this Part except its title page
(recorded in the source block above) and its reference list (merged into this
report's bibliography): its abstract above, its ten sections as
Sections~\ref{ied:par:sec:intro}--\ref{ied:par:sec:further} and its
Appendix~\ref{ied:par:app:checklist}. Its letters are kept
(Section~\ref{ied:par:sec:notation}). This section and the notes marked
[write] are the write's; all other text of Sections~\ref{ied:par:sec:intro}--\ref{ied:par:sec:further}
and of the appendix is the manuscript's.

\subsection{The manuscript}\label{ied:par:sec:provenance}

The archive \texttt{ProveIt\_Research\_2026-10-08\ (1).zip} ($1{,}394{,}668$ bytes;
45 files in one wrapper directory \file{ProveIt_Research_2026-10-08/},
$1{,}636{,}954$ bytes unpacked) arrived in commit \file{b28d0850b}
(8~October 2026, 08:39~PDT). It holds three manuscripts; this is its
manuscript~02, \file{02_parabolic_amplitudes/} (12 files). Its 9 supporting
files were placed in this directory by commit \file{803f4d937}
(9~October 2026) with the prefix \file{04-parabolic-}, the Part number: the
three programs as \file{code/04-parabolic-*.py}, the four data files,
\file{provenance.json} and \file{requirements.txt} as \file{data/04-parabolic-*};
all 9 are byte-identical to the delivery (confirmed at the write). Not shipped,
and retrievable from the arrival commit: the manuscript \file{article.tex}, its
PDF and its \file{README.md}. The archive is not retired: its manuscripts~01
and~03 and its package-level files (among them \file{REVIEW_NOTES.md},
\file{THEOREM_LEDGER.json}, \file{RESEARCH_QUESTIONS.json}, \file{reproduce.py})
belong to their own placements.

The manuscript is \file{article.tex} (1{,}079 lines, dated 8~October 2026); a
build of the delivered text gives 16~pages, as delivered, with no warnings.
Its title page reads ``Research companion prepared for Vladimir Reshetnikov /
Developed with OpenAI ChatGPT'' (PDF author field the same). Its pin
(Section~\ref{ied:par:sec:literature}.2, \file{data/04-parabolic-provenance.json}
and its three bibliography entries) is ProveIt commit
\file{58175ca45563d9ee29875374dd77942f069268a9} (7~October 2026, 20:30~PDT), with
the blobs \file{b5af1ca30b} of this report's \file{article.tex} (Parts~I--III
as merged and checked, unchanged until this write), \file{2b1d8fbe81} of the
iterated Bell report and \file{da8cac4780} of the strict partition-chain
report, both unchanged at HEAD.

\emph{How the manuscript was merged.} Its 56 labels, and 6 labels the write gave
to its unlabelled Sections~1, 4, 7, 9, 10 and Appendix~A, carry the prefix
\file{ied:par:}; its 51 references (41 \verb|\eqref|, 10 \verb|\ref|) were
updated. Its citation keys are entries of the merged bibliography, marked
[Part~IV], except two: \file{BellReport} is \file{IBD} (4 citations re-keyed;
Part~IV's pin and description added) and \file{Prellberg} is the entry of
Parts~I--III (Part~IV's description added); its key \file{Euler} is this
report at the pin, that is, Parts~I--III, and its \file{OEIS} (A005121, listed
but not cited) is \file{A005121}. Its macros \verb|\coef| and \verb|\Repin| were
added to the preamble; its \verb|\C|, \verb|\R|, \verb|\N|, \verb|\dd|,
\verb|\ii|, \verb|\file|, \verb|\Log| are those of Parts~I--III; its
\verb|\Oh| (a calligraphic $\mathcal O$, an upright $\mathrm O$ in
Parts~I--III) is switched at the start of this Part. From this Part on,
equations are numbered within sections, as delivered. The delivered font
(Palatino), running heads and paragraph skip are not reproduced. No other text
was changed. The write labelled the nine items of its
Section~\ref{ied:par:sec:further}. The numbering, against a build of the
delivered text: its Sections~1--10 are
Sections~\ref{ied:par:sec:intro}--\ref{ied:par:sec:further} here, its
Theorems~1.1 and~1.2 are Theorems~\ref{ied:par:thm:identity}
and~\ref{ied:par:thm:aggregate}, Corollary~1.3 is
Corollary~\ref{ied:par:cor:marked}, Corollary~4.2 is
Corollary~\ref{ied:par:cor:endpoint}, Proposition~4.3 is
Proposition~\ref{ied:par:prop:general}, Corollary~7.1 is
Corollary~\ref{ied:par:cor:K}; every statement and equation $k.m$ is
$(k{+}38).m$; Appendix~A is unchanged (all 56 labels compared).

\emph{What the manuscript says about the repository.} Its account of the
three reports is correct: item~\ref{ied:q:ibd} (``Section~37, item~9'') asks
for the identity; the iterated Bell report records the same question (its
dated note at the end of \file{ibd:sec:literature}); the labels it names as
resolved are \file{ied:q:ibd}, the positive-axis part of \file{ied:q:analytic},
the leading-amplitude part of \file{spc:q:depth}, and in part
\file{spc:q:disk}, as the table of Section~\ref{ied:par:sec:answers} states. The
chain report's \file{spc:q:depth} carries the \emph{heuristic}
$C\overset{?}{=}L^{L/3-1}I(L)/2$ of that report's write, with $L=\log2$;
\eqref{ied:par:eq:Lengyel} proves it.

\subsection{The trust boundary: four inputs}\label{ied:par:sec:trust}

Every theorem of this Part is an implication. Its main theorem uses exactly
the four statements of Section~\ref{ied:par:sec:inputs}, each proved in an
unrefereed, unformalized report of this collection:
\begin{itemize}
\item \emph{Input~F}: the normalized Fatou coordinate and the entire orbit,
Part~I's Lemma~\ref{ied:lead:lem:fatou} and~\eqref{ied:lead:eq:fatou};
\item \emph{Input~M}: the positive Laplace measure with continuous, strictly
positive density, Part~I's Proposition~\ref{ied:lead:prop:measure};
\item \emph{Input~L}: the absolute local coefficient transfer
\file{ibd:pd:eq:absolute} of the iterated Bell report, uniform for $n/m$ in
compact subsets of $(0,\infty)$; that report's Part~II is itself conditional on
its Part~I, whose outer arcs rest on an interval certificate of common entry
(Remark~\ref{ied:par:rem:seed} says so);
\item \emph{Input~A}: the holomorphy part of \file{ibd:pd:prop:positive}
there (not its numerical positivity seed $I(1)>1.64$).
\end{itemize}
The corollaries use more: Corollary~\ref{ied:par:cor:endpoint} uses Part~III's
\eqref{ied:log:eq:logdensity}; Corollary~\ref{ied:par:cor:marked}'s
agreement near $q=1$ and Corollary~\ref{ied:par:cor:K} use the strict
partition-chain report's marked expansion and its cumulant expansions
\eqref{ied:par:eq:mean}--\eqref{ied:par:eq:variance}. The manuscript states
this boundary itself (``The present proof checks the implication from their
stated analytic interfaces''), and its package record (\file{REVIEW_NOTES.md},
not shipped here) lists among its remaining trust boundaries that ``The
parabolic source reports are unrefereed'' and that its inverse-Laplace
calculations ``are not interval-certified''. Nothing here makes the inputs
more certain than their sources: the theorems of this Part have exactly the
status of Parts~I--III and of the iterated Bell report together. Prellberg's
2002 seminar, in Mishna's summary~\cite{Prellberg}, is credited for the
bivariate parabolic mechanism and an aggregate contour constant; no
new-equivalent or priority claim is made.

\subsection{What Part~IV answers}\label{ied:par:sec:answers}

\begingroup\small
\renewcommand{\arraystretch}{1.12}
\renewcommand{\theHtable}{iedpar.answers}% no PDF destination clash
\begin{longtable}{@{}>{\raggedright\arraybackslash}p{0.27\textwidth}>{\raggedright\arraybackslash}p{0.67\textwidth}@{}}
\toprule
Question & After Part~IV\\
\midrule
\endhead
Item~\ref{ied:q:ibd} (the amplitude identity) & \textbf{Answered} (Theorem~\ref{ied:par:thm:identity}), by limits of positive coefficient measures, without a comparison of contours.\\
Item~\ref{ied:q:analytic} (analytic structure of $h$) & \textbf{Answered on the positive axis}: $h$ holomorphic on $\Re t>0$ \eqref{ied:par:eq:hhol}. \textbf{Open}: the germ $g$ at zero and its Taylor series.\\
Remark~\ref{ied:rem:labelled} & Hypothesis~(a) proved; (b) open; the remark stays conditional.\\
Remark~\ref{ied:rem:ibdnumbers} & $I(1)=2^{1/3}h(1)$ is now forced; the iterated Bell report's uncertified runs sit $2.9\times10^{-9}$ to $1.6\times10^{-8}$ above the forced value (a record, not a refutation).\\
Iterated Bell report, endpoint $\lambda\to\infty$ & The amplitude half: Corollary~\ref{ied:par:cor:endpoint}; the uniform transfer stays open (Question~\ref{ied:par:q:uniform}).\\
Chain report, \file{spc:q:depth} & The leading-amplitude part: its heuristic $C=L^{L/3-1}I(L)/2$ is \eqref{ied:par:eq:Lengyel}, a theorem; all-order summation and the logarithmic cancellation stay open (Question~\ref{ied:par:q:summation}).\\
Chain report, \file{spc:q:disk} & In part: $C(q)$ holomorphic on $\Re q>0$, positive on $q>0$; the zero set and global uniform asymptotics stay open (Questions~\ref{ied:par:q:zeros}, \ref{ied:par:q:positive}).\\
Items~\ref{ied:q:validated}--\ref{ied:q:other} otherwise & Untouched; items~\ref{ied:q:validated} and~\ref{ied:q:endpoint} reappear as Questions~\ref{ied:par:q:certified} and~\ref{ied:par:q:large}.\\
\bottomrule
\end{longtable}
\addtocounter{table}{-1}% the uncaptioned longtable takes no number
\endgroup

Dated notes at the end of Section~\ref{ied:sec:further} and of the front
matter say the same; the notes for the two neighbouring reports are proposed
separately and not applied.

\emph{What Part~IV proves again}: nothing of Parts~I--III; its Lemma~\ref{ied:par:lem:radius}
re-derives the radius from Input~F, and its Lemma~\ref{ied:par:lem:mixture}
the chain report's exact mixture $z_n=\sum_m2^{-m-1}H(n,m)$ at $q=1$.
\emph{What is new to the report}: everything else.

\subsection{What the write read and checked}\label{ied:par:sec:checks}

\emph{At placement} (the intake dossier, 9~October 2026; own code): from the
delivered exploratory $h(\log2)$, $C(1)=(2L)^{L/3}h(L)/2=1.09868580552518740\ldots$
against Lengyel's constant, OEIS A086053, $1.09868580552518701\ldots$ (agreement
to 15 digits); $I(\log2)=1.65770061282$, consistent with the chain report's
finite-$n$ extrapolations; the exact counts against A005121 and A139383; and
$2^{1/3}h(1)=2.28626478178$ against the iterated Bell report's runs.

\emph{At the write} (9~October 2026; Python~3.14.4, SymPy~1.14.0, mpmath~1.3.0;
own code, kept with the intake records and not shipped). Every proof of
Sections~\ref{ied:par:sec:measures}--\ref{ied:par:sec:constants} was read line
by line, with no error found; and the four inputs were read at their sources
(Part~I's Lemma~\ref{ied:lead:lem:fatou} and
Proposition~\ref{ied:lead:prop:measure}; \file{ibd:pd:eq:absolute}, stated there
for $\eta=n/m\in[b,B]$ with remainder $\Oh(m^{-1}(1+\log m)^D)$, and the
holomorphy part of \file{ibd:pd:prop:positive}), where they say what
Section~\ref{ied:par:sec:inputs} quotes. Rechecked by hand or with SymPy in
particular: the three sensitive normalizations after the proof of
Theorem~\ref{ied:par:thm:identity}; the main term of
Theorem~\ref{ied:par:thm:aggregate} from $(n-1)!\,2^{-n}(n/L)^{-L/3}I(L)$ times
$n!L^{-n-1}$; \eqref{ied:par:eq:Lengyel} and~\eqref{ied:par:eq:Cqh}; the endpoint
expansion~\eqref{ied:par:eq:Iendpoint} and its exponentiation
\eqref{ied:par:eq:Iendfirst}; the chain rule behind
Corollary~\ref{ied:par:cor:K} and both expanded forms; the radius model of
Lemma~\ref{ied:par:lem:radius} (numerically, residual $O(\log m/m)$ at
$m=10^3,10^6$). With exact integer and rational arithmetic: the Newton form
$H(n,m)=\sum_k\binom mk(N^k)_{n1}$ against the recurrence
\eqref{ied:par:eq:Hrec}; the chain counts against A005121; the majorant
\eqref{ied:par:eq:Hmajorant} for $n\le30$, $m\le60$; the mixture
\eqref{ied:par:eq:mixture} for $n\le25$ at five rational $q$, through the closed
form $\sum_ma^mH(n,m)=\sum_k(N^k)_{n1}a^k(1-a)^{-k-1}$. Numerically
(uncertified): Theorem~\ref{ied:par:thm:aggregate} at $L=\log2$ and $L=1$ for
$n\le320$, with ratios to the main term $1.00036$ and $1.00052$ at $n=320$
(note at the end of Section~\ref{ied:sec:further}); the iterated Bell
report's three runs against the forced $I(1)$. It did not rerun the shipped
programs. \emph{Not read by the write}: Prellberg's summary, Lengyel (1984) and
Skau--Kristensen; the manuscript's accounts of them are reported, not
confirmed.

\subsection{What is not claimed}\label{ied:par:sec:nonclaims}

From the manuscript: the theorems are implications from the four inputs, whose
sources ``are unrefereed and not proof-assistant verified'', and the proof
``does not promote their earlier proofs or interval software to formal
verification''; no numerical seed is needed for $I(1)$, but other
computational dependencies of the Bell foundation remain
(Remark~\ref{ied:par:rem:seed}); the leading equivalent and the existence of
Lengyel's constant are classical, and Prellberg's seminar is a direct
antecedent; the mean, variance and central limit theorem of chain lengths are
not new claims; the continuation of $C(q)$ does not give uniform complex
asymptotics or zero-freeness; the endpoint series does not give the uniform
transfer, and the germ is not shown convergent;
Proposition~\ref{ied:par:prop:general} is a transfer principle under stated
hypotheses; the inverse-Laplace values are diagnostics, not certified;
``A literature search and source inspection do not establish that no
equivalent theorem has appeared elsewhere''. The write adds: its checks are
exact finite computations, hand or SymPy algebra and floating evaluations; its
observation on the first correction of the aggregate is numerical; nothing was
submitted anywhere, and the notes for the two neighbouring reports are not
applied.

\subsection{Reading conventions}\label{ied:par:sec:notation}

Part~IV keeps its letters. Where they clash with Parts~I--III, read them as
follows.

\begingroup\small
\renewcommand{\arraystretch}{1.1}
\renewcommand{\theHtable}{iedpar.notation}% no PDF destination clash
\begin{longtable}{@{}>{\raggedright\arraybackslash}p{0.95in}>{\raggedright\arraybackslash}p{3.0in}>{\raggedright\arraybackslash}p{2.2in}@{}}
\toprule
Symbol & Part IV & Parts I--III\\
\midrule
\endhead
$H(n,m)$, $f$ & $n![z^n]f^{\circ m}$, the labelled iterated Bell array; $f(z)=e^z-1$ & $H(z)$ the matching function, $H_j$, $H_{\theta,a}$; $f(w)=e^w-1$\\
$I(t)$, $h$, $P$, $\Phi$ & the iterated Bell amplitude; the density; the orbit; the coordinate (Parts~I--III's) & $I$ (Section~\ref{ied:log:sec:interfaces}); $I$ a compact interval (II)\\
$t$, $\beta$ & $t=n/m$ & $\beta=n/m$\\
$L$, $L(q)$ & $L>0$ the tilt of $S_n(L)$; $L(q)=\log(1+1/q)$; $L=\log2$ at $q=1$ & $L(y)$ a clock, $L=\log n$ (III)\\
$C$, $C(q)$, $K_1$, $K_2$, $K$ & Lengyel's constant, the marked amplitude; cumulant constants; $K=\max(D,2)$ & $c$ the amplitude; $C$, $K$ constants of Parts~II--III\\
$S_n(L)$, $Z_n(q)$, $T_n$, $\mathcal H_n$ & the geometric aggregate; marked chain counts; $T_n=\sum_ma^mH(n,m)$; the number of strict transitions & $T=\pi/2$\\
$r_m$, $u_m$, $\mu_m$, $s_m$ & $P(-m)$; $2/r_m$; the coefficient measures; $m+\Phi(r_me^{s/m})$ & $r$, $r_m$ radii; $s$ a contour parameter\\
$A(x)$, $a$, $\rho$ & $\log h(x)+(x/3)\log(2x)$; $a=q/(q+1)$ or the $a$ of Proposition~\ref{ied:par:prop:general}; $\rho$ the log coefficient there & $A(n,m)$ the Euler array; $a_n$; $\rho_j$\\
$\mathcal W_n(L)$, $w_n$ & the central depth window; $\sqrt n(\log n)^2$ & \\
$c_j$, $D$ & Part~III's log-density coefficients; the logarithmic degree of Input~L & $c_j$; $D$ a band exponent (III)\\
\bottomrule
\end{longtable}
\addtocounter{table}{-1}% the uncaptioned longtable takes no number
\endgroup

\section{Question, result and scope}\label{ied:par:sec:intro}

Let $f(z)=e^z-1$, and let $f^{\circ m}$ denote $m$ iterations of $f$,
with $f^{\circ0}(z)=z$. Define
\begin{equation}\label{ied:par:eq:H}
 H(n,m)=n!\coef n f^{\circ m}(z),\qquad n\ge1,\quad m\ge0.
\end{equation}
The array $H$ is the iterated exponential Bell array.
The principal objects here are functions that describe its asymptotic
coefficients, rather than a new interpretation of these integers.

The report on iterated Euler transforms \cite{Euler} constructs an
entire function $P$ satisfying
\begin{equation}\label{ied:par:eq:P}
 P(s+1)=e^{P(s)}-1,\qquad
 P(s)=\int_0^\infty e^{st}h(t)\dd t.
\end{equation}
Its additive normalization is fixed by the inverse coordinate
\begin{equation}\label{ied:par:eq:Phi}
 \Phi(w)=-\frac2w+\frac13\log w+\Oh(w),\qquad \Phi(P(s))=s
\end{equation}
near the positive real $w$-axis, equivalently for sufficiently negative
real $s$. The density $h$ is continuous and strictly positive on
$(0,\infty)$, and the measure in \eqref{ied:par:eq:P} has every exponential moment.

The iterated-Bell report \cite{IBD} defines a second function $I$
by a fixed-radius contour and proves the local coefficient formula
\begin{equation}\label{ied:par:eq:Bellshape}
 H(n,m)\sim(n-1)!\,2^{-n}m^n m^{-n/(3m)}I(n/m)
\end{equation}
uniformly when $n/m$ ranges over compact subsets of $(0,\infty)$.
Section~37, item~9, of \cite{Euler} explicitly asks whether
\begin{equation}\label{ied:par:eq:mainidentity}
 I(t)=t\,2^{t/3}h(t).
\end{equation}
The same question is recorded in the iterated-Bell report. Both sources
observe agreement of numerical values, but leave the orientation,
normalization and endpoint terms of a proposed contour comparison
unresolved.

\begin{theorem}[Amplitude identification]\label{ied:par:thm:identity}
Under the four analytic input statements recorded in
Section~\ref{ied:par:sec:inputs}, identity \eqref{ied:par:eq:mainidentity} holds for every
real $t>0$. The density therefore has the unique holomorphic extension
\begin{equation}\label{ied:par:eq:hhol}
 h(z)=\frac{2^{-z/3}}{z}I(z),\qquad\Re z>0.
\end{equation}
In particular, $h$ is real analytic on all of $(0,\infty)$.
\end{theorem}

The formulation records ordinary mathematical dependencies: the four
statements are proved in the cited repository reports. Those sources
are unrefereed and not proof-assistant verified. The present proof
checks the implication from their stated analytic interfaces; it does
not promote their earlier proofs or interval software to formal
verification.

Our second main result concerns the geometric aggregate
\[
 S_n(L)=\sum_{m\ge0}e^{-Lm}H(n,m),\qquad L>0.
\]
The sum converges for every $n$. The following theorem justifies the
missing depth summation and identifies its normalization.

\begin{theorem}[Geometric aggregation]\label{ied:par:thm:aggregate}
For every $L>0$,
\begin{equation}\label{ied:par:eq:Sasym}
 S_n(L)\sim
 (n!)^2(2L)^{-n}n^{-1-L/3}L^{L/3-1}I(L).
\end{equation}
The equivalent is uniform for $L$ in any fixed compact subinterval of
$(0,\infty)$.
\end{theorem}

For strict partition chains, put
\[
 Z_1(q)=1,\qquad
 Z_n(q)=q\sum_{j=1}^{n-1}
 \left\{\begin{matrix}n\\j\end{matrix}\right\}Z_j(q)
 \quad(n\ge2),
\]
where the braces denote Stirling numbers of the second kind. Thus $q$
marks the number of strict transitions, not the number of partitions
in the chain. Write
\begin{equation}\label{ied:par:eq:Lq}
 L(q)=\log(1+1/q)\qquad(q>0).
\end{equation}
\begin{corollary}[Marked chain amplitude]\label{ied:par:cor:marked}
For every $q>0$,
\begin{equation}\label{ied:par:eq:Zasym}
 Z_n(q)\sim C(q)(n!)^2[2L(q)]^{-n}n^{-1-L(q)/3},
\end{equation}
uniformly on compact positive $q$-intervals, with
\begin{align}
 C(q)&=\frac{L(q)^{L(q)/3-1}I(L(q))}{q+1}\label{ied:par:eq:CqI}\\
 &=\frac{[2L(q)]^{L(q)/3}h(L(q))}{q+1}.\label{ied:par:eq:Cqh}
\end{align}
This $C(q)$ is positive for real $q>0$ and extends holomorphically to
the half-plane $\Re q>0$.
\end{corollary}
At $q=1$, $Z_n(1)$ is OEIS A005121, whose normalized leading constant is
Lengyel's constant $C$. Thus
\begin{equation}\label{ied:par:eq:Lengyel}
 C=\frac12(\log2)^{(\log2)/3-1}I(\log2)
   =\frac12(2\log2)^{(\log2)/3}h(\log2).
\end{equation}
The leading equivalent and the existence of this constant are classical
\cite{Lengyel,Prellberg}; our contribution is the rigorous identification
with the precise functions and normalization used by the current
repository reports.

\subsection{What changes in the research record}
\begin{center}\small
\begin{tabular}{@{}p{.40\textwidth}p{.51\textwidth}@{}}
\toprule
Question or statement & Status in this companion\\
\midrule
Euler--Bell amplitude identity, \cite{Euler}, item 9
 & Proved by positive coefficient measures.\\
Real analyticity of $h$ on $(0,\infty)$
 & Proved; holomorphic extension to $\Re z>0$.\\
Strict-chain amplitude from the depth theorem
 & Proved with a uniform bound for every depth.\\
Marked amplitude beyond a neighborhood of $q=1$
 & Holomorphic on $\Re q>0$; positive on the positive axis.\\
Bell proportional-depth leading form
 & An input from \cite{IBD}, with older antecedents credited.\\
All fixed-order marked-chain expansions near $q=1$
 & Prior repository result \cite{Chains}; not reproved here.\\
Complex zero-freeness and global uniform complex expansions
 & Remain open; amplitude continuation alone does not imply them.\\
\bottomrule
\end{tabular}
\end{center}

\section{Four analytic inputs and exact conventions}\label{ied:par:sec:inputs}

We separate the new argument from the analytic foundations it uses.
This permits an audit without reconstructing the long Euler-transform
clock estimates or the iterated-Bell contour certificate.

\subsection*{Input F: normalized Fatou coordinate and orbit}
Lemma \file{ied:lead:lem:fatou} and equation
\file{ied:lead:eq:fatou} of \cite{Euler} give an analytic $\Phi$ on a
sufficiently small sector about the positive real axis with
\eqref{ied:par:eq:Phi}. Its local inverse $P_{\rm loc}$ extends to an entire
function $P$, agrees with that inverse in its sector, and satisfies
the functional equation in \eqref{ied:par:eq:P}. In particular, for all
sufficiently large integers $m$,
\[
 r_m=P(-m)>0,\qquad\Phi(r_m)=-m.
\]
Only this inverse statement near the negative real $s$-axis and the
entire functional equation will be needed.

\subsection*{Input M: positive Laplace measure}
Proposition \file{ied:lead:prop:measure} of \cite{Euler} proves that $P$
is the Laplace transform of a finite positive measure on $[0,\infty)$.
It has no atom at zero, every exponential moment, and a continuous
strictly positive density $h$ on $(0,\infty)$.
The identification needs only finiteness, positivity, the Laplace
formula and continuity. Strict positivity is used for later relative
aggregation. Smoothness, proved in the source's
Corollary \file{ied:lead:cor:smooth}, is not needed as an input.

\subsection*{Input L: absolute local coefficient transfer}
Equation \file{ibd:pd:eq:absolute} of \cite{IBD} states that, for
each fixed $0<a<b<\infty$, there is a finite $D$ such that
\begin{equation}\label{ied:par:eq:LLT}
 \frac{H(n,m)}
 {(n-1)!\,2^{-n}m^n m^{-n/(3m)}}
 =I(n/m)+\Oh\!\left(\frac{(1+\log m)^D}{m}\right)
\end{equation}
uniformly for integers $n,m\ge1$ with $a\le n/m\le b$.
Here $m\to\infty$, and the implicit constant may depend on $a,b$.
This is an \emph{absolute} estimate; it appears before the source's
positivity proof and involves no division by $I$.

\subsection*{Input A: holomorphy of the contour amplitude}
The first part of the proof of Proposition
\file{ibd:pd:prop:positive} of \cite{IBD} shows that $I$ is
holomorphic on $\Re t>0$. It uses a complex-linear contour functional
and locally uniform domination. The later numerical seed for
strict positivity is not needed for holomorphy, and is not used here.

\begin{remark}[Dependency and numerical boundaries]\label{ied:par:rem:seed}
The amplitude identification gives a route from positivity of $h$ to
positivity of $I$ without a numerical enclosure for $I(1)$. This does
not erase other computational dependencies of the earlier Bell
foundation. That report uses an interval certificate for common entry
on a fixed semicircle. We import its resulting absolute transfer
statement rather than independently formalizing that certificate.
\end{remark}

\subsection{Nonnegative coefficients}
Every coefficient of $f$ is nonnegative and $f(0)=0$, so every
coefficient of $f^{\circ m}$ is nonnegative. The exact recurrence is
\begin{equation}\label{ied:par:eq:Hrec}
 H(n,m+1)=\sum_{j=1}^{n}
 \left\{\begin{matrix}n\\j\end{matrix}\right\}H(j,m),
 \qquad H(n,0)=\mathbf1_{\{n=1\}}.
\end{equation}
Indeed, iterates commute, and
\[
 \frac{f(z)^j}{j!}=
 \sum_{n\ge j}\left\{\begin{matrix}n\\j\end{matrix}\right\}
 \frac{z^n}{n!}.
\]
These facts ensure positivity of the measures below. The use of
$(n-1)!$ rather than $n!$ in \eqref{ied:par:eq:LLT} is essential: after coefficient
extraction it supplies $1/n$, ultimately producing the factor $t$
in \eqref{ied:par:eq:mainidentity}.

\section{Positive coefficient measures and their limit}\label{ied:par:sec:measures}

\begin{lemma}[The exact radius and its expansion]\label{ied:par:lem:radius}
Let $r_m=P(-m)$ for sufficiently large integers $m$. Then
\begin{align}
 \frac2{r_m}
 &=m-\frac13\log m+\frac13\log2
 +\Oh\!\left(\frac{\log m}{m}\right),\label{ied:par:eq:radius}\\
 \log\frac{mr_m}{2}
 &=\frac{\log m-\log2}{3m}
 +\Oh\!\left(\frac{(\log m)^2}{m^2}\right).\label{ied:par:eq:radiuslog}
\end{align}
\end{lemma}
\begin{proof}
Put $u_m=2/r_m$. Since $\Phi(r_m)=-m$ and
$\log r_m=\log2-\log u_m$, Input F gives
\[
 u_m+\frac13\log u_m=m+\frac13\log2+\Oh(1/u_m).
\]
It first follows that $u_m/m\to1$: the logarithmic term is $o(u_m)$,
and the identity forces $u_m\to\infty$. Consequently
$u_m=m+\Oh(\log m)$ and $\log u_m=\log m+\Oh(\log m/m)$.
Substitution proves \eqref{ied:par:eq:radius}. Taylor expansion of
$-\log(1+(u_m-m)/m)$ proves \eqref{ied:par:eq:radiuslog}.
\end{proof}

For every sufficiently large integer $m$, define
\begin{equation}\label{ied:par:eq:mu}
 \mu_m=\sum_{n\ge1}
 \bigl(\coef n f^{\circ m}(z)\bigr)r_m^n\,\delta_{n/m}.
\end{equation}
This is a finite positive measure because $f^{\circ m}$ is entire.
The functional equation gives its exact normalization:
\begin{equation}\label{ied:par:eq:mass}
 \mu_m([0,\infty))=f^{\circ m}(r_m)
 =f^{\circ m}(P(-m))=P(0).
\end{equation}
Thus no asymptotic renormalization of these measures is needed.

\begin{lemma}[Laplace convergence]\label{ied:par:lem:laplace}
For each fixed real $s$,
\begin{equation}\label{ied:par:eq:Lapconv}
 \int_0^\infty e^{st}\dd\mu_m(t)
 =f^{\circ m}(r_m e^{s/m})\longrightarrow P(s).
\end{equation}
The convergence of the right-hand side is locally uniform for complex $s$.
\end{lemma}
\begin{proof}
For $s$ in a fixed compact subset of $\C$ and $m$ large,
$z_m=r_m e^{s/m}$ lies in the local inverse sector. Set
$s_m=m+\Phi(z_m)$. The local identity $P(\Phi(z_m))=z_m$ and
iteration of \eqref{ied:par:eq:P} give exactly
\[
 f^{\circ m}(z_m)=P(s_m).
\]
Subtracting the two coordinate expansions and writing $u_m=2/r_m$,
\begin{align*}
 s_m
 &=\Phi(r_m e^{s/m})-\Phi(r_m)\\
 &=u_m(1-e^{-s/m})+\frac{s}{3m}+\Oh(1/m)\\
 &=s+\Oh\!\left(\frac{1+\log m}{m}\right).
\end{align*}
The estimates are uniform on the chosen compact set. The last line
uses $u_m=m+\Oh(\log m)$. Entire analyticity of $P$ proves convergence.
The Laplace equality follows from the positive power series in
\eqref{ied:par:eq:mu}.
\end{proof}

\begin{lemma}[Compact test convergence]\label{ied:par:lem:weak}
For every continuous function $\varphi$ with compact support in
$(0,\infty)$,
\begin{equation}\label{ied:par:eq:weak}
 \int\varphi(t)\dd\mu_m(t)\longrightarrow
 \int_0^\infty\varphi(t)h(t)\dd t.
\end{equation}
\end{lemma}
\begin{proof}
Push $\mu_m$ and $h(t)\dd t$ forward under $x=e^{-t}$ to finite
measures on $[0,1]$, giving them zero mass at $x=0$.
Their total masses are $P(0)$. For every integer $k\ge0$,
Lemma~\ref{ied:par:lem:laplace} at $s=-k$ gives convergence of their $k$th
moments. Hence their integrals converge for each polynomial.

Uniform polynomial approximation on $[0,1]$ gives convergence for
every continuous function there: an approximation error at most
$\epsilon$ contributes at most $2P(0)\epsilon$, uniformly in $m$.
The function
\[
 g(x)=\varphi(-\log x)\quad(0<x\le1),\qquad g(0)=0
\]
is continuous on $[0,1]$ because $\varphi$ has compact support inside
$(0,\infty)$. Applying the conclusion to $g$ proves \eqref{ied:par:eq:weak}.
\end{proof}

\begin{remark}
A compact asymptotic formula for individual coefficients alone need not
identify a global transform limit. Here the masses are finite and
fixed, and their Laplace transforms determine polynomial moments after
compactification. The same measures can then be tested in the local
coefficient regime. No positive-measure interpretation of the signed
Bell contour functional is required.
\end{remark}

\section{Identification, analyticity and endpoint consequences}\label{ied:par:sec:identification}

\begin{proof}[Proof of Theorem~\ref{ied:par:thm:identity}]
Fix $0<a<b<\infty$ and put $t=n/m$. Dividing \eqref{ied:par:eq:LLT} by $n!$
and multiplying by $r_m^n$ gives, uniformly for $a\le t\le b$,
\begin{equation}\label{ied:par:eq:atomfirst}
 \bigl(\coef n f^{\circ m}(z)\bigr)r_m^n
 =\frac1n\left(\frac{mr_m}{2}\right)^n m^{-t/3}
 \left[I(t)+\Oh\!\left(\frac{(1+\log m)^D}{m}\right)\right].
\end{equation}
By Lemma~\ref{ied:par:lem:radius},
\[
 n\log\frac{mr_m}{2}-\frac t3\log m
 =-\frac t3\log2+
 \Oh\!\left(\frac{(1+\log m)^2}{m}\right).
\]
Therefore
\begin{equation}\label{ied:par:eq:atom}
 \bigl(\coef n f^{\circ m}(z)\bigr)r_m^n
 =\frac1m\left[
 \frac{2^{-t/3}I(t)}t+
 \Oh\!\left(\frac{(1+\log m)^K}{m}\right)\right],
\end{equation}
where $K=\max\{D,2\}$ suffices. Only boundedness of $I$ on $[a,b]$
and $t\ge a$ have been used, not positivity of $I$.

Let $\varphi$ have compact support in $(a,b)$. Summing
\eqref{ied:par:eq:atom} against $\varphi(n/m)$ is a Riemann sum.
There are $\Oh(m)$ relevant atoms, so their accumulated error is
$\Oh((1+\log m)^K/m)=o(1)$. Thus
\[
 \int\varphi\dd\mu_m\longrightarrow
 \int_a^b\varphi(t)\frac{2^{-t/3}I(t)}t\dd t.
\]
Lemma~\ref{ied:par:lem:weak} gives another expression for the same limit.
The continuous function $h(t)-2^{-t/3}I(t)/t$ has zero integral against
every such test function and therefore vanishes pointwise.
Arbitrariness of $a,b$ proves \eqref{ied:par:eq:mainidentity} on $(0,\infty)$.

Input A makes the right-hand side of \eqref{ied:par:eq:hhol} holomorphic
on $\Re z>0$. It agrees with $h$ on the positive axis.
Uniqueness follows from the identity theorem.
\end{proof}

There are three sensitive normalizations. The denominator in
\eqref{ied:par:eq:radius} is $m-\frac13\log m+\frac13\log2$.
The factor $1/n$ in \eqref{ied:par:eq:atomfirst} comes from $(n-1)!/n!$.
The exponential tilt cancels $m^{-t/3}$ but leaves $2^{-t/3}$.
Their combination gives exactly the stated identity.

Equality was established on each positive interval by test functions.
It was not inferred from matching finitely many endpoint coefficients
or decimal values. The proof also never takes $n/m\to0$ in a theorem
whose hypotheses require a fixed positive compact ratio range.

\begin{corollary}[Analytic positivity]\label{ied:par:cor:positive}
$I(t)>0$ for every $t>0$.
\end{corollary}
\begin{proof}
Input M gives $h(t)>0$, and \eqref{ied:par:eq:mainidentity} has a positive multiplier.
\end{proof}

\subsection{The small-parameter amplitude}
The Euler report supplies a separate endpoint input: for every fixed
$N\ge2$,
\begin{equation}\label{ied:par:eq:hendpoint}
 \log\frac{h(t)}2
 =-\frac t3\log t+\frac{1-\gamma-\log2}{3}t
 +\sum_{j=2}^N c_jt^j+\Oh(t^{N+1}),\qquad t\downarrow0,
\end{equation}
with fixed differentiated remainders and $c_2=(9-\pi^2)/108$.
This is equation \file{ied:log:eq:logdensity} and its coefficient
discussion in \cite{Euler}; it is not required for the main theorem.

\begin{corollary}[Endpoint transseries for $I$]\label{ied:par:cor:endpoint}
For every fixed $N\ge2$,
\begin{equation}\label{ied:par:eq:Iendpoint}
 \log\frac{I(t)}{2t}
 =-\frac t3\log t+\frac{1-\gamma}{3}t
 +\sum_{j=2}^N c_jt^j+\Oh(t^{N+1}).
\end{equation}
The fixed differentiated remainder statements transfer unchanged.
In particular,
\begin{align}
 I(t)=2t\bigg\{1
 &+\frac t3(1-\gamma-\log t)\notag\\
 &+t^2\left[
 \frac{(1-\gamma-\log t)^2}{18}
 +\frac{9-\pi^2}{108}\right]
 +\Oh\bigl(t^3(1+|\log t|)^3\bigr)\bigg\}.
 \label{ied:par:eq:Iendfirst}
\end{align}
\end{corollary}
\begin{proof}
Take logarithms of \eqref{ied:par:eq:mainidentity} and substitute
\eqref{ied:par:eq:hendpoint}. The $\log2$ term cancels exactly.
Exponentiation through degree two gives \eqref{ied:par:eq:Iendfirst}.
The logarithms differ by an exactly linear function, so the
differentiated remainders introduce no new error.
\end{proof}

This solves the amplitude part of the Bell large-depth endpoint
question. Coefficient asymptotics when depths grow like $n\log n$
still need a uniform transfer as $n/m$ decreases. Neither the endpoint
series nor holomorphic continuation supplies that uniformity.
The smooth renormalized germ at zero is not shown here to have a
convergent Taylor series.

\subsection{A reusable identification principle}
\begin{proposition}[General positive-measure matching]\label{ied:par:prop:general}
Suppose a positive-coefficient iterated family has an entire orbit $P$
and local inverse coordinate
\[
 \Phi(w)=-\frac1{aw}+\rho\log w+\Oh(w),
 \qquad a>0,\quad\rho\in\R,
\]
with the analogues of Inputs F and M. Suppose its absolute local transfer
for $t=n/m$ in positive compact sets is
\[
 \coef n f^{\circ m}(z)
 =\frac1n(am)^n m^{-\rho t}\bigl(J(t)+o(1)\bigr),
\]
where $J$ is continuous and the error uniform. Then
\[
 J(t)=t\,a^{-\rho t}h(t)\qquad(t>0).
\]
\end{proposition}
\begin{proof}
Set $r_m=P(-m)$ and $u_m=1/(ar_m)$.
The inverse identity yields
$u_m=m-\rho\log m-\rho\log a+o(1)$, and therefore
\[
 m^{-\rho t}(amr_m)^n\longrightarrow a^{\rho t}
\]
uniformly on positive compact ratio ranges.
The Laplace convergence of the coefficient measures follows as before,
since $u_m(1-e^{-s/m})\to s$.
Their local density is $a^{\rho t}J(t)/t$, which compact test
convergence identifies with $h$.
\end{proof}
This is a transfer principle with stated hypotheses, not a claim that
every parabolic map has a positive Laplace orbit or a local coefficient
theorem. The exponential map has $a=1/2$ and $\rho=1/3$.

\section{A bound valid for every iteration depth}\label{ied:par:sec:majorant}

For power series with real coefficients, write $A\preceq B$ when every
coefficient of $B-A$ is nonnegative. If the series have nonnegative
coefficients and zero constant term, this order is preserved by
composition.

\begin{lemma}[Rational composition majorant]\label{ied:par:lem:majorant}
For every integer $m\ge0$,
\begin{equation}\label{ied:par:eq:majorantseries}
 f^{\circ m}(z)\preceq\frac{z}{1-mz/2}.
\end{equation}
In particular, for $m,n\ge1$,
\begin{equation}\label{ied:par:eq:Hmajorant}
 0\le H(n,m)\le n!\left(\frac m2\right)^{n-1}.
\end{equation}
\end{lemma}
\begin{proof}
For every $j\ge1$, $j!\ge2^{j-1}$: each factor $2,\ldots,j$ is at least
$2$. Hence
\[
 f(z)=\sum_{j\ge1}\frac{z^j}{j!}\preceq
 g(z):=\frac{z}{1-z/2}.
\]
Composition and induction give $f^{\circ m}\preceq g^{\circ m}$.
Direct substitution shows $g^{\circ m}(z)=z/(1-mz/2)$.
Extract coefficients and multiply by $n!$.
\end{proof}

This elementary inequality supplies what compact-ratio asymptotics
omit: a bound at depths much smaller or larger than $n$.
Near the main saddle it loses a factor of order $n^{L/3}$, but
exponential tail decay dominates that polynomial loss.
We give the needed summation estimates explicitly.

\begin{lemma}[Discrete gamma sums]\label{ied:par:lem:gamma}
Let $L$ range over a fixed compact interval
$[L_0,L_1]\subset(0,\infty)$. Uniformly as $n\to\infty$,
\begin{equation}\label{ied:par:eq:gammasum}
 \sum_{m\ge1}e^{-Lm}m^n
 =\frac{\Gamma(n+1)}{L^{n+1}}
 \bigl(1+\Oh(n^{-1/2})\bigr).
\end{equation}
Put $w_n=\sqrt n(\log n)^2$. For $p=n$ or $p=n-1$, there are constants
$c,C>0$ such that
\begin{equation}\label{ied:par:eq:gammatail}
 \sum_{\substack{m\ge1\\|m-n/L|>w_n}}e^{-Lm}m^p
 \le C\frac{\Gamma(p+1)}{L^{p+1}}e^{-c(\log n)^4}
\end{equation}
for all sufficiently large $n$, uniformly in $L$.
\end{lemma}
\begin{proof}
For a nonnegative unimodal integrable function $v$ on $[0,\infty)$,
comparison with unit intervals on each side of its maximum gives
\[
 \left|\sum_{m\ge1}v(m)-\int_0^\infty v(x)\dd x\right|
 \le2\sup_{x\ge0}v(x).
\]
Take $v(x)=e^{-Lx}x^p$. Its integral is
$\Gamma(p+1)L^{-p-1}$ and its maximum is $(p/(eL))^p$.
Stirling's formula shows that their ratio is $\Oh(p^{-1/2})$
uniformly on the stated $L$-interval. This proves
\eqref{ied:par:eq:gammasum} and bounds the complete sum by a uniform constant
times its integral.

For the upper tail, set $M=n/L+w_n$ and
$\theta=L^2w_n/(4n)$. For large $n$, $L\pm\theta$ lie in a fixed
positive compact interval. Exponential tilting gives
\begin{align*}
 \sum_{m\ge M}e^{-Lm}m^p
 &\le e^{-\theta M}\sum_{m\ge1}e^{-(L-\theta)m}m^p\\
 &\le C\Gamma(p+1)L^{-p-1}
 e^{-\theta M}\left(\frac{L}{L-\theta}\right)^{p+1}.
\end{align*}
Let $x=\theta/L$. Because $p+1=n+\Oh(1)$ and
$-\log(1-x)\le x+x^2$ for $0\le x\le1/2$, the logarithm of the last
factor is at most
\[
 -\theta w_n+\Oh(\theta)+(n+\Oh(1))\theta^2/L^2.
\]
Its principal terms are
$-L^2w_n^2/(4n)+L^2w_n^2/(16n)$.
Thus it is bounded above by $-c w_n^2/n$ uniformly for large $n$.

For the lower tail, when $m\le n/L-w_n$,
\[
 e^{-Lm}\le e^{\theta(n/L-w_n)}e^{-(L+\theta)m}.
\]
The same complete-sum bound leads to the logarithm
\[
 \theta(n/L-w_n)-(p+1)\log(1+\theta/L).
\]
Using $\log(1+x)\ge x-x^2/2$ gives the same negative bound.
Finally $w_n^2/n=(\log n)^4$.
Rounding the real cutoffs to integers only introduces factors
$e^{\Oh(\theta)}$, which do not affect the estimates.
\end{proof}

\section{Summing the Bell local limit}\label{ied:par:sec:aggregation}

\begin{proof}[Proof of Theorem~\ref{ied:par:thm:aggregate}]
Fix a compact positive $L$-interval and use
\[
 \mathcal W_n(L)=
 \{m\in\N:|m-n/L|\le w_n\},
 \qquad w_n=\sqrt n(\log n)^2.
\]
Throughout this window, uniformly,
\[
 \frac nm=L+\Oh\!\left(\frac{(\log n)^2}{\sqrt n}\right),
 \qquad
 \log m=\log(n/L)+
 \Oh\!\left(\frac{(\log n)^2}{\sqrt n}\right).
\]
By the amplitude identification, $I$ is continuous and strictly
positive. It has a positive lower bound on the compact range of these
ratios. Therefore
\begin{equation}\label{ied:par:eq:slow}
 m^{-n/(3m)}I(n/m)
 =(n/L)^{-L/3}I(L)(1+o(1))
\end{equation}
uniformly on the window and in $L$. In particular, the difference
between the logarithm of the power on the left and its central value is
$\Oh((\log n)^3/\sqrt n)=o(1)$. Its dependence on $n$ does not
obstruct this replacement.

Input L now applies within one fixed compact ratio range and gives
\begin{align}
 \sum_{m\in\mathcal W_n(L)}e^{-Lm}H(n,m)
 &=(n-1)!\,2^{-n}(n/L)^{-L/3}I(L)(1+o(1))\notag\\
 &\quad\times\sum_{m\in\mathcal W_n(L)}e^{-Lm}m^n.
 \label{ied:par:eq:central}
\end{align}
The absolute transfer error is a relative $o(1)$ here because
$m\asymp n$ and $I$ has a positive lower bound.
Lemma~\ref{ied:par:lem:gamma} shows that the last sum equals
$n!L^{-n-1}(1+o(1))$.
Substitution gives the main term in \eqref{ied:par:eq:Sasym}.

For the omitted depths, both preceding lemmas give
\begin{align*}
 \sum_{m\notin\mathcal W_n(L)}e^{-Lm}H(n,m)
 &\le n!\,2^{1-n}
 \sum_{m\notin\mathcal W_n(L)}e^{-Lm}m^{n-1}\\
 &\le Cn!\,2^{1-n}(n-1)!L^{-n}e^{-c(\log n)^4}.
\end{align*}
Dividing by the main scale gives at most a uniform constant times
$n^{L_1/3}e^{-c(\log n)^4}\to0$.
The positive amplitude $L^{L/3-1}I(L)$ is bounded away from zero on
the fixed interval. The $m=0$ term is zero for $n>1$.
All estimates were uniform, proving the theorem.
\end{proof}

\subsection{An exact mixture for strict chains}
\begin{lemma}\label{ied:par:lem:mixture}
For $q>0$ and $a=q/(q+1)$, for every $n\ge1$,
\begin{equation}\label{ied:par:eq:mixture}
 Z_n(q)=(1-a)\sum_{m\ge0}a^mH(n,m)
       =\frac1{q+1}S_n(L(q)).
\end{equation}
\end{lemma}
\begin{proof}
The majorant proves absolute convergence.
Put $T_n=\sum_{m\ge0}a^mH(n,m)$ and sum \eqref{ied:par:eq:Hrec}:
\[
 T_n=\mathbf1_{\{n=1\}}+
 a\sum_{j=1}^{n}
 \left\{\begin{matrix}n\\j\end{matrix}\right\}T_j.
\]
The diagonal Stirling number equals $1$, so $T_1=(1-a)^{-1}$ and
\[
 T_n=\frac{a}{1-a}\sum_{j<n}
 \left\{\begin{matrix}n\\j\end{matrix}\right\}T_j
 \quad(n\ge2).
\]
Multiplying by $1-a$ and using $a/(1-a)=q$ gives the defining recurrence
and initial value of $Z_n(q)$. Uniqueness by induction proves the result.
\end{proof}

\begin{proof}[Proof of Corollary~\ref{ied:par:cor:marked}]
Apply the aggregation theorem to \eqref{ied:par:eq:mixture} and then the
amplitude identification. This proves \eqref{ied:par:eq:Zasym},
\eqref{ied:par:eq:CqI} and \eqref{ied:par:eq:Cqh}, their positivity and compact real
uniformity.

For the complex continuation define
$L(q)=\Log(1+1/q)$ on $\Re q>0$, using the principal logarithm.
There $\Re(1+1/q)>1$, so $L$ is holomorphic and
$\Re L(q)=\log|1+1/q|>0$.
In particular $L$ is nonzero and $\Log L$ is holomorphic as well.
The function
\[
 \frac1{q+1}
 \exp\!\left[\left(\frac{L(q)}3-1\right)\Log L(q)\right]I(L(q))
\]
is holomorphic throughout $\Re q>0$ and agrees with the identified
positive-real amplitude. On a common neighborhood of $q=1$ it agrees,
by the identity theorem, with the prior local marked amplitude
constructed in \cite{Chains}.
\end{proof}

The continuation concerns the amplitude. A uniform expansion for
$Z_n(q)$ on a complex region additionally needs control of cancellation
and of its remainders; it does not follow from amplitude continuation.

\section{Consequences for chain-length constants}\label{ied:par:sec:constants}

The prior local marked-chain expansion \cite{Chains} gives, with
$L=\log2$ and $\mathcal H_n$ the number of strict transitions in a
uniformly chosen chain,
\begin{align}
 \mathbb E\,\mathcal H_n
 &=\frac n{2L}+\frac{\log n}{6}+K_1-\frac1{12n}
 +\Oh(n^{-2}),\label{ied:par:eq:mean}\\
 \operatorname{Var}(\mathcal H_n)
 &=\frac{1-L}{4L^2}n-\frac{\log n}{12}+K_2+\frac1{24n}
 +\Oh(n^{-2}).\label{ied:par:eq:variance}
\end{align}
That source defines the constant terms by
\[
 K_1=\left.\frac{\dd}{\dd t}\log C(e^t)\right|_{t=0},
 \qquad
 K_2=\left.\frac{\dd^2}{\dd t^2}\log C(e^t)\right|_{t=0}.
\]
The new identity expresses these in the common density and its
derivatives.

\begin{corollary}[Density formulas for constant cumulant terms]
\label{ied:par:cor:K}
Put $A(x)=\log h(x)+(x/3)\log(2x)$ for $x>0$. Then
\begin{equation}\label{ied:par:eq:Kcompact}
 K_1=-\frac12-\frac12A'(L),\qquad
 K_2=-\frac14+\frac14\bigl(A''(L)+A'(L)\bigr).
\end{equation}
Equivalently,
\begin{align}
 K_1={}&-\frac12-\frac12\left[
 \frac{h'(L)}{h(L)}+\frac{\log(2L)+1}{3}\right],
 \label{ied:par:eq:K1}\\
 K_2={}&-\frac14+\frac14\left[
 \frac{h''(L)}{h(L)}-\left(\frac{h'(L)}{h(L)}\right)^2
 +\frac{h'(L)}{h(L)}
 +\frac1{3L}+\frac{\log(2L)+1}{3}\right].
 \label{ied:par:eq:K2}
\end{align}
\end{corollary}
\begin{proof}
Near $t=0$, \eqref{ied:par:eq:Cqh} gives
\[
 \log C(e^t)=-\log(1+e^t)+
 A\bigl(\log(1+e^{-t})\bigr).
\]
The inner function has first derivative $-1/2$ and second derivative
$1/4$ at zero; the first two derivatives of $-\log(1+e^t)$ there are
$-1/2$ and $-1/4$. The chain rule yields \eqref{ied:par:eq:Kcompact}.
Differentiating $A$ gives the expanded expressions.
\end{proof}

The mean, variance and central limit theorem are not new claims of
this note. These formulas identify their amplitude derivatives in the
normalization used by the three reports.
Counting partitions instead of strict transitions would add one to
the mean and to $K_1$; we retain the source's convention.

\section{Computation and reproducibility}\label{ied:par:sec:computation}

No analytic proof above uses a measured constant or a finite experiment.
The accompanying programs check exact normalizations and compare
independent numerical constructions.

\subsection{Exact finite checks}
The script \file{code/verify_exact.py} uses exact integer and rational
arithmetic for:
\begin{itemize}[leftmargin=*,itemsep=2pt]
\item $960$ instances of \eqref{ied:par:eq:Hmajorant},
for $1\le n\le24$, $1\le m\le40$;
\item $144$ evaluations of the depth polynomial in the Newton basis;
\item $96$ rational instances of the geometric mixture,
at $q=1/2,1,2,3/2$;
\item a symbolic check that the four-term Fatou coordinate has zero
Abel defect through degree five.
\end{itemize}
All checks passed. The Newton-basis check uses that the Stirling matrix
on degrees at most $n$ is $I+N$ with $N$ strictly triangular:
\[
 H(n,m)=\sum_{k=0}^{n-1}\binom mk(N^k)_{n1}.
\]
The coefficient $(N^k)_{n1}$ counts strict chains of length $k$,
independently checking the convention in \eqref{ied:par:eq:mixture}.
These verifications concern finite algebra, not bounds for asymptotic
remainders.

\subsection{Independent Laplace inversion}
The script \file{code/explore_density.py} approximates $P(-p)$ directly.
It solves a truncated inverse-coordinate equation far along the
negative direction and iterates $w\mapsto e^w-1$ forward.
The four coefficients are
\[
 \Phi(w)=-\frac2w+\frac13\log w-\frac w{36}
 +\frac{w^2}{540}+\frac{w^3}{7776}
 -\frac{71w^4}{435456}+\Oh(w^5).
\]
The exact script checks those coefficients symbolically.
Numerical inverse Laplace transformation then estimates $h(t)$.
This route neither uses the Bell fixed-contour amplitude nor estimates
the density from Bell counts.

Two runs used $160$ and $320$ forward steps and inversion degrees
$48$ and $64$. Their values agree to at least twelve decimal places
at the four displayed arguments. Because no interval bound for
truncation or quadrature is supplied, that agreement is a diagnostic,
not a certified precision claim.
\begin{center}
\begin{tabular}{@{}rrr@{}}
\toprule
$t$ & Exploratory $h(t)$ & $t2^{t/3}h(t)$\\
\midrule
$1/2$ & $2.14193747174$ & $1.20212176094$\\
$\log2$ & $2.03764226105$ & $1.65770061282$\\
$1$ & $1.81460955983$ & $2.28626478178$\\
$2$ & $1.02595377154$ & $3.25720019242$\\
\bottomrule
\end{tabular}
\end{center}

The script \file{code/compare_counts.py} computes Bell counts by exact
Stirling recurrence through degree $160$ and depth $320$.
It compares normalized counts at $n/m=1/2,1,2$ with the independently
inverted density. It also computes strict-chain counts through $n=320$
and compares their normalized constant, after the prior report's four
known corrections, with \eqref{ied:par:eq:Lengyel}.
The complete outputs are in \file{data/count_comparison.json}.
Finite agreement is never used as a proof premise.

\subsection{Replay}
From the companion directory run:
\begin{verbatim}
python code/verify_exact.py --output data/exact_checks.json
python code/explore_density.py --steps 160 --degree 48 \
  --output data/density_160.json
python code/explore_density.py --steps 320 --degree 64 \
  --output data/density_320.json
python code/compare_counts.py --density data/density_320.json \
  --output data/count_comparison.json
pdflatex -halt-on-error article.tex
pdflatex -halt-on-error article.tex
\end{verbatim}
Counting uses the Python standard library; the symbolic coordinate
check uses SymPy, and the optional decimal calculations use mpmath.
No script accesses the network. Outputs distinguish exact checks
from floating diagnostics.
\begin{writenote}[The package here, 9 October 2026]
In this report the three programs are \file{code/04-parabolic-verify_exact.py},
\file{code/04-parabolic-explore_density.py} and
\file{code/04-parabolic-compare_counts.py}, and the data files named above are
\file{data/04-parabolic-*}; they read and write the delivered names, so the
replay runs only on a copy with the delivered layout (the README gives the
recipe). The write did not rerun them. Its own exact checks of the majorant,
the mixture and the Newton basis, and its numerical check of the aggregate, are
recorded in Section~\ref{ied:par:sec:checks}; the table's
$t2^{t/3}h(t)$ column is the inverse-Laplace $h$ times an exact factor, so its
digits are as uncertified as those of $h$.
\end{writenote}

\section{Literature, provenance and research boundaries}\label{ied:par:sec:literature}

\subsection{Prior analytic work}
Prellberg's 2002 seminar, summarized by Mishna in 2005
\cite{Prellberg}, applies parabolic iteration and complex analysis to
linear composition recurrences. Its general setting is
\[
 X(z)=a(z)X(f(z))+b(z),\qquad
 f(z)=z+cz^2+dz^3+\cdots.
\]
The summary contains a two-index coefficient expression before
summing over depth and a final partition-chain equivalent with an
inverse-coordinate contour amplitude. For $f=e^z-1$, the quantity
$1-d/c^2$ is $1/3$. These are direct antecedents. The leading scale
and the use of Fatou coordinates are not discoveries of this note.

The four-page primary seminar summary was opened and read for this
work. The separate FPSAC slide URL cited by the Bell report could not
be retrieved in this session and is not represented as independently
read. The repository distinguishes its detailed compact-uniform
transfer theorem from the older outline.

Lengyel's 1984 recurrence paper is the historical source identified
by the summary \cite{Lengyel}. The fixed-degree, large-iteration
leading coefficient of the Bell array is also discussed by Skau and
Kristensen \cite{Skau}; its abstract was inspected here. That
information alone does not justify a stronger claim about the full
scope of the paper. Our arguments do not use it as a theorem.

\subsection{Repository pin and exact questions}
The inspected ProveIt revision is
\[
 \Repin.
\]
The source reports are under
\begin{center}\small
\file{SetTheory/Cardinals/docs/reports/}\\
\file{generating-functions-and-asymptotics/oeis-sequence-asymptotics/}.
\end{center}
Their exact source identifiers are:
\begin{center}\small
\begin{tabular}{@{}p{.47\textwidth}p{.45\textwidth}@{}}
\toprule
Report & Git blob of \file{article.tex}\\
\midrule
\file{a290354-iterated-euler-diagonals}
 & \nolinkurl{b5af1ca30b21cd89d90488258d8722aa8a155b0e}\\
\file{a139383-iterated-bell-diagonals}
 & \nolinkurl{2b1d8fbe811b11925a0879552b7319500f16cfd2}\\
\file{a005121-strict-partition-chains}
 & \nolinkurl{da8cac4780378f0a0a1af35df4ffebdee4b014d2}\\
\bottomrule
\end{tabular}
\end{center}

The precise open-question labels resolved here are
\file{ied:q:ibd}, the all-positive-axis portion of
\file{ied:q:analytic}, and the leading-amplitude portion of
\file{spc:q:depth}. The marked continuation also bears on
\file{spc:q:disk}, but does not settle the complex zero set or global
uniform asymptotic estimates. The mean-constant identities identify
amplitude derivatives requested beside the depth question.
\begin{writenote}[The pin and the labels, 9 October 2026]
At the pin the three blobs are as stated, and the reports' texts at HEAD are
the same, except that this report now has Part~IV. The labels named in the
preceding paragraph exist with the meanings given:
\file{ied:q:ibd} and \file{ied:q:analytic} are items~\ref{ied:q:ibd}
and~\ref{ied:q:analytic} of Section~\ref{ied:sec:further}, and \file{spc:q:depth},
\file{spc:q:disk} are items of the strict partition-chain report's
``Further questions and research''. What this Part answers in each is stated in
Section~\ref{ied:par:sec:answers}; ``the leading-amplitude portion of
\file{spc:q:depth}'' is that report's heuristic
$C\overset{?}{=}L^{L/3-1}I(L)/2$, which \eqref{ied:par:eq:Lengyel} proves.
\end{writenote}

The new work proves connections between existing normalizations and
supplies an explicit tail argument. Its repository-relative novelty
is verifiable at the pin. A literature search and source inspection do
not establish that no equivalent theorem has appeared elsewhere, and
this note makes no such claim. The connection is useful because
previously separate amplitude functions become one positive analytic
density, allowing parameter transfer without guessing contour constants.

\section{Further research questions}\label{ied:par:sec:further}
The following problems are concrete continuations of this work.
\begin{enumerate}[leftmargin=*,itemsep=5pt]
\item\label{ied:par:q:uniform} \textbf{Uniform Bell transfer toward the endpoint.}
Prove a version of \eqref{ied:par:eq:LLT} for $n/m$ decreasing on logarithmic
and then polynomial scales. The endpoint amplitude is now known.
The remaining task is a uniform coefficient remainder, especially
for the long-delayed parts of the fixed cut.

\item\label{ied:par:q:summation} \textbf{Summation of every correction.}
The Bell expansion has polynomials in $\log n$ at each inverse power
of $n$, whereas the chain recurrence expansion has no such logarithms.
The factor $m^{-n/(3m)}$ slightly displaces the central depth saddle.
Carry its expansion to every fixed order and prove cancellation of
its logarithms against the Bell corrections.
Theorem~\ref{ied:par:thm:aggregate} proves only the leading summation.

\item\label{ied:par:q:zeros} \textbf{Complex zeros of the marked amplitude.}
Determine whether \eqref{ied:par:eq:CqI} has zeros on $\Re q>0$; equivalently,
study zeros of $I$ or $h$ on the image of $L(q)$.
Real positivity gives no complex zero-free theorem.

\item\label{ied:par:q:positive} \textbf{A common positive-parameter asymptotic theory.}
Extend the recurrence proof of every fixed order from a small disk
about $q=1$ to neighborhoods of every positive compact interval.
Compare its limiting constant directly with \eqref{ied:par:eq:Cqh}.
Continuation of the amplitude alone is insufficient.

\item\label{ied:par:q:certified} \textbf{Certified density and cumulant constants.}
Give interval enclosures for $h(\log2)$, $h'(\log2)$ and
$h''(\log2)$ from the normalized orbit and explicit inversion tails.
Equations \eqref{ied:par:eq:Lengyel}, \eqref{ied:par:eq:K1} and \eqref{ied:par:eq:K2} then
yield certified $C,K_1,K_2$. The present decimal calculations stop
before that certification.

\item\label{ied:par:q:large} \textbf{Large-argument density asymptotics.}
Determine $h(t)$ as $t\to\infty$. This also determines the Bell
amplitude at small proportional depth and, through $L(q)\to\infty$
as $q\downarrow0$, the marked-chain amplitude at small transition
weight. Compact-parameter results do not control this endpoint.

\item\label{ied:par:q:germ} \textbf{Convergence or summability of the endpoint germ.}
The renormalization $t^{t/3}h(t)$ has finite-order Taylor asymptotics.
Decide convergence, or determine coefficient growth and summability.
Holomorphy on $\Re t>0$ does not decide behavior at its boundary
point zero.

\item\label{ied:par:q:other} \textbf{Other positive parabolic maps.}
Apply Proposition~\ref{ied:par:prop:general} to a family for which both the
positive Laplace orbit and local coefficient theorem can be proved.
The main difficulty is establishing those hypotheses, rather than
identifying the measures once they hold.

\item\label{ied:par:q:formal} \textbf{Formal verification of the bridge.}
Formalize the implication from the four explicit analytic inputs:
radius asymptotics, finite positive measures, polynomial approximation
and Riemann sums. This gives a small auditable theorem boundary before
the much larger earlier contour foundations are formalized.
\end{enumerate}
\begin{writenote}[Further questions of Part~IV, 9 October 2026]
Following Vladimir Reshetnikov's standing rule of 4~October 2026, every claim
of the manuscript that is not proved in full is stated as an open question;
these nine items are the manuscript's own, as delivered, and the write labelled
them. No claim of the manuscript was found wrong; nothing is moved or refuted.
Relations to Section~\ref{ied:sec:further} and to the neighbouring reports:
Question~\ref{ied:par:q:uniform} is the coefficient half of the iterated Bell
report's endpoint $\lambda\to\infty$ (the amplitude half is
Corollary~\ref{ied:par:cor:endpoint}) and hypothesis~(b) of
Remark~\ref{ied:rem:labelled}; Question~\ref{ied:par:q:summation} is the
open part of the chain report's \file{spc:q:depth} (the write's numerical
observation at the end of Section~\ref{ied:sec:further} suggests a first
correction $L/(6n)$ without logarithm); Questions~\ref{ied:par:q:zeros}
and~\ref{ied:par:q:positive} are the open part of \file{spc:q:disk};
Question~\ref{ied:par:q:certified} contains item~\ref{ied:q:validated} at
$\log2$; Question~\ref{ied:par:q:large} is the $\beta\to\infty$ half of
item~\ref{ied:q:endpoint}; Question~\ref{ied:par:q:germ} is the open part of
item~\ref{ied:q:analytic}.
\end{writenote}

\appendix
\section{A short proof checklist for independent review}\label{ied:par:app:checklist}
A reviewer can assess the new reasoning by checking this chain:
\begin{enumerate}[leftmargin=*,itemsep=3pt]
\item $r_m=P(-m)$ is positive and gives the exact mass
$f^{\circ m}(r_m)=P(0)$.
\item The additive constant of $\Phi$ in \eqref{ied:par:eq:Phi} is zero.
It gives the $\log2/3$ term in \eqref{ied:par:eq:radius}.
\item The inverse identity is applied only near $r_m$, where
$r_m e^{s/m}$ stays in the stated sector.
\item Negative integer Laplace values and equal finite masses give
compact test convergence after $x=e^{-t}$.
\item Input L is absolute and uniform on each fixed positive compact
interval; such an interval contains $\Oh(m)$ atoms.
\item Continuity converts equality of the tested densities into a
pointwise identity.
\item The central aggregation window lies in one compact ratio
region; all outside depths obey the rational majorant.
\item The majorant's loss $n^{L/3}$ is retained before its tail is
declared negligible.
\item For complex continuation, $L(q)$ stays in $\Re L>0$ on
$\Re q>0$, with every logarithm branch specified.
\end{enumerate}
None of these steps uses the decimal values.
The earlier analytic interfaces remain the inherited dependency.

\begin{thebibliography}{99}
\bibitem{A290353} The OEIS Foundation Inc., \emph{A290353}, square array of iterated Euler transforms, entry by Alois P. Heinz, 28 July 2017. Retrieved 5 October 2026. \url{https://oeis.org/A290353}.\small
\\{}[As cited in Report~228: \emph{A290353}, iterated Euler-transform array, entry by Alois P. Heinz, 28 July 2017. Source checked 5 October 2026. Report~226: Source checked for Report 225 on 5 October 2026.]
\bibitem{A290354} The OEIS Foundation Inc., \emph{A290354}, entry by Alois P. Heinz, 28 July 2017; asymptotic conjecture by V\'aclav Kot\v{e}\v{s}ovec, 14 August 2017. Retrieved 5 October 2026. \url{https://oeis.org/A290354}.
\\{}[As cited in Report~226: leading asymptotic conjecture by V\'aclav Kot\v{e}\v{s}ovec, 14 August 2017. Source checked for Report 225 on 5 October 2026.]
\bibitem{Kaneiwa} R. Kaneiwa, \emph{On r-fold Partitions and a Certain Form of Infinite Products}, Otaru University of Commerce, Humanities Research \textbf{61} (1980), 11--26. \url{https://barrel.repo.nii.ac.jp/records/2331}; persistent identifier \url{https://hdl.handle.net/10252/2056}.
\\{}[Report~228 adds: Propositions 1 and 2, printed p.~19.]
\bibitem{BW} M. Bechtloff Weising, \emph{Higher Order Bell Symmetric Functions}, arXiv:2505.18504v2. \url{https://arxiv.org/abs/2505.18504v2}. Definitions and statements checked in the versioned HTML on 5 October 2026.
\bibitem{IBD} \emph{Asymptotics of Iterated Bell Diagonals (OEIS A139383, A261280)}, unrefereed research report in VladimirReshetnikov/ProveIt, combined fixed-shift and proportional-depth manuscripts, 2 October 2026. Catalogue checked 5 October 2026; indexed commit \texttt{65969409cfc007047182e9303e1d2d0106324dea}. \href{https://github.com/VladimirReshetnikov/ProveIt/tree/main/SetTheory/Cardinals/docs/reports/generating-functions-and-asymptotics/oeis-sequence-asymptotics/a139383-iterated-bell-diagonals}{Repository report}.
\\{}[As cited in Report~228: unrefereed report in VladimirReshetnikov/ProveIt, 2 October 2026. Snapshot at main commit \texttt{5fdc0d68a062c815100ac27bd173a8c4b901c08d}, checked 5 October 2026. \href{https://github.com/VladimirReshetnikov/ProveIt/blob/5fdc0d68a062c815100ac27bd173a8c4b901c08d/SetTheory/Cardinals/docs/reports/generating-functions-and-asymptotics/oeis-sequence-asymptotics/a139383-iterated-bell-diagonals/article.tex}{Pinned public source}. Sections 4, 12, 13, and 17.] [\wtag\ In this repository: \file{SetTheory/Cardinals/docs/reports/generating-functions-and-asymptotics/oeis-sequence-asymptotics/a139383-iterated-bell-diagonals}; its \file{article.tex} is the same blob at both pins and at HEAD, and its Sections 4, 12, 13 and 17 are the sections Report~228 names.]
\\{}[Part~IV (key \texttt{BellReport}): ProveIt research report \emph{Iterated Bell diagonals}, with the proportional-depth addendum; pinned revision \Repin, blob \texttt{2b1d8fbe81}.]
\bibitem{R225} \emph{Diagonal and proportional height asymptotics for iterated Euler transforms}, Report 225, 5 October 2026. [\wtag\ Printed in this report as Part~\ref{ied:lead:part}.] As cited in Report~226: preceding unrefereed research report; unchanged PDF and editable TeX included in the archive's \texttt{supporting} directory; exact provenance in \texttt{sources/SOURCE\_LEDGER.md}. As cited in Report~228: unrefereed research report, 5 October 2026; unchanged complete TeX and PDF in the accompanying archive; provenance in \texttt{sources/SOURCE\_LEDGER.md}. [\wtag\ The included copies are byte-identical to Part~I and are not shipped; the ledgers are shipped as \file{226-firstcorr-sources-SOURCE_LEDGER.md} and \file{228-logheight-sources-SOURCE_LEDGER.md}.]
\bibitem{R226} \emph{First logarithmic correction for iterated Euler diagonals}, Report 226, unrefereed research report, 5 October 2026. Unchanged complete TeX and PDF in the accompanying archive; provenance in \texttt{sources/SOURCE\_LEDGER.md}. [\wtag\ Printed in this report as Part~\ref{ied:corr:part}; the included copy is byte-identical to it and is not shipped.]
\bibitem{DS} A. Dudko and D. Sauzin, \emph{The resurgent character of the Fatou coordinates of a simple parabolic germ}, C. R. Acad. Sci. Paris I \textbf{352} (2014), 255--261. \url{https://doi.org/10.1016/j.crma.2013.12.005}.
\bibitem{Prellberg} T. Prellberg, \emph{On the Asymptotic Analysis of a Class of Linear Recurrences}, September 2002 seminar, summary by M. Mishna, Algorithms Seminar 2002--2004 (2005), 47--50. Sections 2.2 and 3.1. \url{https://algo.inria.fr/seminars/summary/Prellberg2002a.pdf}.
\\{}[Part~IV: seminar of 23 September 2002, summary by M.~Mishna, in F.~Chyzak (ed.), \emph{Algorithms Seminar 2002--2004}, INRIA, 2005, pp.~47--50; ``The four-page primary seminar summary was opened and read''.]
\bibitem{NV} S. V. Nagaev and V. I. Vakhtel, \emph{Limit theorems for probabilities of large deviations of a Galton--Watson process}, Discrete Math. Appl. \textbf{13}(1) (2003), 1--26. Theorems 1--2 and their Poisson specialization; the source-comparison limitations are recorded in the ledger. \url{https://doi.org/10.1515/156939203321669537}; \url{https://www.mathnet.ru/eng/dm183}.
\bibitem{TSvol} [\wtag] \emph{Transseries and Inversion}, the transseries volume of this repository, \file{Analysis/Transseries/docs/series-and-transseries/Transseries_And_Inversion/transseries_and_inversion.tex}: Proposition \texttt{p0:prop:factorial-core}, Theorems \texttt{p0:thm:lambert-core}, \texttt{p0:thm:perturbed-inversion} and \texttt{p0:thm:staircase}.
\bibitem{Euler} [Part~IV] ProveIt research report, \emph{Diagonal and proportional height asymptotics for iterated Euler transforms}, with Parts II--III and editorial notes, \file{a290354-iterated-euler-diagonals/article.tex}. Pinned revision \Repin\ (blob \texttt{b5af1ca30b}). [This report's Parts~I--III, unchanged at the pin.]
\bibitem{Chains} [Part~IV] ProveIt research report, \emph{Strict partition chains and Lengyel numbers}, \file{a005121-strict-partition-chains/article.tex}. Pinned revision \Repin\ (blob \texttt{da8cac4780}).
\bibitem{Lengyel} [Part~IV] T. Lengyel, \emph{On a recurrence involving Stirling numbers}, European Journal of Combinatorics \textbf{5} (1984), 313--321. Bibliographic attribution checked in \cite{Prellberg}; the original paper was not independently reread for the companion.
\bibitem{Skau} [Part~IV] I. H. Skau and K. F. Kristensen, \emph{An asymptotic formula for the iterated exponential Bell numbers}, arXiv:1903.07979, 2019. \url{https://arxiv.org/abs/1903.07979}.
\bibitem{A005121} [Part~IV; listed, not cited in its text] The OEIS Foundation Inc., \emph{A005121}, \url{https://oeis.org/A005121}. Accessed 8 October 2026. Sequence metadata are context; the checks independently recompute every count.
\end{thebibliography}
\end{document}
